% Compilation driver for manuscript pentagon_v005.tex.
%%
% Author Thomas C. Hales et al.
% LaTeX Format


%!TEX TS-program = latex    
%% This line is for TexShop. 

% started June 2015
% submitted to arxiv Feb 23, 2016

\documentclass{amsart}
\usepackage{graphicx}
\usepackage{amsfonts}
\usepackage{amscd}
\usepackage{amssymb}
\usepackage{alltt}
\usepackage{multicol}
\usepackage{amsmath}
\usepackage{amsthm}
\usepackage{amscd}
%\usepackage{mnsymbol}

%%\documentclass[spanningrule]{cambridge7A}
%% Cambridge University Press Macros from
%% https://authornet.cambridge.org/information/productionguide/laTex_files/


% required by CUP
\usepackage[numbers]{natbib}
\usepackage{rotating}
\usepackage{floatpag}
 \rotfloatpagestyle{empty}
\usepackage{amsthm}
\usepackage{graphicx}
%\usepackage{multind}\ProvidesPackage{multind}
\usepackage{times}


% my additions
\usepackage{verbatim}
\usepackage{latexsym}
\usepackage{amsfonts}
\usepackage{amsmath}
\usepackage{crop}
\usepackage{fancyhdr}
\pagestyle{fancy}
\usepackage{txfonts}
%\usepackage[hyphens]{url}
%\usepackage{url}
\usepackage{setspace}
\usepackage{ellipsis} % 
% http://www.ctan.org/tex-archive/macros/latex/contrib/ellipsis/ellipsis.pdf 
% wrapping graphics
% http://en.wikibooks.org/wiki/LaTeX/Floats,_Figures_and_Captions
\usepackage{wrapfig}
\usepackage{listings}


%tikz graphics
\usepackage{xcolor} % to remove color.
\usepackage{tikz} % Needs pgf version >= 2.10.
\usetikzlibrary{chains,shapes,arrows,trees,matrix,positioning,backgrounds,fit,calc,fadings}

% fonts
\usepackage[mathscr]{euscript} % powerset.
\usepackage{pifont} %ding
\usepackage[displaymath]{lineno}


%%%%%%%%%%%%%%%%%%%%%%%%%%%%%%%%%%%%%%%%%
\def\tikzfig#1#2#3{%
\begin{figure}[htb]%
  \centering
\begin{tikzpicture}#3
\end{tikzpicture}
  \caption{#2}
  \label{fig:#1}%
\end{figure}%
}

\def\tikzwrap#1#2#3#4{%
\begin{wrapfigure}{r}{#4\textwidth}
  \begin{center}
\begin{tikzpicture}#3
\end{tikzpicture}
\end{center}
  \caption{#2}
  \label{fig:#1}%
\end{wrapfigure}%
}

%-%
%--Redefining--
%-%
\def\emptyset{\varnothing}
\def\ups{\upsilonup} % Needs txfonts; else use \upsilon

%% % vector notation
\def\v{{\mathbf v}}
\def\u{{\mathbf u}}
\def\w{{\mathbf w}}
\def\p{{\mathbf p}}
\def\q{{\mathbf q}}

% operatorname
\def\op#1{{\operatorname{#1}}}
\def\optt#1{{\operatorname{{\texttt{#1}}}}}
\def\card{\op{card}}
\def\arc{\operatorname{arc}}
\newcommand{\ring}[1]{\mathbb{#1}}
\def\norm#1#2{\hbox{\ensuremath{\|#1\unskip-\unskip{#2}\|}}}
\def\normo#1{{\|#1\|}}

\def\red#1{\textcolor{red}{#1}}
\def\calc{$^\dagger$}
\def\calcref#1{$^\dagger$}
% line numbers
\def\lll{\resetlinenumber[1]}
\def\linenumberfont{\normalfont\small\sffamily}

\def\tocpart#1{
  \addcontentsline{toc}{part}{\Large{#1}}}
% (or even \LARGE)
  
\crop
%\makeindex
%\makeindex{index/Notation}
%\makeindex{index/Index}

\def\linput#1{\lll\input{#1}}


\raggedbottom  % for now.
%\raggedright  % don't worry for now.

%% \renewcommand{\baselinestretch}{2}  %% turns on double spacking

%%% end notes -kill this.
%%\maketextnotes

%%%%%%%%%%%%%%%%%%%%%%%%%%%%%%%%%%%%%%%%

% new theorems

% One section-based counter shared by equations and theorem-like environments.
\numberwithin{equation}{section}
% Figures and tables have their own section-based sequences.
\numberwithin{figure}{section}
\numberwithin{table}{section}

\theoremstyle{plain}
\newtheorem{theorem}[equation]{Theorem}
%\newtheorem{theorem*}[equation]{Theorem$^*$}
\newtheorem{lemma}[equation]{Lemma}
%\newtheorem{lemdag}[equation]{Lemma$^\dagger$}
\newtheorem{background}[equation]{Background}
\newtheorem{corollary}[equation]{Corollary}
\newtheorem{example}[equation]{Example}
\newtheorem{conjecture}[equation]{Conjecture}
%---
\theoremstyle{definition}
\newtheorem{definition}[equation]{Definition}
%---
\theoremstyle{remark}
\newtheorem{remark}[equation]{Remark}
\newtheorem{notation}[equation]{Notation}
\newtheorem{exer}[equation]{Exercise}

\renewcommand\footnotemark{}
%\renewcommand\footnoterule{}

\def\fig#1#2#3{%
\begin{figure}[htb]%
  \centering
#3
  \caption{#2}
  \label{fig:#1}%
\end{figure}%
}


%%%%%%%%%%%%%%%%%%%%%%%%%%%%%%%%%%%%%%%%%

\begin{document}
\title
    {Packings of Regular Pentagons in the Plane}
\author{Thomas Hales and W\"oden Kusner}
\date{Revision history: 2/23/2016 (arXiv v1), 9/13/2016 (arXiv v2). This version: 6 November 2026}
\thanks{Research supported by NSF grant 1104102 and partially supported by Austrian Science Fund (FWF) Project 5503. We thank a referee for a detailed report. ChatGPT-6 Astra assisted in the
2026 revision of the written paper and code.}

\begin{abstract}  
We show that every packing of congruent regular pentagons in the Euclidean plane has
density at most $(5-\sqrt5)/3\approx 0.92$.     More specifically,
this article proves the pentagonal ice-ray
conjecture of Henley (1986), and Kuperberg and Kuperberg (1990), 
which asserts that an optimal packing of congruent regular pentagons in the plane is a double lattice,
formed by aligned vertical columns of upward pointing pentagons alternating
with aligned vertical columns of downward pointing pentagons.  The strategy is based on estimates
of the areas of Delaunay triangles.  Our strategy reduces the pentagonal ice-ray conjecture to
area minimization problems that involve at most four Delaunay triangles.  These minimization problems
are solved by computer.  The computer-assisted portions of the proof use techniques such as
interval arithmetic, automatic differentiation, and a meet-in-the-middle algorithm.
\end{abstract}

 \lhead{Hales and Kusner}
\rhead{Pentagon Packings}

\parskip=\baselineskip

 \maketitle

%%%


    
%%%%%%%%%%%%%%%%%%%%%%%%%%%%%%%%%%%%%%%%%
%%% FRONT

% Packings of Regular Pentagons in the Plane -- manuscript v005.
% Compile header_v005.tex. Source-only COMPUTATION codes index the companion code.
% The mathematical text is based on pentagon_v004g.tex; see provenance/paper-changes.md.

\def\area{\op{area}}
\def\areta{\op{area}^\circ}
\def\ao{a_\text{min}}
\def\acrit{a_\text{crit}}
\def\loc{\op{loc}}

\def\Ra{\Rightarrow}
\def\nRa{\nRightarrow}
\def\rab{\Ra_{b}}

\def\C{\mathcal C}
\def\SS{\mathcal S}
\def\N{\mathcal N}
\def\M{\mathcal M}
\def\T{\mathcal T}
\def\PD{\Psi D}
\def\cong{\equiv}
\def\r{{\mathbf r}}
\def\s{{\mathbf s}}
\def\bl{\ell}
\def\epsM{\epsilon_{\M}}
\def\epsN{\epsilon_{\N}}
\def\c{{\mathbf c}}
\def\K{\op{K}}

\def\dx#1{d_{#1}} 
\def\dd#1#2{\norm{\c_{#1}}{\c_{#2}}}
\def\tb{\underline{b}}
\def\tn{\bar n}

\def\libel#1{\label{#1}}
\def\rif#1{\ref{#1}}

\def\smalldot#1{\draw[fill=black] (#1) node [inner sep=0.8pt,shape=circle,fill=black] {}}
\def\graydot#1{\draw[fill=gray] (#1) node [inner sep=1.3pt,shape=circle,fill=gray] {}}
\def\whitedot#1{\draw[fill=gray] (#1) node [inner sep=1.3pt,shape=circle,fill=white,draw=black] {}}
\tikzset{dartstyle/.style={fill=black,rotate=-90,inner sep=0.7pt,dart,shape border uses incircle}}
\tikzset{grayfatpath/.style={line width=1ex,line cap=round,line join=round,draw=gray}}

\def\pent#1#2#3#4{%
\draw[red] (#1,#2) + (#3:#4cm) -- + (#3+72:#4cm) -- 
+(#3+144:#4cm) -- +(#3+216:#4cm) -- + (#3+288:#4cm) -- cycle
}

\def\pentpink#1#2#3#4#5{%
\draw[#5] (#1,#2) + (#3:#4cm) -- + (#3+72:#4cm) -- 
+(#3+144:#4cm) -- +(#3+216:#4cm) -- + (#3+288:#4cm) -- cycle
}

\def\pen#1#2#3{%
\draw[red] (#1,#2) + (#3:1cm) -- + (#3+72:1cm) -- 
  +(#3+144:1cm) -- +(#3+216:1cm) -- + (#3+288:1cm) -- cycle
}

\def\penp#1#2{%
\draw[red] (#1) + (#2:1cm) -- + (#2+72:1cm) -- 
  +(#2+144:1cm) -- +(#2+216:1cm) -- + (#2+288:1cm) -- cycle
}

\def\threepent#1#2#3#4#5#6#7#8#9{%
\pen{#1}{#2}{#3};
\pen{#4}{#5}{#6};
\pen{#7}{#8}{#9};
\draw[blue] (#1,#2) -- (#4,#5) -- (#7,#8) -- cycle
}

\def\threepentnoD#1#2#3#4#5#6#7#8#9{%
\pen{#1}{#2}{#3};
\pen{#4}{#5}{#6};
\pen{#7}{#8}{#9}
}

\def\blankfig#1{\tikzfig{#1}{Insert graphic}
{
[scale=1]
\draw(0,) circle(1cm);
}}

\centerline{\it We dedicate this article to W. Kuperberg.}

% COMPUTATION VNDQIXO | Supplementary exact paper checks: mathematica/PentagonChecks.wl
% COMPUTATION YQRQXVD | OCaml function-contract tests: tests/ocaml/
\section{Introduction}\libel{sec:intro} 

A fundamental problem in discrete geometry is to determine the highest
density of a packing in Euclidean space by congruent copies of a
convex body $C$.  For example, when $C$ is a ball of given radius,
this problem reduces to the sphere-packing problem in Euclidean space.
Besides a sphere, the simplest shape to consider is a regular polygon
$C$ in the plane.  If the regular polygon is an equilateral
triangle, square, or hexagon, then copies of $C$ tile the plane.  In
these cases, the packing problem is trivial.  The first nontrivial
case is the packing problem for congruent regular pentagons.  This
article solves that problem.

Henley and Kuperberg and Kuperberg have conjectured that the densest
packing of congruent regular pentagons in the plane is achieved by a
double-lattice arrangement: vertical columns of aligned pentagons
pointing upward, alternating with vertical columns of aligned
pentagons pointing downward (Figure~\rif{fig:double-lattice})
\cite{Kup} \cite[p.801]{henley}.  This packing of pentagons is called
the {\it pentagonal ice-ray} in Dye's book on Chinese lattice designs
\cite{dye}.  Two plates (Y3b and Y3c) in Dye's book depict the
pentagonal ice-ray, originating from Chengdu, China around 1900 CE.

We call this the {\it pentagonal ice-ray  conjecture}.
This packing has density
\[
\frac{5 - \sqrt{5}}3 \approx 0.921311.
\] %
Before our work, the best known bound on the density of packings of
regular pentagons was $0.98103$, obtained through the representation
theory of the group of isometries of the plane \cite{Val}.  Our
research is a continuation of W. Kusner's thesis \cite{Kus}, which
proves the local optimality of the pentagonal ice-ray.

As far as we know, our methods are adequate for the solution of other
related problems in geometric optimization.  The limiting factor seems
to be the availability of sufficient computer resources.

\tikzfig{double-lattice}{The pentagonal ice-ray.  All figures show
pentagons in red and Delaunay triangles in blue.}
{
\pent{0.0}{0.0}{-90}{0.4};  %
\pent{0.0}{0.724}{-90}{0.4};  %
\pent{0.0}{2*0.724}{-90}{0.4};  
\pent{0.57}{-0.3854}{90}{0.4};  %
\pent{0.57}{-0.3854+0.724}{90}{0.4}; 
\pent{0.57}{-0.3854+2*0.724}{90}{0.4};  
\pent{0.57}{-0.3854+3*0.724}{90}{0.4};  
\pent{-0.57}{-0.3854}{90}{0.4};  %
\pent{-0.57}{-0.3854+0.724}{90}{0.4}; 
\pent{-0.57}{-0.3854+2*0.724}{90}{0.4};  
\pent{-0.57}{-0.3854+3*0.724}{90}{0.4};  
\pent{2*0.57}{0.0}{-90}{0.4};  %
\pent{2*0.57}{0.724}{-90}{0.4};  %
\pent{2*0.57}{2*0.724}{-90}{0.4}; 
\pent{3*0.57}{-0.3854}{90}{0.4};  %
\pent{3*0.57}{-0.3854+0.724}{90}{0.4}; 
\pent{3*0.57}{-0.3854+2*0.724}{90}{0.4};  
\pent{3*0.57}{-0.3854+3*0.724}{90}{0.4};   
}

This article gives a computer-assisted proof of the pentagonal ice-ray
conjecture.  A discussion of the computer calculations appears in
Appendix~\rif{appendix:cc}.
The proof of the theorem appears at the end of Section~\rif{sec:nonobtuse}.

\begin{theorem}\libel{thm:main}  
  No packing of congruent regular pentagons in the Euclidean plane has
  density greater than that of the pentagonal ice-ray.  The pentagonal
  ice-ray is the unique periodic packing of congruent regular
  pentagons that attains optimal density.
\end{theorem}

In this article, we consider packings of congruent regular pentagons
in the Euclidean plane.  By a packing, we mean that the pentagons have
pairwise disjoint interiors.  By density, we mean specifically the
upper density: if ${\mathcal P}$ is a packing and if $S$ is an arbitrary square
of edge-length $r(S)$, then the density of $\mathcal P$ in the bounded
window $S$ is
\[
\delta({\mathcal P},S) = \sum_{P\in {\mathcal P},P\subset S}{\op{area}(P)}/{r(S)^2}.
\]
The upper density is defined to be
\[
\bar\delta({\mathcal P}):=\lim_{r\to\infty}
\sup_{S~\text{square},~r(S)\ge r}\delta({\mathcal P},S).
\]
In the limit, the edge effects of the window do not matter:  we obtain
the same limit by summing with the slightly more relaxed condition
$P\cap S\ne\emptyset$.

The upper density can be approximated by periodic packings.  If
$\epsilon >0$, we can find a finite window $S$ such that
$\delta({\mathcal P},S)>\bar\delta({\mathcal P})-\epsilon$.  Then the
periodic extension of $\{P\in {\mathcal P}\mid P\subset S\}$ obtained
by tiling the plane with tile $S$ has upper density at least
$\bar\delta({\mathcal P})-\epsilon$.

Density is scale-invariant.  Without loss of
generality, we may assume that all pentagons are regular pentagons of
fixed circumradius $1$.  The inradius of each pentagon is $\kappa:=
\cos (\pi/5) = (1+\sqrt{5})/4 \approx 0.809$. We set $\sigma :=
\sin(\pi/5) \approx 0.5878$.  The length of each pentagon edge is
$2\sigma$.  (Checks of the numerical constants and elementary inequalities
are collected in the supplementary Wolfram Language package
{\tt mathematica/\allowbreak PentagonChecks.wl} in the companion code.)

All pentagon packings will be assumed to be saturated; that is, no
further regular pentagons can be added to the packing without overlap.
Concretely, a packing can be saturated as follows.  Let $S_i$,
$i=1,2,\ldots$ be a sequence of nested squares whose union is the
plane.  Adding finitely many pentagons at each step, we first saturate
the packing in $S_1$, then continue with successive $S_i$.
When starting from a periodic packing, we may assume that the saturation
is periodic, by saturating with finitely many pentagons on a quotient torus.
The
assumption of saturation can be made without loss of generality,
because our ultimate aim is to give upper bounds on the density of
pentagon packings, and the saturation of a packing cannot decrease its
density.

We form the Delaunay triangulations of the pentagon packings.  The
vertices of the triangles are always taken to be the centers of the
pentagons.  The Delaunay triangulation can have degeneracies when two
or more triangles have the same circumcircle.  Fix any
translation-invariant degeneracy-breaking rule in such cases.  The
saturation hypothesis implies that no Delaunay triangle has
circumradius greater than $2$, because whenever a circumradius is at
least $2$, a new pentagon can be inserted at the circumcenter.  This
property of Delaunay triangles is crucial.  Every edge of a Delaunay
triangle has length at most $4$; that is, the maximum diameter of a
Delaunay triangle.

Each Delaunay triangle in a pentagon packing has edge lengths at least
$2\kappa$.  This minimum Delaunay edge length is attained exactly when
the two pentagons have a full edge in common.

For most of the article, we consider a fixed saturated packing and its
Delaunay triangulation.  Generally, unless otherwise stated,
every triangle is a  Delaunay triangle.
Statements of lemmas and theorems implicitly
assume this fixed context.

Initially, pentagons in a packing play two roles: they constrain the
shapes of the Delaunay triangles and they carry mass for the density.
We prefer to change our model slightly so pentagons are only used
to constrain the shapes of triangles.  We replace the mass of each
pentagon by a small, massive, circular disk (each of the same small
radius) centered at the center of the pentagon, and of uniform density
and the same total mass as the pentagon.  In this model, each Delaunay
triangle contains exactly one-half the mass of a pentagon.  By
distributing the pentagon mass uniformly among the Delaunay triangles,
we may replace density maximization with Delaunay triangle area
minimization.

We write 
\[
\acrit := \frac{3}{2}{ \sigma \kappa(1+\kappa)} =
\frac{3}{16}\sqrt{5(5+2\sqrt5)}\approx 1.29036
\] %
for the common (critical) area of every Delaunay triangle in the 
pentagonal ice-ray.

This article is organized as follows.  The aim is to show that on
average, every Delaunay triangle has area at least $\acrit$.  There
are two primary averaging strategies.  One strategy is to transfer
small areas from one triangle to another in a way that preserves the
total area.  Areas of two different sizes $\epsN$ and $\epsM$ are
moved between triangles. (See Figure~\rif{fig:visual}.)  Section~\rif{sec:cluster}
defines the constants $\epsN,\epsM$ and gives the rules that govern
these transfers of area.  The modified area of a triangle after area
transfer is denoted $b(T)$.  A second strategy is to define finite
disjoint clusters of triangles such that the average of the modified
areas $b(T)$ over each cluster is at least $\acrit$.  Lemma~\rif{lemma:main}
states the main inequality:
lower bounds on averages of modified areas over finite
clusters imply the pentagonal ice-ray conjecture.
Section~\rif{sec:cluster} also gives the definition of these clusters.
Two particularly important types of pairs of adjacent triangles are
called dimer pairs and pseudo-dimer pairs.  A dimer pair is a pair of
triangles that share a common longest edge and that have properties
similar to pairs of triangles in the ice-ray packing.
A pseudo-dimer pair is a pair of adjacent triangles that has properties
somewhat less similar to triangles in the ice-ray packing.  Yet the area
of pseudo-dimers is so small that they require special attention.
Section~\rif{sec:cluster} also gives a long list of examples.

A $P$-triangle is a Delaunay triangle together with the three pentagons
at the three vertices. Section~\rif{sec:contact} makes a detailed
study of the structure of $P$-triangles.  Of particular interest are
$P$-triangles in which each pentagon is in contact with the other two.
The section gives a classification of the possible combinatorial types
of such $P$-triangles.

The purpose of Section~\rif{sec:area} is to give lower bounds
on the areas of Delaunay triangles -- especially acute triangles.
To simplify the area estimates, 
the section introduces some local deformations of $P$-triangles that
are monotonic in area.  The section also contains estimates on
the lengths of edges of various Delaunay triangle.

Section~\rif{sec:calc} states a series of lemmas whose proofs are
computer assisted using interval arithmetic.  The proofs are delayed until 
Appendix~\rif{appendix:cc}.  The computer part of these proofs has
been organized as a github repository~\cite{Git}.

Section~\rif{sec:dimer-pair} proves the local optimality of the
pentagonal ice-ray.  This is stated as optimization problem involving
pairs of acute Delaunay triangles (called dimer pairs), which share
their common longest edge.  The unique dimer pair that minimizes area
is the dimer pair found in the pentagonal ice-ray packing.

This local optimality result does not in itself give local optimality
in the form that is needed for the pentagonal ice-ray conjecture.  We
need to supplement this result with lemmas that state the
compatibility of dimer pairs with the partition into finite clusters
of Delaunay triangles made in Section~\rif{sec:cluster}.  We also need
to show that the modified area $b(T)$ equals the unmodified area
$\op{area}(T)$ for triangles in dimer pairs.  These supplementary
results about dimer pairs are established in
Section~\rif{sec:dimer-main}.

Section~\rif{sec:pseudo} proves a series of preliminary lemmas about
pseudo-dimers.  Pseudo-dimers are pairs of Delaunay triangles that
fail to satisfy all the properties of a dimer pair, yet they have such
small area that they require special consideration.  Among other
results, the section proves that pseudo-dimer pairs are disjoint from
one another, that the area of one of the triangles in the pair is
unmodified: $b(T_1)=\area(T_1)$, and that the other triangle in the
pair has modified area greater than the target area: $b(T_0)>\acrit$.
These lemmas will be used in Section~\rif{sec:nonobtuse}.

The final two sections~\rif{sec:obtuse} and \rif{sec:nonobtuse}
complete the proof of the pentagonal ice-ray conjecture by proving the
main inequality: the modified area averages over clusters of Delaunay
triangles are at least the area of the Delaunay triangle in the
pentagonal-ice ray.  Section~\rif{sec:obtuse} treats the case when at
least one Delaunay triangle in the cluster is obtuse, and
Section~\rif{sec:nonobtuse} treats the remaining case: no Delaunay
triangle in the cluster is obtuse.

There are two appendices.  The first, Appendix~\rif{sec:coords},
introduces explicit coordinate systems for configurations of
pentagons, including pairs of pentagons, and $P$-triangles (triples of
pentagons).  These coordinate systems are not explicitly used in this
article.  However, they are used heavily in the computer code to give
explicit implementations of geometric objects and functions that are
needed for computer proofs.  This section serves as
documentation of geometric aspects of the computer code.

Appendix~\rif{appendix:cc} also serves as documentation of the
computer code.  It describes how the computer code has been
organized. It describes computer algorithms such as interval
arithmetic, automatic differentiation, and meet-in-the-middle.  A
final subsection~\rif{sec:calc-prep} briefly discusses each of the
computer-assisted lemmas that are stated (without proof) in
Section~\rif{sec:calc}.  This subsection gives reduction arguments,
showing how simpler computer calculations can be used to prove the
lemmas in their stated generality.
Additionally, all floating point
approximations marked $\approx$ and all inequality comparisons with a floating-point number have been checked by computer.

\section{Clusters of Delaunay Triangles}\libel{sec:cluster}

The area of a Delaunay triangle in a saturated pentagon packing can be
smaller than $\acrit$.  Our strategy for proving the pentagonal
ice-ray conjecture is to collect triangles into finite clusters such
that the average area over each cluster is at least $\acrit$.

We say that a Delaunay triangle (in any pentagon packing) is {\it
  subcritical} if its area is at most $\acrit$.  We will obtain a
lower bound $\ao := 1.237$ on the area of a nonobtuse Delaunay
triangle (Lemma~\rif{lemma:a0}).  This is a very good bound.  It is
very close to the numerically smallest achievable area,%
\footnote{In the notation of the appendix, the
  numerical minimum is achieved by a pinwheel with parameters
  $\alpha=\beta=0$ and $x_\gamma\approx 0.16246$.}
 which is
approximately $1.23719$.  We write 
\begin{equation}\libel{eqn:epsN}
\epsN :=
\acrit - \ao \approx 0.05336,
\end{equation} for the difference between the desired
bound $\acrit$ on averages of triangles and the bound $\ao$ for a
single nonobtuse triangle.  Set
\begin{equation}\libel{eqn:epsM}
\epsM := 0.008.
\end{equation}

\bigskip
\fig{constant}{Key real constants}%
{
\begin{tabular}{llll}
{\bf Constant} & {\bf Approx. Value} & {\bf Note} \\
& 0.92 & ice-ray density\\
$\kappa$ & $\cos(\pi/5)\approx 0.809$ & pentagon inradius \\
$\sigma$ &$\sin(\pi/5)\approx 0.5878$ & pentagon half-edge \\
$\acrit$ &1.29 & ice-ray triangle area\\
$\ao$ & 1.237 & non-obtuse triangle area lower bound\\
$\epsN$ & 0.05336 & correction term\\
$\epsM$ & 0.008 & correction term\\
\end{tabular}
}

\subsection{Examples}

While reading this article, it is useful to carry along a series of examples,
illustrated in Figures~\rif{fig:KKtriangle} -- \rif{fig:obtuse}.
Otherwise, later definitions such as the modified area function $b(T)$
and the construction of clusters might appear to be unmotivated.  Some
of these examples serve as counterexamples to naive approaches to this
problem.  These different examples can interact with one another in
potentially complicated ways.  The proof of the main theorem must sort
through these interactions.

\tikzfig{KKtriangle}{Ice-ray triangles and ice-ray dimers.
  All Delaunay triangles in the pentagonal ice-ray are congruent and
  have area $\acrit$.  We call them ice-ray triangles. 
  The pentagonal ice-ray can be partitioned into pairs of Delaunay
  triangles (called ice-ray dimers) in which the two triangles in each
  dimer share their common longest edge.  
}
{
\begin{scope}
\pentpink{0.0}{0.0}{-90}{0.4}{red!30};  %
\pentpink{0.0}{0.724}{-90}{0.4}{red!30};  %
\pentpink{0.0}{2*0.724}{-90}{0.4}{red!30};  
\pentpink{0.57}{-0.3854}{90}{0.4}{red!30};  %
\pentpink{0.57}{-0.3854+0.724}{90}{0.4}{red!30}; 
\pentpink{0.57}{-0.3854+2*0.724}{90}{0.4}{red!30};  
\pentpink{0.57}{-0.3854+3*0.724}{90}{0.4}{red!30};  
\pentpink{-0.57}{-0.3854}{90}{0.4}{red!30};  %
\pentpink{-0.57}{-0.3854+0.724}{90}{0.4}{red!30}; 
\pentpink{-0.57}{-0.3854+2*0.724}{90}{0.4}{red!30};  
\pentpink{-0.57}{-0.3854+3*0.724}{90}{0.4}{red!30};  
\pentpink{2*0.57}{0.0}{-90}{0.4}{red!30};  %
\pentpink{2*0.57}{0.724}{-90}{0.4}{red!30};  %
\pentpink{2*0.57}{2*0.724}{-90}{0.4}{red!30}; 
\pentpink{3*0.57}{-0.3854}{90}{0.4}{red!30};  %
\pentpink{3*0.57}{-0.3854+0.724}{90}{0.4}{red!30}; 
\pentpink{3*0.57}{-0.3854+2*0.724}{90}{0.4}{red!30};  
\pentpink{3*0.57}{-0.3854+3*0.724}{90}{0.4}{red!30};   
\coordinate (T) at (0.5706,-0.38541);
\coordinate (U) at (0,0.7236);
\coordinate (T') at ($(T)+(U)$);
\draw[blue] (0,0) -- ++ ($2*(U)$);
\draw[blue] ($-1.0*(T')$) -- ++ ($3*(U)$);
\draw[blue] (T) -- ++ ($3*(U)$);
\draw[blue] ($(T)+(T')$) -- ++ ($2*(U)$);
\draw[blue] ($(T)+(T')+(T)$) -- ++ ($3*(U)$);
\draw[blue] ($-1.0*(T)$) -- ++ ($2.0*(T)$);
\draw[blue] ($-1.0*(T) + (U)$) -- ++ ($4.0*(T)$);
\draw[blue] ($-1.0*(T) + 2.0*(U)$) -- ++ ($4.0*(T)$);
\draw[blue] ($1.0*(T) + 3.0*(U)$) -- ++ ($2.0*(T)$);
\draw[blue] ($(T)$) -- ++ ($2.0*(T')$);
\draw[blue] ($1.0*(T) +1.0*(U) -2.0 *(T')$) -- ++ ($4.0*(T')$);
\draw[blue] ($1.0*(T) +2.0*(U) -2.0 *(T')$) -- ++ ($4.0*(T')$);
\draw[blue] ($1.0*(T) +3.0*(U) -2.0 *(T')$) -- ++ ($2.0*(T')$);
\end{scope}
\begin{scope}[xshift=5cm]
\pentpink{0.0}{0.0}{-90}{0.4}{red!30};  %
\pentpink{0.0}{0.724}{-90}{0.4}{red!30};  %
\pentpink{0.0}{2*0.724}{-90}{0.4}{red!30};  
\pentpink{0.57}{-0.3854}{90}{0.4}{red!30};  %
\pentpink{0.57}{-0.3854+0.724}{90}{0.4}{red!30}; 
\pentpink{0.57}{-0.3854+2*0.724}{90}{0.4}{red!30};  
\pentpink{0.57}{-0.3854+3*0.724}{90}{0.4}{red!30};  
\pentpink{-0.57}{-0.3854}{90}{0.4}{red!30};  %
\pentpink{-0.57}{-0.3854+0.724}{90}{0.4}{red!30}; 
\pentpink{-0.57}{-0.3854+2*0.724}{90}{0.4}{red!30};  
\pentpink{-0.57}{-0.3854+3*0.724}{90}{0.4}{red!30};  
\pentpink{2*0.57}{0.0}{-90}{0.4}{red!30};  %
\pentpink{2*0.57}{0.724}{-90}{0.4}{red!30};  %
\pentpink{2*0.57}{2*0.724}{-90}{0.4}{red!30}; 
\pentpink{3*0.57}{-0.3854}{90}{0.4}{red!30};  %
\pentpink{3*0.57}{-0.3854+0.724}{90}{0.4}{red!30}; 
\pentpink{3*0.57}{-0.3854+2*0.724}{90}{0.4}{red!30};  
\pentpink{3*0.57}{-0.3854+3*0.724}{90}{0.4}{red!30};   
\coordinate (T) at (0.5706,-0.38541);
\coordinate (U) at (0,0.7236);
\coordinate (T') at ($(T)+(U)$);
\draw[blue] ($-1.0*(T)$) -- ++ ($2.0*(T)$);
\draw[blue] ($-1.0*(T) + (U)$) -- ++ ($4.0*(T)$);
\draw[blue] ($-1.0*(T) + 2.0*(U)$) -- ++ ($4.0*(T)$);
\draw[blue] ($1.0*(T) + 3.0*(U)$) -- ++ ($2.0*(T)$);
\draw[blue] ($(T)$) -- ++ ($2.0*(T')$);
\draw[blue] ($1.0*(T) +1.0*(U) -2.0 *(T')$) -- ++ ($4.0*(T')$);
\draw[blue] ($1.0*(T) +2.0*(U) -2.0 *(T')$) -- ++ ($4.0*(T')$);
\draw[blue] ($1.0*(T) +3.0*(U) -2.0 *(T')$) -- ++ ($2.0*(T')$);
\draw[blue] ($-1.0*(T) + 2.0*(U)$) -- ++ ($-1.0*(T')$) -- ++ (T)
  -- ++ ($-1.0*(T')$) -- ++ (T)
  -- ++ ($-1.0*(T')$) -- ++ (T);
\draw[blue] ($3.0*(T) + 4.0*(U)$) -- ++ (T) -- ++ ($-1.0*(T')$) 
  -- ++ (T)
  -- ++ ($-1.0*(T')$) 
  -- ++ (T)
  -- ++ ($-1.0*(T')$);
\end{scope}
}

\tikzfig{subcritical}{Subcritical triangles.
Experimentally, the first triangle has the smallest area among all
acute Delaunay triangles.  The second triangle is a cloverleaf that
has two edges of minimal length $2\kappa$.  Experimentally, the
third triangle minimizes the longest edge among subcritical acute Delaunay triangles.
It is equilateral with edge length about $1.72256$.}
{
\begin{scope}[scale=0.5]
\threepent{0.00}{0.00}{85.23}{0.93}{1.33}{265.23}{1.86}{0.00}{229.23};
\draw (0.5,-1.5) node {$\area\approx 1.23719$};
\end{scope}
\begin{scope}[scale=0.5,xshift=5cm]
\threepent{0.00}{0.00}{90.00}{0.95}{1.31}{270.00}{1.90}{0.00}{234.00};
\draw (0.5,-1.5) node {$\area\approx 1.24$};
\end{scope}
\begin{scope}[scale=0.5,xshift=10cm]
\threepent{0.00}{0.00}{96.03}{0.86}{1.49}{264.03}{1.72}{0.00}{216.03};
\draw (0.5,-1.5) node {$\area\approx 1.285$};
\end{scope}
\begin{scope}[scale=0.5,xshift=15cm] %
\threepent{0.00}{0.00}{109.36}{0.83}{1.59}{235.36}{1.62}{0.00}{145.36};
\draw (0.5,-1.5) node {$\area\approx 1.286$};
\end{scope}
}

\tikzfig{KK2min}{Ice-ray triangle deformation.
The ice-ray triangle is not a local minimum of the area function.
The  ice-ray triangle (left) can be continuously deformed
along an area decreasing curve to the subcritical acute triangle of 
numerically minimum area (right). }
{
\begin{scope}[scale=0.5]
\threepent{0.00}{0.00}{72.00}{0.96}{1.43}{252.00}{1.81}{0.00}{216.00};
\draw (0.5,-1.5) node {$\area=\acrit$};
\end{scope}
\begin{scope}[scale=0.5,xshift=6cm]
\threepent{0.00}{0.00}{79.49}{0.93}{1.37}{259.49}{1.82}{0.00}{223.49};
\draw (0.5,-1.5) node {$\area\approx 1.248$};
\end{scope}
\begin{scope}[scale=0.5,xshift=12cm]
\threepent{0.00}{0.00}{85.23}{0.93}{1.33}{265.23}{1.86}{0.00}{229.23};
\draw (0.5,-1.5) node {$\area\approx 1.23719$};
\end{scope}
}

\tikzfig{Gamma}{Ice-ray dimer deformation.  The ice-ray
  dimer (center) admits a shear motion that preserves all edges of
  contact.  This deformation increases the area of the dimer.  It is
  obvious by the symmetry of the figures on the left and right that
  the area function along this deformation has a critical point at the
  ice-ray dimer.  It is known that the ice-ray dimer is a local
  minimum of the area function \cite{Kus}.  }
{
\begin{scope}[scale=0.5]
\threepent{0.00}{0.00}{72.00}{0.96}{1.43}{252.00}{1.81}{0.00}{216.00};
\threepent{0.00}{0.00}{-72.00}{0.96}{-1.43}{-252.00}{1.81}{0.00}{-216.00};
\draw (0.5,-3.5) node {$\area=2\acrit$};
\end{scope}
\begin{scope}[scale=0.5,xshift=6cm]
\threepent{0.00}{0.00}{79.49}{0.93}{1.37}{259.49}{1.82}{0.00}{223.49};
\threepent{0.00}{0.00}{-64.51}{1.03}{-1.50}{-244.51}{1.82}{0.00}{-208.51};
\draw (0.5,-3.5) node {$\area\approx 2\acrit+0.03$};
\end{scope}
\begin{scope}[scale=0.5,xshift=-6cm]
\threepent{0.00}{0.00}{-79.49}{0.93}{-1.37}{-259.49}{1.82}{0.00}{-223.49};
\threepent{0.00}{0.00}{64.51}{1.03}{1.50}{244.51}{1.82}{0.00}{208.51};
\draw (0.5,-3.5) node {$\area\approx 2\acrit+0.03$};
\end{scope}
}

\tikzfig{pseudo-dimer}{Pseudo-dimer.  There exist pairs of acute
  triangles $(T_1,T_0)$ such that the sum of their two areas is at
  most $2\acrit$ and such that (1) $T_1$ is subcritical and is
  adjacent to $T_0$ along the longest edge of $T_1$, but such that (2)
  the longest edge of $T_0$ is not the edge shared with $T_1$.  We
  call such pairs {\it pseudo-dimers}.    The illustrated
  pseudo-dimer has area about $2\acrit - 10^{-5}$, and the triangle $T_0$
  has longest edge length about $1.84$.  The modification
  term $\epsM = 0.008$ is based on this and closely related examples
  of pseudo-dimers.  Pseudo-dimers add
  significant complications to the proof. }
{
\begin{scope}[scale=1.2,rotate=90]
\threepent{0.00}{0.00}{93.42}{0.90}{1.48}{259.93}{1.73}{0.00}{214.93};
\threepent{0.00}{0.00}{-50.58}{1.10}{-1.49}{-181.07}{1.73}{0.00}{-217.07};
\draw (1.3,-0.1) node {$T_1\Rightarrow T_0$};
\draw (0.8,-1) node {$\Searrow$};
\end{scope}
}

\tikzfig{obtuse}{Obtuse Delaunay triangles can have small area.  A
  triangle with circumradius about $2$ and area about $0.98$
  is shown. The adjacent Delaunay triangle cannot have a vertex
  inside the circumcircle of the first.  The neighbor of a subcritical
  obtuse Delaunay triangle tends to have large area.}
{
\begin{scope}[scale=0.5]
\threepent{0}{0}{90}
{1.486}{-0.661}{-90}
{-1.486}{-0.661}{-90};
\end{scope}
\begin{scope}[scale=0.5,xshift=6cm]
\draw[blue] (0,0)-- (1.486,-0.661) -- (-1.486,-0.661) -- cycle;
\draw[blue] (1.486,-0.661) -- ($(0,-2) + (-6:2)$) -- (-1.486,-0.661);
\draw (0,-2) circle (2cm);
\end{scope}
}

\subsection{Attachment, modified area, and clusters}

As mentioned above, our strategy for proving the pentagonal ice-ray
conjecture is to collect triangles into finite clusters such that the
average area over each cluster is at least $\acrit$.  The clusters are
defined by an equivalence relation.  The equivalence relation, in
turn, is defined as the reflexive, symmetric, transitive closure of a
further relation on the set of Delaunay triangles in a pentagon
packing.

A second strategy is to replace the area function $\area(T)$ on
Delaunay triangles with a modified area function $b(T)$.  The
modification steals area from nearby triangles that have area to spare
and gives to triangles in need.  It will be sufficient to prove that
the average of $b(T)$ over each cluster is at least $\acrit$.

In more detail, below, we define a particular relation $(\rab)$,
viewing a relation in the usual way as a set of ordered pairs.  For
this relation $(\rab)$, we write ${(\equiv_{b})}$ for the equivalence
relation obtained as the reflexive, symmetric, transitive closure of
$(\rab)$.  We call a corresponding equivalence classes $\C$ a {\it
  cluster}.  In other words, the relation $(\rab)$ defines a directed
graph whose nodes are the Delaunay triangles of a pentagon packing,
with directed edges given as arrows $T_1\rab T_0$.  A cluster is the
set of nodes in a connected component of the underlying undirected
graph.

We say that Delaunay triangle $T_1$ {\it attaches to Delaunay
  triangle} $T_0$ when the following condition holds: $T_0$ is the
adjacent triangle to $T_1$ along the longest edge of $T_1$.  (If the
triangle $T_1$ has more than one equally longest edge, fix once and
for all a choice among them, and use this choice to determine the
triangle that $T_1$ attaches to.  Thus, $T_1$ always attaches to
exactly one triangle $T_0$.  We can assume that the tie-breaking
choices are made according to a translation invariant rule.)  We write
$T_1 \Ra T_0$ for the attachment relation {\it $T_1$ attaches to
  $T_0$}.  Although it is not always possible to adhere to the
convention, note our general convention to use descending subscripts
$i > j$ for attachment $T_i\Ra T_j$ of triangles.  We also follow a
general convention of letting $T_0$ denote a triangle that is the
target of other triangles in the same cluster.

If $\T$ is a finite set of Delaunay triangles, we set
$\area(\T):=\sum_{T\in\T} \area(T)$.  The following are key
definitions of this article: dimer pair, pseudo-dimer, $\N$, $\M$,
$n_\pm$, $m_\pm$, $b(T)$, $(\rab)$, and cluster.

\begin{definition}[dimer pair]
  We define a dimer pair to be an ordered pair $(T_1,T_0)$ of Delaunay
  triangles such that
\begin{enumerate}
\item $T_0$ and $T_1$ are both nonobtuse.
\item $T_1\Ra T_0$, and $T_1$ is subcritical.
\item $T_0 \Ra T_1$.
\item $\area\{T_1,T_0\} \le 2\acrit$.
\end{enumerate}
We write $DP$ for the set of dimer pairs.
\end{definition}
In particular, the two triangles in a dimer pair attach to each other.

\begin{definition}[pseudo-dimer]
  We define a pseudo-dimer to be an ordered pair $(T_1,T_0)$ of
  Delaunay triangles such that
\begin{enumerate}
\item $T_0$ and $T_1$ are both nonobtuse;
\item $T_1\Ra T_0$, and  $T_1$ is subcritical.
\item $T_0 \nRightarrow T_1$;
\item $\area\{T_1,T_0\} \le 2\acrit$;
\end{enumerate}
We write $\PD$ for the set of pseudo-dimers.
\end{definition}

We observe that dimer pairs differ from pseudo-dimers in the third
defining condition, which is a condition on the location of the
longest edge of $T_0$.  Every pseudo-dimer determines a third triangle
$T_-$ by the condition $T_0\Ra T_-$.  The shared edge of $T_0$ and
$T_-$, leading out of the pseudo-dimer is called the {\it egressive} edge of
the pseudo-dimer or of the triangle $T_0$.

For any set $\SS$ of ordered pairs, and any Delaunay triangles $T_+$
and $T_-$, let
\[
n_+(T_+,\SS) = \card \{T_-\mid (T_+,T_-)\in \SS\}
\quad\text{ and }\quad
n_-(T_-,\SS)
= \card \{T_+\mid (T_+,T_-)\in \SS\}.
\]

We define $\N$ as the disjoint union of two sets of ordered pairs: $\N
= \N_{\text{obtuse}} \sqcup \N_{\text{nonobtuse}}$.  The set
$\N_{\text{obtuse}}$ consists of those pairs with obtuse target, and
$\N_{\text{nonobtuse}}$ are those pairs with nonobtuse target $T_-$,
as follows.

Define $\N_{\text{obtuse}}$ to be the set of pairs $(T_+,T_-)$ of
Delaunay triangles such that
\begin{enumerate}
\item $T_-$ is obtuse;
\item $T_+\Ra T_-$;
\item $T_+$ is nonobtuse, and the longest edge of $T_+$ has length at
  least $1.72$.
\end{enumerate}

Define  $\N_{\text{nonobtuse}}$ to 
be the set of pairs $(T_+,T_-)$ of Delaunay triangles such that
\begin{enumerate}
\item $T_-$ is nonobtuse;
\item $T_+\Ra T_-$;
\item $T_+$ is nonobtuse and the longest edge of $T_+$ has length at
  least $1.72$;
\item there exists an obtuse triangle $T$ such that $T\Ra T_-$ and 
\[
\area(T) -   n_-(T,\N_{\text{obtuse}})\epsN\le \acrit.
\]
\end{enumerate}

Let $\M$ be the set of pairs $(T_+,T_-)$ of Delaunay triangles such that
\begin{enumerate}
\item $(T_+,T_-)\not\in \N$;
\item $T_+\Ra T_-$;
\item The longest edge of $T_+$ has length  at least $1.72$.
\item There exists a unique $T_1$ such that $(T_1,T_+)\in \PD$.
\end{enumerate}

We abbreviate 
\[
m_+(T) = n_+(T,\M), 
\quad m_-(T) = n_-(T,\M), 
\quad n_+(T) =n_+(T,\N), 
\ \text{ and } \ 
n_-(T) = n_-(T,\N).
\]

\begin{remark}
  Note  that $n_+(T)\le 1$, because each triangle attaches to
  exactly one other triangle.  Also, $n_-(T)\le 3$, because each
  attachment forms along an edge of the triangle $T$.  Similarly,
  $m_+(T)\le 1$ and $m_-(T)\le 3$.
\end{remark}

We define the modified area function using the constants $\epsN,\epsM$
in Equations \eqref{eqn:epsN}, \eqref{eqn:epsM}:
\begin{equation}
b(T) := \op{area}(T) + \epsN (n_+(T) - n_-(T)) + \epsM (m_+(T) - m_-(T)).
\end{equation}
We say that $T$ is {\it $b$-subcritical} if $b(T) \le \acrit$.  We define
a relation $(\rab)$, writing
$T_1\rab T_0$, if $T_1\Ra T_0$ and $T_1$ is $b$-subcritical.  An
equivalence class of triangles under the corresponding equivalence
relation $(\equiv_b)$ is called a {\it cluster}.

We note that $(\rab)$ is given by a translation-invariant rule.  The
function $b(T)$ is also translation invariant and depends only on
local information in the pentagon packing near the triangle $T$.

The intuitive basis of using $b(T)\le \acrit$ as the condition for
cluster formation with $(\rab)$ is the following.  Eventually, we wish
to show that the average of the modified areas $b(T)$ over each
cluster is greater than $\acrit$.  (See Lemma~\rif{lemma:main}.)  If
$b(T)\le \acrit$, then this means that its modified area $b(T)$ is
less than our desired goal for the average over the cluster, so that
$T$ needs to be part of a larger cluster.  This suggests we should
define $(\rab)$ in such a way that a further triangle is added
whenever $b(T)\le \acrit$. That is what our definition of $(\rab)$
does.

We remark that the modification $b(T)$ of the area function has two
terms.  The first term $\epsN (n_+(T) - n_-(T))$ takes away
from obtuse triangles (and their neighbors) and gives to nonobtuse
triangles.  The intuition behind this modification is that the
triangle adjacent to an obtuse triangle along its long edge has a very
large surplus area that can be beneficially redistributed
(Figure~\rif{fig:obtuse}).  It allows us to make a clean separation of
the proof of the main inequality into two cases: clusters that contain
an obtuse triangle and clusters that do not.

The second term $\epsM (m_+(T) - m_-(T))$ augments the area of a
pseudo-dimer, by taking from a neighboring triangle.  The intuition
behind this modification is that a pseudo-dimer can have area
strictly less than the ice-ray dimer, and we need to boost its area
with a modification to make it satisfy the main inequality
(Figure~\rif{fig:pseudo-dimer}).  A calculation given below shows that
the neighbor of the pseudo-dimer has area to spare
(Corollary~\rif{lemma:egress'}).

\tikzfig{visual}{The area modification $b(T)$ can
be visualized as a small semicircular 
knob (of area $\epsN$) or triangular notch (of area $\epsM$) 
that transfers area
from one Delaunay triangle to an adjacent one. 
}
{
\begin{scope}[scale=2,xshift=0cm]
\coordinate (T) at (0.5706,-0.38541);
\coordinate (U) at (0,0.7236);
\coordinate (T') at ($(T)+(U)$);
\draw[blue] (U) -- ++ (T) -- ++ ($-1.0*(T')$) -- ++ ($-1.0*(T)$) -- ++ (T');
\draw[blue] (0,0) -- ++ ($0.5*(U)-(0,0.08)$);
\draw[blue] (0, 0.5*0.7236 - 0.08)
  arc[start angle=-90, end angle=90, radius=0.08];
\draw[blue] ($0.5*(U)+(0,0.08)$) -- (U);
\end{scope}
\begin{scope}[scale=2,xshift=1.5cm]
\coordinate (T) at (0.5706,-0.38541);
\coordinate (U) at (0,0.7236);
\coordinate (T') at ($(T)+(U)$);
\draw[blue] (U) -- ++ (T) -- ++ ($-1.0*(T')$) -- ++ ($-1.0*(T)$) -- ++ (T');
\draw[blue] (0,0) -- ++ ($0.5*(U)-(0,0.055/2.0)$) -- ++ (30:0.055) --++ (150:0.055)--(U) ;
\end{scope}
}

When $(T_+,T_-)\in\N$, the formula for $b(T)$ adds to $\op{area}(T_+)$
and subtracts from $\op{area}(T_-)$. 
We can visualize the modified area geometrically as an actual
transfer in the shape of a knob (or some other shape) of area $\epsN$.
Similarly, when $(T_+,T_-)\in\M$, we visualize a transfer of area in
the shape of a notch of area $\epsM$.  (See Figure~\rif{fig:visual}.)

The modification terms allow us to keep the size of each cluster small.
Eventually, we show that each cluster contains at most four triangles
(Lemma~\rif{lemma:card4}).  This small size will be helpful when we
turn to the computer interval-arithmetic calculations.  If $(T_+,T_-)$ is a member of $\N$
or $\M$, the rough expectation is that there should not be an arrow
$T_+\rab T_-$ and that $T_+$ and $T_-$ should belong to different
clusters.  That is, $\N$ and $\M$ are designed to mark cluster
boundaries.  Some lemmas in this article make this expectation more
precise (for example, Lemma~\rif{lemma:Nb}).

We give some simple consequences of our definitions.

\begin{lemma}\libel{lemma:M-nonobtuse}
  If $(T_+,T_-)\in \M$, then both $T_+$ and $T_-$ are nonobtuse.
\end{lemma}

\begin{proof} 
  By the definition of pseudo-dimer, $T_+$ is nonobtuse, because
  $(T_1,T_+)\in\PD$ for some $T_1$.  If $T_-$ were obtuse, then we
  would satisfy all the membership conditions for $(T_+,T_-)\in \N$
  (obtuse target), which is impossible because $\N\cap \M =
  \emptyset$.
\end{proof}

\begin{lemma}[obtuse $b$]\libel{lemma:obtuse-b} 
  If $T$ is an obtuse Delaunay triangle, then
  $m_+(T)=m_-(T)=n_+(T)=0$.  Thus, $b(T) = \area(T) - \epsN n_-(T)$.
\end{lemma}

\begin{proof} 
  The previous lemma gives $m_+(T)=m_-(T)=0$.  The first component of
  $\N$ is nonobtuse by definition, so $n_+(T)=0$.
\end{proof}

\begin{corollary}\libel{cor:obtuse-b-arrow} 
  If $(T_+,T_-)\in\N$ with nonobtuse target $T_-$, then there exists
  an obtuse triangle $T$ such that $T\rab T_-$.
\end{corollary}

\begin{proof} 
  Let $T$ be the obtuse triangle such that $T\Ra T_-$ given by $\N$
  (nonobtuse target), condition 4.  Then $n_-(T,\N_{\text{obtuse}}) =
  n_-(T)$, and by Lemma~\rif{lemma:obtuse-b}, we have $b(T) =
  \area(T)-\epsN n_-(T)\le\acrit$.  The result follows from the
  definition of $(\rab)$.
\end{proof}

\begin{lemma}\libel{lemma:n1m1}  
  Suppose that $n_+(T)>0$. Then $m_+(T)=0$.
\end{lemma}

\begin{proof} 
  This follows directly from the disjointness of $\M$ and $\N$.
\end{proof}

\begin{lemma}\libel{lemma:n2m2}  
  Suppose $n_-(T_-)>0$.  Then $m_-(T_-)=0$.
\end{lemma}

\begin{proof}
Assume for a contradiction that $n_-(T_-)>0$ and $m_-(T_-)>0$.  The
inequality $m_-(T_-)>0$ gives the existence of a Delaunay triangle
$T_+$ such that $(T_+,T_-)\in\M$.  By Lemma~\rif{lemma:M-nonobtuse},
triangles $T_+$ and $T_-$ are nonobtuse.  We check
that $(T_+,T_-)$ satisfies each membership condition of $\N$
(nonobtuse target), so that $(T_+,T_-)\in\N$.  All membership
conditions except (4) are elementary. To obtain condition (4), we use
the definition of $\N$ (nonobtuse target) and the inequality
$n_-(T_-)>0$ to obtain the existence of some $(T'_+,T_-)\in\N$ and of
an obtuse $T$ with $T\Ra T_-$ (by $\N$ condition 4 for $(T'_+,T_-)$).  Thus,
$(T_+,T_-)\in\N$. Disjointness gives the contradiction
$(T_+,T_-)\in\M\cap\N=\emptyset$.
\end{proof}

\subsection{The main inequality}

\begin{lemma}\libel{lemma:main}  
If  every cluster $\C$ in every saturated packing of regular
  pentagons is finite, and if for some $a$ every cluster average
  satisfies
\begin{equation}\libel{eqn:main}
\frac{\sum_{T\in \C} b(T)}{\card(\C)} \ge a,
\end{equation}
then the density of a packing of regular pentagons never exceeds 
\[
\frac{\area_P} {2 a},
\]
where $\area_P = 5\kappa\sigma$ is the area of a regular pentagon of
circumradius $1$.  In particular, if the inequality holds for
$a=\acrit$, then the density never exceeds
\[
\frac{\area_P}{2 \acrit} = \frac{5 - \sqrt{5}}{3},
\] %
the density of the pentagonal ice-ray.
\end{lemma}

For any finite set $\C$ of triangles, we will call the inequality
(\rif{eqn:main}) with the constant $a=\acrit$ the {\it main inequality}
(for $\C$).  We call the strict inequality,
\begin{equation}\libel{eqn:strict-main}
\frac{\sum_{T\in \C} b(T)}{\card(\C)} > \acrit,
\end{equation}
the {\it strict main inequality} (for $\C$).

\begin{proof} The maximum density can be obtained as the limit of the
  densities of a sequence of saturated periodic packings.  Thus, it is
  enough to consider the case when the packing is periodic.  A
  periodic packing descends to a packing on a flat torus
  $\ring{R}^2/\Lambda$, for some lattice $\Lambda$. We may assume that
  the lattice is large enough to avoid degenerate identifications, such
  as an identification of two vertices on a single Delaunay triangle.
  The rule defining
  $\N$ is translation invariant, and $\N$ descends to the torus.  On
  the torus, the set of pentagons, the set of triangles, and the set
  $\N$ are finite.  The equivalence relation $(\equiv_b)$ defining
  clusters is translation invariant, and each cluster is finite by
  assumption, so that no cluster contains both a triangle and a
  translate of the triangle under a nonzero element of $\Lambda$.
  Thus, each cluster $\C$ in $\ring{R}^2$ maps bijectively to a
  cluster $\C$ in the flat torus.  The functions $b$, $n_\pm$, $m_\pm$
  are the same whether computed on $\ring{R}^2$ or
  $\ring{R}^2/\Lambda$.  Let $p$ be the number of pentagons in the
  torus.  By the Euler formula for a torus triangulation, the number
  of Delaunay triangles is $2p$.  We have
\[
\sum_{T} n_-(T) =  \sum_{T} n_+(T) = \card(\N);\qquad
\sum_{T} m_-(T) =  \sum_{T} m_+(T) = \card(\M).
\]
Thus, the terms in $b(T)$ involving $n_+(T)$, $n_-(T)$, $m_+(T)$,
and $m_-(T)$ cancel:
\[
\area(\ring{R}^2/\Lambda) = \sum_T \area(T) = \sum_T b(T).
\]    
Let $\area_P$ be the area of a regular pentagon.  Making use of the
hypothesis of the lemma, we see that the density is
\[
\frac{p\, \area_P}{\sum_T \area(T)} 
=\frac{p\, \area_P}{\sum_T b(T)} \le\frac{p\, \area_P}{2 p\, a} 
= \frac{\area_P}{2\, a}.
\]
When $a=\acrit$, the term on the right is the density of the
pentagonal ice-ray, as desired.
\end{proof}

This article gives a proof of the following theorem. In view of
Lemma~\rif{lemma:main}, it implies the main result,
Theorem~\rif{thm:main}.  The proof of this result appears at the end
of Section~\rif{sec:nonobtuse}.

\begin{theorem}\libel{thm2:main}
  Let $\C$ be a cluster of Delaunay triangles in a saturated packing
  of regular pentagons.  Then $\C$ is finite and the average of $b(T)$
  over the cluster is at least $\acrit$.  That is, $\C$ satisfies the
  main inequality.  Equality holds exactly when $\C$ consists of two
  adjacent Delaunay triangles from the pentagonal ice-ray, attached
  along their common longest edge, forming an ice-ray dimer pair.
\end{theorem}

\begin{corollary}\label{cor:equal}
  In a periodic packing, the average of $b(T)$ over the set of
  all triangles modulo periods is equal to $\acrit$ exactly when
  the packing is the pentagonal ice-ray.
  Thus, the pentagonal ice-ray  is the
unique periodic packing that achieves maximal density.
\end{corollary}

\begin{proof}
  The set of triangles modulo periods is finite and is given as union
  of clusters.  Analyzing the proof of Lemma~\rif{lemma:main}, we see
  that for a periodic packing, the maximum density is achieved exactly
  when each cluster in the packing gives exact equality in the main
  inequality.  The average of $b(T)$ is equal to the average of
  $\area(T)$.  By the theorem, each cluster is an ice-ray dimer pair.
  Each dimer pair is a kite (a quadrilateral with two pairs of equal-length
  edges in
  which equal edges are adjacent).  The partition by
  Delaunay triangles is edge-to-edge.  This forces the tiling by kites
  to be edge-to-edge.  To match one kite edge with another, an adjacent
  kite must be reflected or rotated by 180.  However, since the
  edges of the kite do not lie along the axes of reflectional symmetry of pentagons,
  edge reflection fails to preserve
  the pentagons.  Hence, rotations must be used;
  all four kites adjacent
  to a given kite must be rotated by 180 degrees relative to the given
  kite.  By transitive extension, all kites in the tiling alternate in
  rotational orientation.  Then the pentagons align in columns in the
  ice-ray pattern.  This gives global uniqueness, up to isometry.
\end{proof}

\section{Pentagons in Contact}\libel{sec:contact}

\subsection{Notation}

By way of general notation, we use uppercase $A,B,C,\ldots$ for
pentagons; $\v_A,\v_B,\ldots$ for the vertices of pentagons;
$\c_A,\c_B,\ldots$ for centers of pentagons; $\p,\q,\ldots$ for
general points in the plane; $\normo{\p}$ for the Euclidean norm;
$\dx{AB}=\dd{A}{B}$ for center-to-center distances;
$\alpha,\beta,\gamma,\phi,\psi,\ldots$ for angles; and
$T,T',T_+,T_-,T_0,T_1,T_2,\ldots$ for Delaunay triangles.

We let $\eta(T) = \eta(d_1,d_2,d_3)$ be the circumradius of a triangle
$T$ with edge lengths $d_1,d_2$, and $d_3$.

Let $\angle(\p,\q,\r)$ be the angle at $\p$ of the triangle with
vertices $\p$, $\q$, and $\r$.  Note the nonstandard convention of
making the \emph{first} argument the vertex of the angle.  Let $\arc(d_1,d_2,d_3)$ be the angle
of a triangle (when it exists) with edge lengths $d_1$, $d_2$, and
$d_3$, where $d_3$ is the edge length of the edge opposite the
calculated angle.  We write $\area(T)= \area(d_1,d_2,d_3)$ for the
area of triangle $T$ with edge lengths $d_1,d_2,d_3$.

\subsection{Triple contact}

In this subsection, we describe possible contacts between pentagons.

We consider a single Delaunay triangle and the three nonoverlapping
pentagons centered at the triangle's vertices
(Figure~\rif{fig:3C-type}).  We call such a configuration a {\it
  $P$-triangle}.  A $P$-triangle is determined up to congruence by six
parameters: the lengths of the edges of the Delaunay triangle and the
rotation angles of the regular pentagons.  When we refer to the area
or edges of a $P$-triangle, we mean the area or edges of the
underlying Delaunay triangle.  More generally, we allow $P$-triangles
to inherit properties from Delaunay triangles, such as obtuseness or
nonobtuseness, the relation $(\rab)$. We can form clusters of $P$-triangles, and so forth. In
clusters of $P$-triangles it is to be understood that the pentagons
agree at coincident vertices of the triangles.  When there is a
fixed backdrop of a Delaunay triangulation of a pentagon packing, it
is not necessary to make a careful distinction between a $P$-triangle
and its underlying Delaunay triangle.

When two pentagons touch each other, some vertex of one meets an edge
of the other.  We call the pentagon with the vertex contact the {\it
  pointer} pentagon, and the pentagon with the edge contact the {\it
  receptor} pentagon (Figure~\rif{fig:receptor}).  We also call the
vertex in contact the {\it pointer vertex} of the pointer
pentagon. There are degenerate cases, when the contact set between two
pentagons contains of a vertex of both pentagons.  In these degenerate
cases, the designation of one pentagon as a pointer and the other as a
receptor is ambiguous.

\tikzfig{receptor}{Pointer and receptor pairs of pentagons.  
In each pair, the pentagon on the
left can be considered a pointer pentagon, 
with receptor on the right.  The first pair is nondegenerate,
and the other two pairs are degenerate.}{
[scale=0.5]
\pent{0}{0}{0}{1};
\pent{1.8}{0}{0}{1};
\pent{5}{0}{0}{1};
\pent{7}{0}{180}{1};
\pent{10+0.176}{0.243}{0}{1};  %
\pent{12-0.172}{-0.243}{180}{1};
}

We say that a $P$-triangle is $3C$ (triple contact),
if each of the three pentagons contacts the other two.

We may direct the edges of a $3C$ triangle by drawing an arrow from
the pointer pentagon to the receptor pentagon.  We may classify $3C$
triangles according to the types of triangles with directed edges.
There are two possibilities for the directed graph.
\begin{enumerate}
\item Some vertex of the triangle is a source of two directed edges
  and another vertex is the target of two directed edges
  ($LJ$-junction, $TJ$-junction or $\Delta$-junction).
\item Every vertex of the triangle is both a source and a target
  (pinwheel, pin-$T$).
\end{enumerate}

As indicated in parentheses, we have named each of the various contact
types.  An example of each of the contact types is shown in
Figure~\rif{fig:3C-type}.  An exact description of these contact types
appears later.  The \emph{inner polygon} of a $3C$ $P$-triangle refers to
the bounded region delimited by the three pentagons.  The inner
polygon is a $(3+n)$-gon, and has a nonconvex angle at $n$ vertices
(formed by freestanding pentagon vertices). The name $LJ$-junction is
suggested by the $L$-shaped inner polygon.  Similarly, the name
$TJ$-junction is suggested by the concave $T$-shaped inner polygon.
Similarly, the name $\Delta$-junction refers to the triangular inner polygon.
This section shows that the types in the figure exhaust the geometric
types of $3C$ contact.

\tikzfig{3C-type}{Types of $3C$-contact from top-to-bottom:
a pinwheel, a pin-$T$ junction, 
a $\Delta$-junction, an $LJ$-junction, and a $TJ$-junction.
In the first two types, arrows form a directed cycle. In the other
three types, both arrows point outward from the pentagon on top.
The second column gives \emph{distorted} representations of the
inner polygons.
In the pin-$T$, the inner polygon, which is
a nonconvex pentagon, is difficult to discern because of a 
necessarily small angle and
short edge.   
}%
{
[scale=0.8] 
\begin{scope}[yshift=20cm]
\threepentnoD{0.00}{0.00}{46.69}{0.82}{1.53}{218.09}{1.73}{0.00}{163.28};
\coordinate (A) at (0,0);
\coordinate (B) at (0.82,1.53);
\coordinate (C) at (1.73,0);
\draw[->] ($0.8*(A) + 0.2*(B)$) -- ($0.2*(A) + 0.8*(B)$);
\draw[->] ($0.8*(B) + 0.2*(C)$) -- ($0.2*(B) + 0.8*(C)$);
\draw[->] ($0.8*(C) + 0.2*(A)$) -- ($0.2*(C) + 0.8*(A)$);
\end{scope}
\begin{scope}[xshift=4.5cm,yshift=20cm,scale=0.8]
\draw[red] (0,0) -- ++ (30:1cm)  
-- ++ (150:1cm) node[black,anchor=south] {\hbox{pinwheel}} -- (0,0);
\end{scope}
\begin{scope}[yshift=17.5cm]
\threepentnoD{0.00}{0.00}{90}{1.40}{-1.377}{0}{-0.35}{-1.628}{-18.89};
\coordinate (A) at (0,0);
\coordinate (B) at (1.40,-1.377);
\coordinate (C) at (-0.35,-1.628);
\draw[<-] ($0.8*(A) + 0.2*(B)$) -- ($0.2*(A) + 0.8*(B)$);
\draw[<-] ($0.8*(B) + 0.2*(C)$) -- ($0.2*(B) + 0.8*(C)$);
\draw[<-] ($0.8*(C) + 0.2*(A)$) -- ($0.2*(C) + 0.8*(A)$);
\end{scope}
\begin{scope}[xshift=5.5cm,yshift=15.5cm,scale=0.8]
\draw[red] (0,0) -- ++ (90:2cm) 
-- ++ (245:0.5cm) 
-- node[anchor=south,black] {pin-$T$} ++ (180:1.5cm) -- ++ (350:1.2cm) -- (0,0);
\end{scope}
\begin{scope}[yshift=12cm]
\threepentnoD{0.00}{0.00}{-5.16}{0.99}{1.70}{235.38}{1.98}{0.00}{183.43};
\coordinate (A) at (0,0);
\coordinate (B) at (0.99,1.7);
\coordinate (C) at (1.98,0);
\draw[<-] ($0.8*(A) + 0.2*(B)$) -- ($0.2*(A) + 0.8*(B)$);
\draw[->] ($0.8*(B) + 0.2*(C)$) -- ($0.2*(B) + 0.8*(C)$);
\draw[->] ($0.8*(C) + 0.2*(A)$) -- ($0.2*(C) + 0.8*(A)$);
\end{scope}
\begin{scope}[xshift=5cm,yshift=12cm,scale=1.3]
\draw[red] (0,0) -- ++ (60:1cm)  
-- node[black,anchor=south] {$\Delta$} ++ (180:1cm)  -- (0,0);
\end{scope}
\begin{scope}[yshift=8cm]
\threepentnoD{0.00}{0.00}{80.86}{0.97}{1.58}{232.21}{1.82}{0.00}{223.24};
\coordinate (A) at (0,0);
\coordinate (B) at (0.97,1.58);
\coordinate (C) at (1.82,0);
\draw[<-] ($0.8*(A) + 0.2*(B)$) -- ($0.2*(A) + 0.8*(B)$);
\draw[->] ($0.8*(B) + 0.2*(C)$) -- ($0.2*(B) + 0.8*(C)$);
\draw[<-] ($0.8*(C) + 0.2*(A)$) -- ($0.2*(C) + 0.8*(A)$);
\end{scope}
\begin{scope}[xshift=5cm,yshift=8cm,scale=0.7]
\draw[red] (0,0) -- ++ (95:0.8cm)  
-- ++ (30:1cm) 
-- node[black,anchor=south] {\hbox{LJ}} ++ (180:2cm) 
-- (0,0);
\end{scope}
\begin{scope}[yshift=4cm]
\threepentnoD{0.00}{0.00}{114.48}{0.90}{1.59}{237.18}{1.66}{0.00}{219.24};
\coordinate (A) at (0,0);
\coordinate (B) at (0.9,1.59);
\coordinate (C) at (1.66,0);
\draw[<-] ($0.8*(A) + 0.2*(B)$) -- ($0.2*(A) + 0.8*(B)$);
\draw[->] ($0.8*(B) + 0.2*(C)$) -- ($0.2*(B) + 0.8*(C)$);
\draw[<-] ($0.8*(C) + 0.2*(A)$) -- ($0.2*(C) + 0.8*(A)$);
\end{scope}
\begin{scope}[xshift=5cm,yshift=3.5cm,scale=0.65]
\draw[red] (0,0) -- ++ (90:1.8cm)  
-- ++ (20:1.2cm) 
-- node[black,anchor=south] {\hbox{TJ}} ++ (180:2.5cm) 
-- ++ (345:1.2cm)
-- (0,0);
\end{scope}
}

A {\it cloverleaf} arrangement is a $3C$ triangle that has a point at
which vertices from all three pentagons meet
(Figure~\rif{fig:clover}).  This is degenerate because this shared
vertex can be considered as a pointer or receptor.

\tikzfig{clover}{A cloverleaf (degenerate pinwheel).  The Delaunay triangle
in this particular example is not subcritical.}
{
[scale=0.6]
\begin{scope}[xshift=4cm]
\threepentnoD{0.00}{0.00}{31.70}{0.76}{1.52}{203.11}{1.70}{0.00}{148.30};
\end{scope}
}

In general, in this article, a {\it non-anomaly} lemma refers to a
geometrical lemma that shows that certain geometric configurations are
impossible.  Generally, it is obvious from the informal pictures that
various configurations cannot exist.  The non-anomaly lemmas then
translate the intuitive impossibilities into mathematically precise
statements.  We give a few non-anomaly lemmas as follows.
They are expressed as separation results, asserting that
two pentagons $A$ and $C$ do not touch.

\begin{lemma}\libel{lemma:sep1} 
Let $T$ be a $3C$-triangle with pentagons $A$, $B$, and $C$ such that
$B$ is a pointer to both of the other pentagons $A$ and $C$.  
Assume that $T$ is not a cloverleaf.
Then the two pointer vertices $\v_B$ and $\v_B'$ are adjacent
vertices of $B$.
\end{lemma}

\begin{lemma}\libel{lemma:sep2}  
Let $T$ be a $3C$-triangle with pentagons $A$, $B$, and $C$.  Suppose
that pentagon $A$ is a pointer to $B$ at $\v_A$ and that $B$ is a
pointer to $C$ at $\v_B$.  Then some edge of $B$ containing $\v_A$ is
equal or adjacent to some edge of $B$ containing $\v_B$.
\end{lemma}

\begin{lemma}\libel{lemma:sep3} 
  Let $T$ be a $3C$-triangle with pentagons $A$, $B$, and $C$ such
  that $B$ is a receptor of both of the other pentagons.  Then the two
  pointer vertices $\v_A$ and $\v_C$ lie on the same edge or adjacent
  pentagon edges of $B$.
\end{lemma}

\tikzfig{sep}{A line through the center 
of the middle pentagon $B$
through one of its vertices 
separates the two extremal 
pentagons $A$ and $C$.}{
[scale=0.5]
\pen{0}{0}{90};
\pen{1.72}{0.56}{90};  %
\pen{-1.72}{0.56}{90};  
\draw (0,-1.5) -- (0,1.5) node[black,anchor=west] {$\v$};
\begin{scope}[xshift=6cm]
\pen{0}{0}{90};
\pen{1.72}{0.56}{90};  
\pen{-1.72}{-0.56}{90};  
\draw (0,-1.5) -- (0,1.5) node[black,anchor=west] {$\v$};
\end{scope}
\begin{scope}[xshift=12cm]
\pen{0}{0}{90};
\pen{1.72}{-0.56}{90};  
\pen{-1.72}{-0.56}{90};  
\draw (0,-1.5) -- (0,1.5) node[black,anchor=west] {$\v$};
\end{scope}
}

\begin{proof} The Lemmas~\rif{lemma:sep1}, \rif{lemma:sep2},
  and~\rif{lemma:sep3} can be proved in the same way.  In each case,
  we prove the contrapositive, assuming the negation of the geometric
  conclusion, assuming that $B$ has the indicated contacts with $A$ and $C$,
  and proving that the configuration is not $3C$.  We show
  that the configuration is not $3C$ by constructing a separating
  hyperplane between the pentagons $A$ and $C$.  In each case, the
  separating hyperplane is a line through the center of the middle
  pentagon $B$ and passing through a vertex $\v$ of that pentagon.
  See Figure~\rif{fig:sep}.
  In the case of Lemma~\rif{lemma:sep1}, assuming the two pointer
  vertices are not adjacent,
   the vertex $\v$ is the vertex adjacent to both
  pointer vertices of $B$.
  There is a degenerate case of a cloverleaf, where all three
  pentagons meet at the vertex $\v$ on the separating line.
  Among the three cases, this is the only degeneracy.
  In the case of Lemma~\rif{lemma:sep2}, the vertex $\v$ can be
  chosen to any vertex other than three exclusions: the two endpoints
  of the edge of $B$ containing $\v_A$ and the vertex $\v_B$.  The
  pentagon $C$ can meet the separating line, but the pentagon $A$ must
  lie in the open half-plane.
  In the case of Lemma~\ref{lemma:sep3}, the contrapositive assumption
  is that the two pointer vertices lie in the relative interiors of two
  nonadjacent edges.  Excluding the four endpoints of these two edges,
  the vertex $\v$ is chosen to be the remaining vertex of $B$. The two pentagons
  lie in opposite open halfplanes formed by the separating line.
\end{proof}

\begin{definition}[$\Delta$]
We say that a $3C$-triangle has type $\Delta$ if we are in the first
  case of Lemma~\rif{lemma:sep3} (that is, the two pointer vertices
  $\v_A$ and $\v_C$ of $A$ and $C$ lie on the same edge of $B$) and if
  the line $\lambda$ through that edge of $B$ separates $B$ from $A$
  and $C$. (See Figure~\rif{fig:delta}.)
\end{definition}

\tikzfig{delta}{In type $\Delta$, a line separates pentagon 
$B$ from
the other two pentagons. 
The second figure (which is degenerate of type $LJ$) does not have
type $\Delta$.}{
[scale=0.5]
\pen{0}{0}{90};
\draw (0,0) node {$B$};
\draw (-2.0,-0.809) node[above] {$\lambda$} --  (2.0,-0.809);
\pen{-1.05}{-2*0.809}{-90};
\pen{1.9 - 1.05}{-2*0.809}{-90};
\begin{scope}[xshift=6cm]
\pen{0}{0}{90};
\draw (0,0) node {$B$};
\draw (-2.0,-0.809) node[above] {$\lambda$} --  (2.0,-0.809);
\pen{-0.428}{-2*0.809}{-90};
\pen{1.36}{ -1.44}{ 212.704};
\end{scope}
}

In type $\Delta$, say $A$ is a pointer into $C$ at $\v$.  Then $\v_A$
and $\v$ are the two endpoints of some edge of $A$.  Also, $\v_C$ and
$\v$ lie on the same edge of $C$.  If the line $\lambda$ does not
separate $B$ from $A$ and $C$, then $\v_C$ is a shared vertex of $B$
and $C$, and we have a degeneracy that can also be viewed as $\v_A$
and $\v_C$ on adjacent pentagon edges of $B$.  This case will be
classified as a degenerate $LJ$-junction below.

\begin{definition} Let $T$ be a $P$-triangle with pentagons $A$, $B$,
  and $C$.  Assume that $A$ points to $B$ at $\v_{AB}$, and $B$ points
  to $C$ at $\v_{BC}$.  An {\it inner vertex} $\v$ of $B$ is a vertex
  $\v\ne \v_{BC}$ of $B$ such that $\v$ lies between $\v_{AB}$ and
  $\v_{BC}$ (along the short run of the perimeter of $B$ from
  $\v_{AB}$ to $\v_{BC}$).  We allow the degeneracy $\v=\v_{AB}$.
  See Figure~\rif{fig:pin1}.
\end{definition}

\begin{definition} Let $T$ be a $3C$-triangle.  Fix the pointer
  directions on $T$, if ambiguous. We say $T$ has type {\it pin}-$k$,
  for $k\in \{0,1,2,3\}$ if $A$ points into $B$, $B$ points into $C$,
  $C$ points into $A$, and if there are exactly $k$ pentagons among
  $A,B,C$ that have an inner vertex.  We use {\it pinwheel} as a
  synonym for pin-$0$ and {\it pin-T} as a synonym for pin-$2$.
\end{definition}

\begin{lemma} There does not exist a $3C$-triangle $T$ of type pin-$1$
  with pentagons $A$, $B$, and $C$.
\end{lemma}

\tikzfig{pin1}{A distorted pin-$1$ configuration.}{
\begin{scope}[scale=1.0]
\pen{0}{0}{-90};
\pen{-0.974}{ 1.543}{ 32.70422};
\draw[red] ++ (-0.974,1.543) 
  ++ (0.842,0.532) node[black,anchor=west] {$\v_{BC}$} -- 
  ++ (115:0.2) -- ++ (-65:1.6) node[black,anchor=south west] {$\v_{CA}$}
   -- ++ (7:1.17);
\draw (0,0) node {$A$};
\draw (0,0) + (3*72-90:1) node[anchor=north west] {$\v_{AB}$};
\draw (-0.974,1.543) node {$B$};
\draw (1.1,1.543) node {$C$};
\smalldot{3*72-90:1}; %
\draw (-0.974,1.543)+(32.7-72:1) node[anchor=south west]{$\v$};
\smalldot{$(-0.974,1.543)+(32.7-72:1)$}; %
\smalldot{$(-0.974,1.543)+(32.7:1)$}; %
\smalldot{$(-0.974,1.543)+(0.842,0.532)+(115:0.2)+(-65:1.6)$}; %
\end{scope}
}

\begin{proof} 
  For a contradiction, we draw a (distorted) picture of a pin-$1$
  configuration (Figure~\rif{fig:pin1}).  We let $\v$ be the inner
  vertex of $B$; that is, the vertex that is interior to the triangle
  $(\v_{CA},\v_{AB},\v_{BC})$ with vertices at the pointers $X$ into
  $Y$, where $(X,Y)=(A,B),(B,C),(C,A)$.  The inner vertex is an endpoint of the edge of the pentagon $B$ containing
  $\v_{AB}$.  We have
\[
\angle (\v,\v_{AB},\v_{CA}) \le \pi,\quad 
\angle(\v,\v_{BC},\v_{AB}) = 3\pi/5,\quad
\angle(\v_{CA},\v_{BC},\v_{AB})\le 2\pi/5.
\]
(The last inequality uses the fact that $T$ is not a degenerate
pin-$2$, so that $\v_{CA}$ is not a vertex of $A$.)  We also have
\begin{align*}
  \angle(\v,\v_{BC},\v_{CA}) &= 2\pi - \angle(\v,\v_{AB},\v_{BC}) - \angle(\v,\v_{AB},\v_{CA})\\
  &\ge 2\pi /5 
\ge \angle(\v_{CA},\v_{BC},\v_{AB}) \ge \angle(\v_{CA},\v_{BC},\v).
\end{align*}
The law of sines applied to the triangle $(\v,\v_{BC},\v_{CA})$ then
gives
\[
2\sigma =\norm{\v_{BC}}{\v}\le\norm{\v_{BC}}{\v_{CA}}\le 2\sigma.
\]
Thus, we have equality everywhere.  In particular,
$\angle(\v,\v_{AB},\v_{CA})=\pi$, forcing $\v$ to lie on an edge of $A$,
and $\v_{BC}$ to be a vertex of $C$.
Hence $\v_{BC}$ is a degenerate inner vertex of $C$, and $T$ has type
pin-$k$, for some $k\ge 2$.
\end{proof}

\begin{lemma}  The type pin-$3$ does not exist.  
\end{lemma}

\begin{proof} 
  Suppose for a contradiction that a $P$-triangle $T$ of type pin-$3$
  exists.  The region $X$ bounded by the three pentagons is a
  nonconvex star-shaped hexagon, with interior angles $\alpha'$,
  $7\pi/5$, $\beta'$, $7\pi/5$, $\gamma'$, and $7\pi/5$.  The vertices
  of $X$ with angles $7\pi/5$ are the inner vertices of the three
  pentagons of $T$.  The sum of the interior angles in a hexagon is
  $4\pi$:
\[
4\pi = \alpha'+\beta'+\gamma' + 3 (7\pi/5),
\]
which implies that $\alpha'+\beta'+\gamma' = -\pi/5$, which is
impossible.
\end{proof}

\begin{definition}[$TJ$ and $LJ$-junction]
  We say that a $3C$-triangle is a type $TJ$- or $LJ$-junction if it
  is not type $\Delta$ and if we are in the second case of
  Lemma~\rif{lemma:sep3} (both $A$ and $C$ point into $B$, and the two
  pointer vertices $\v_A$ and $\v_C$ lie on adjacent edges of $B$).
  Say $A$ is a pointer into $C$ at $\v$.  We say that it has type {\it
    $LJ$-junction} if $\v_C$ and $\v$ lie on the same pentagon edge of
  $C$, and otherwise we say it has type {\it $TJ$-junction}.
\end{definition}

We can be more precise about the structure of a $TJ$-junction.  In the
context of the definition, Lemma~\rif{lemma:sep2} implies that $\v$
and $\v_C$ lie on adjacent pentagon edges of $C$.  An $LJ$-junction
has one inner vertex, and the inner polygon is a nonconvex quadrilateral.
A $TJ$-junction has two inner vertices, and the inner polygon is a
nonconvex pentagon.

This completes the classification of $3C$-triangles: $\Delta$,
pinwheel, pin-$T$, $LJ$, and $TJ$.  Useful coordinate systems for the
various types can be found in the appendix
(Section~\rif{sec:appendix}).

\section{Delaunay Triangle Areas}\libel{sec:area}

As an application of the classification from the
Section~\rif{sec:contact}, this section makes a computer interval-arithmetic
calculation of a lower bound on the longest edge length of a
subcritical triangle.  We also obtain a lower bound on the area of a
nonobtuse Delaunay triangle.

Some proofs will deform $P$-triangles in a way that is monotonic in
relevant quantities such as area.  Whenever, we make deformations, it
is to be understood that the deformations are local; a single
$P$-triangle (or otherwise identified finite cluster) is deformed, and
we do not attempt to extend the deformation globally to the full
packing. The exact constraints that a deformation must satisfy will
depend on the context, but at the very least, a deformation must not
lead to overlapping pentagons. Except where indicated otherwise,
deformations are required to maintain combinatorial structures such
as cluster structure, arrows $(\rab)$ $(\Ra)$, membership of pairs in
$\M$ and $\N$, combinatorially relevant lengths comparisons $1.72$,
dimer and pseudo-dimer structures.  The specific Delaunay triangle
conditions that a deformation must preserve (circumradius at most $2$,
empty circumcircle constraint) will be specified in a case-by-case
basis.  We specify the side, circumradius, and angle constraints that
are retained in each such problem.  The side bound $2\kappa$ follows
from nonoverlap.  Pentagons at a common vertex are the same pentagon.

\begin{lemma}\libel{lemma:21} 
  A nonobtuse subcritical Delaunay triangle has edge lengths at most $2.1$.
\end{lemma}

\begin{proof} By the monotonicity of area as a function of edge length
  for nonobtuse triangles, a triangle with an edge length at least
  $2.1$ has area at least
\[
\area(2.1,2\kappa,2\kappa) > \acrit,
\] %
which is not subcritical.
\end{proof}

\begin{lemma}\libel{lemma:right}
  The edge opposite an obtuse angle of a Delaunay triangle has length greater than $\kappa\sqrt8$.
  A nonobtuse subcritical triangle has edge lengths less than
  $\kappa\sqrt8$.  In particular, a right-angled Delaunay triangle is
  not subcritical.
\end{lemma}

\begin{proof}  
  This is a corollary of the previous lemma, because $2.1 <
  \kappa\sqrt8 \approx 2.288$.
  A right triangle has legs each of length at least $2\kappa$ and hence
  hypotenuse length at least $\kappa\sqrt8$.
\end{proof}

\begin{remark}
  A motion of a pentagon in the plane can be described by an element
  of the isometry group of the plane, which is a semidirect product of
  a translation group and an orthogonal group.  Because of the
  dihedral symmetries of the regular pentagon, each motion can be
  realized as a translation followed by a rotation by angle between
  $0$ and $2\pi/5$.  In particular, a translation of a pentagon is a
  motion of a pentagon such that the rotational part is the identity.
\end{remark}

\begin{definition}
In a $P$-triangle, we say that a pentagon $A$ has {\it primary
  contact} if one or more of the following three conditions hold:
\begin{enumerate}
\item (slider contact) The pentagon $A$ and one of the other two pentagons
   share a positive length edge segment;
\item (midpointer contact) A vertex of one of the other two pentagons
  is the midpoint of one of the edges of the pentagon $A$; or
\item (double contact) The pentagon $A$ is in contact with both of the
  other pentagons.
\end{enumerate}
\end{definition}

The next lemma is used to give area estimates when an edge has length
at most $\kappa\sqrt{8}$.  We give two forms of the lemma.  We prove
them together.

\begin{lemma}\libel{lemma:primary} 
  Let $A$ be a pentagon in a nonobtuse $P$-triangle.  Assume that the
  triangle edge opposite $\c_A$ has length at most $\kappa\sqrt{8}$.
  Keeping $B$ and $C$ fixed, we can choose a position and orientation
of $A$ such that $A$ has primary contact, the triangle is nonobtuse,
and neither its area nor any of its edge lengths exceeds the
corresponding original value.
\end{lemma}

\begin{lemma}\libel{lemma:primary2} 
  Let $A,B,C$ be pentagons in a nonobtuse $P$-triangle.  Assume that
  the triangle edges opposite $\c_A$ and $\c_C$ have length at most
  $\kappa\sqrt{8}$.
  Keeping $B$ and $C$ fixed, we can choose a position and orientation
of $A$ such that $A$ has primary contact, the triangle is nonobtuse,
$\dx{AB}$ has its original value, and the area does not exceed
its original value.
\end{lemma}

  \begin{proof}
Fix $B$ and $C$. Consider all positions and orientations of $A$
for which the pentagons have pairwise disjoint interiors, the
triangle is nonobtuse, and its area is at most the original area.
Require $\c_A$ to lie on the original side of the line through
$\c_B,\c_C$. Also impose the stated constraints of the lemma.
This set is nonempty and compact. Indeed, nonobtuseness
and the area constraint bound the admissible region for $\c_A$,
and the constraints are closed. Choose a
configuration of least area.

We claim that $A$ has primary contact. Otherwise, $A$ contacts
at most one of the fixed pentagons, and any such contact is a
single point. A small rotation about $\c_A$ can break this
contact. Thus, without changing the centers or the area, we may arrange
that $A$ is separated from both $B$ and $C$.

In the first lemma, translate $A$ a small distance towards the
line through $\c_B,\c_C$, perpendicular to that line. This
strictly decreases the area and both incident edge lengths.
The angles at $\c_B$ and $\c_C$ remain nonobtuse. The angle at
$\c_A$ is acute: if it were right, the bound
$\dx{BC}\le\kappa\sqrt8$ and the lower bounds
$\dx{AB},\dx{AC}\ge2\kappa$ would force
$\dx{AB}=\dx{AC}=2\kappa$, giving primary contact.
Consequently a sufficiently small translation preserves
nonobtuseness and nonoverlap, contradicting minimality.

In the second lemma, move $\c_A$ a small distance along the
circle centered at $\c_B$, towards the line through
$\c_B,\c_C$, keeping the orientation of $A$ fixed. This preserves
$\dx{AB}$ and strictly decreases the area. The angle at $\c_B$
decreases, and the angle at $\c_A$ is acute by the preceding
argument. The angle at $\c_C$ is also acute: a right angle there
would give
\[
\dx{AB}^2=\dx{BC}^2+\dx{AC}^2\ge8\kappa^2.
\]
The upper bound on $\dx{AB}$ would then force
$\dx{AC}=2\kappa$, giving slider contact between $A$ and $C$.
Again, a sufficiently small displacement preserves all the
constraints and contradicts minimality.
Thus in both lemmas, the minimizing configuration has primary contact.
\end{proof}

% COMPUTATION QBQQTVY | lemma:delta | Delta-junction area
\begin{lemma}\libel{lemma:delta} 
  A subcritical $3C$-triangle does not have type $\Delta$.  In fact,
  such a $P$-triangle $T$ has area greater than $1.5$.
\end{lemma}

\begin{proof} The proof uses computer interval arithmetic. The $3C$-triangles of
  type $\Delta$ form a three-dimensional configuration space.  The
  appendix (Section~\rif{sec:appendix}) introduces good coordinate
  systems for each of the various $3C$-triangle types.  We make a
  computer interval-arithmetic calculation of the area of the Delaunay triangle as a
  function of these coordinates.
\end{proof}

\begin{lemma}\libel{lemma:mid-172}  
  If a pentagon $A$ has midpointer contact with a pentagon $B$, then
  $\dx{AB} > 1.72$.
\end{lemma}

\begin{proof} Suppose a pointer vertex $\v_A$ of $A$ is the midpoint of an
  edge of pentagon $B$.  Rotating $A$ about the vertex $\v_A$, keeping
  $B$ fixed, we may decrease $\dx{AB}$ until $A$ and $B$
  have slider contact.  This determines the configuration of $A$ and
  $B$ up to rigid motion.  By the Pythagorean theorem, the
  distance between pentagon centers is
\[
\dx{AB} = \sqrt{(2\kappa)^2 + \sigma^2} \approx 1.72149 > 1.72.
\] %
\end{proof}

% COMPUTATION BVICHFW | lemma:172 | Longest edge of a subcritical triangle
\begin{lemma}\libel{lemma:172}
  If every edge of a $P$-triangle $T$ is at most $1.72$, then the
  triangle is not subcritical.
\end{lemma}

\begin{proof} This is a computer interval-arithmetic assisted proof.\footnote{The constant
    $1.72$ is nearly optimal.  For example, in the notation of the
    appendix, the pinwheel with parameters $\alpha=\beta=\pi/15$,
    $x_\gamma = 0.18$ is subcritical equilateral with edge lengths
    approximately $1.72256$.} Such a triangle is nonobtuse.  After
  fixing an overall rigid motion, minimize the area over all configurations
  whose three edge lengths do not exceed their original values. The
  admissible set is compact. If a pentagon does not have primary contact,
  the local contact-breaking rotation and height-decreasing motion in the
  proof of Lemma~\rif{lemma:primary} strictly decrease the area while
  preserving these bounds. Thus each pentagon has primary contact at a
  minimum.
  By the previous lemma, we may assume that the
  contact is not midpointer contact.  Thus, each pentagon has double
  contact or slider contact with the other pentagons.

  The pentagon contact graph of the $P$-triangle does not have an isolated vertex, so some
  vertex of the contact graph has degree two. This vertex $\c_A$ of the
  contact graph corresponds to a pentagon $A$ with double contact.  If the
  $P$-triangle does not have $3C$ contact, then the contact graph is
  acyclic, and the other two pentagons have slider contact
  (Figure~\rif{fig:172-slider}).

  The centers $\c_B$, $\c_C$ lie in a disk at $\c_A$ of radius $1.72$ and
  on the pentagon obtained by a homothety of $A$ of factor $2$.  These
  constraints, together with $2\kappa\le d_{BC}\le1.72$, force $\c_B$
  and $\c_C$ to lie on adjacent edges of the dilated copy of $A$.  (If
  $\c_B$ and $\c_C$ lie on the same edge the pentagons $B,C$ overlap, and if they
  lie on nonadjacent edges the distance satisfies $d_{BC}>1.72$.) Let
  $e_B$, $e_C$ be the parts of these dilated edges inside the disk of
  radius $1.72$ at $\c_A$.  Bisect $e_B$, $e_C$ into $e_B',e_B''$,
  $e_C',e_C''$ with $e_B''$ $e_C''$ the halves that are closer to each
  other.  If $\c_B\in e_B''$ and $\c_C\in e_C''$, then the pentagons $B$
  and $C$ intersect at a vertex of $A$, and the contact is $3C$ (or
  they overlap).  Hence we may assume without loss of generality that
  $\c_B\in e_B'$.  Sliding $\c_B$ along $e_B'$ in the direction of the
  midpoint of $e_B$ is monotonic in all edge lengths and decreases
  area.  Similarly, if $\c_C\in e_C'$, it can slide in a similar
  manner.  We continue to slide until $3C$ contact.

  \tikzfig{172-slider}{The first frame shows pentagon $A$ and segments (on a dilated pentagon) inside
    a circle of radius $1.72$ where centers $\c_B$ and $\c_C$ can be placed in
    slider contact with pentagon $A$.  The second frame shows sliding pentagons $B$
    centered at $\c_B\in e'_B$ and $C$ centered at $\c_C\in e''_C$.
  }{
\begin{scope}[scale=0.5,xshift=8cm]
  \threepent{0}{0}{-90}{0.63}{-1.54}{90}{-1.27}{-1.07}{90};
    \begin{scope}
    \clip (0,0) circle[radius=1.72cm];
  \end{scope}
\end{scope}
\begin{scope}[scale=0.5]
  \pent{0}{0}{-90}{1};
  \begin{scope}
    \clip (0,0) circle[radius=1.72cm];
    \pent{0}{0}{-90}{2};
  \end{scope}
  \draw[blue] (0,0) circle[radius=1.72cm];
\end{scope}
}

Now assume that the $P$-triangle has $3C$ contact.  We have classified
all $3C$ triangles.  We obtain the proof by expressing each type of
triangle in terms of explicit coordinates from the appendix
(Section~\rif{sec:appendix}) and computing bounds on the areas and
edge lengths of the triangles using interval arithmetic.  The result
follows.
\end{proof}

\begin{lemma}\libel{lemma:2C} 
  Let $T$ be a subcritical nonobtuse $P$-triangle with pentagons $A$,
  $B$, and $C$. Keeping $B$ and $C$ fixed, we may choose a position
  and orientation of $A$ of no greater area such that $A$ has double
  contact with $B$ and $C$ and the triangle remains nonobtuse.
\end{lemma}

\begin{proof}
  Fix $B$ and $C$, and choose a position and orientation of $A$ that
  minimizes area among the nonobtuse configurations of no greater area,
  with $A$ on the original side of the line through $\c_B,\c_C$.
  As in the proof of Lemma~\rif{lemma:primary}, this set is compact,
  and a minimum has primary contact. By Lemma~\rif{lemma:right}, a
  subcritical nonobtuse triangle is acute, so no right-angle boundary
  obstructs a sufficiently small area-decreasing motion.

  Suppose that $A$ does not have double contact. If its primary
  contact is slider contact, slide $A$ along the contact edge in the
  direction to decrease its height above the fixed base. A horizontal
  contact edge would put $\c_A$ at height $2\kappa$, giving area at
  least $2\kappa^2>\acrit$. Thus the slide strictly decreases area.

  If the primary contact is midpointer contact, rotate $A$ about the
  pointer vertex of the fixed pentagon. Its center moves on a circle
  of radius $\kappa$. Unless the center is at the lowest point of
  this circle, a sufficiently small rotation decreases its height.
  If that lowest point lies above the fixed base, it lies inside
  the circumdisk of the fixed pointer pentagon, contradicting the
  lower bound $2\kappa>1$ on center separation. It therefore cannot
  be a minimum either.

  In both cases a sufficiently small motion preserves nonoverlap
  with the other fixed pentagon and contradicts minimality. Hence
  $A$ has double contact.
\end{proof}

% COMPUTATION QGPVOYX | lemma:a0 | Minimum nonobtuse triangle area
\begin{lemma}\libel{lemma:a0}  
A nonobtuse $P$-triangle $T$ has area greater than $\ao$.
\end{lemma}

\begin{proof} 
  This is a computer interval-arithmetic assisted proof.  We may assume for a contradiction
  that $T$ has area at most $\ao$.  In particular, it is
  subcritical. Modulo rigid motions, choose a minimum over all
  nonobtuse three-pentagon configurations of no greater area. The
  subcritical edge bounds make this set compact. The local
  strict-decrease argument in the proof of Lemma~\rif{lemma:2C}
  shows that every pentagon has double contact at this minimum.
  Thus $T$ is $3C$.  We have classified all $3C$ triangles.  We obtain the proof
  by expressing each type of triangle in terms of explicit
  coordinates from the appendix (Section~\rif{sec:appendix}) and
  computing bounds on the areas and edge lengths of the triangles
  using interval arithmetic.  The result follows.
\end{proof}

\section{Computer Calculations}\libel{sec:calc}

We make use of the following lemmas, which are discussed in the
appendix and proved by computer interval arithmetic.  (Although
further discussion appears at the end of the article, there is no
circular reasoning involved in using those calculations here.)

% COMPUTATION NKQNXUN | calc:pseudo1 | Pseudo-dimer shared edge
\begin{lemma}\libel{NKQNXUN} \libel{calc:pseudo1} 
Let $(T_1,T_0)\in \PD$.  The edge shared between
$T_0$ and $T_1$ has length less than $1.8$.
  The conclusion also holds if the third defining condition of $\PD$ is
  weakened to allow the shared edge to be equal in length to a nonshared edge
  of $T_0$.
\end{lemma}

\begin{proof} The proof uses computer interval arithmetic. See Section~\rif{sec:pseudo1}.
\end{proof}

% COMPUTATION RWWHLQT | calc:pseudo2 | Pseudo-dimer egressive edge
\begin{lemma}\libel{RWWHLQT}\libel{calc:pseudo2} 
  Let $(T_1,T_0)\in \PD$.  The longest edge of $T_0$ (that is, its
  egressive edge) has length greater than $1.8$.
  The conclusion also holds if the third defining condition of $\PD$ is
  weakened to allow the shared edge to be equal in length to a nonshared edge
  of $T_0$.
\end{lemma}

\begin{proof} The proof uses computer interval arithmetic. See Section~\rif{sec:pseudo2}.
\end{proof}

\begin{lemma}\libel{calc:pseudo3}  
  Let $(T_1,T_0)\in \PD$.  The two edges of $T_0$ other than the
  longest edge have lengths less than $1.8$.
  The conclusion also holds if the third defining condition of $\PD$ is
  weakened to allow the shared edge to be equal in length to a nonshared edge
  of $T_0$.
\end{lemma}

\begin{proof}
  The shared edge between $T_0$ and $T_1$ has length less than $1.8$
  by Lemma~\rif{calc:pseudo1}.  It has length at least $1.72$ by
  Lemma~\rif{lemma:172}.  If (for a contradiction) the third edge of
  $T_0$ has length at least $1.8$, then by the previous lemmas, its
  three edges have lengths at least $1.72$, $1.8$, $1.8$.  
By monotonicity of area in edge lengths for acute triangles, a calculation gives
\[
\area\{T_1,T_0\} > \ao + \area(1.72,1.8,1.8) 
> \ao + (2\acrit -\ao) = 2\acrit.
\]
This area inequality contradicts the subcritical property of pseudo-dimers.
\end{proof}

% COMPUTATION BXZBPJW | calc:pseudo-area | Pseudo-dimer pair area
\begin{lemma}\libel{BXZBPJW}\libel{calc:pseudo-area}
Let $(T_1,T_0)\in\PD$.  Then $\area\{T_0,T_1\} > 2\acrit - \epsM$.
\end{lemma}

\begin{proof} The proof uses computer interval arithmetic. See Section~\rif{sec:pseudo-area}.
\end{proof}

% COMPUTATION JQMRXTH | calc:pseudo-area3 | Pseudo-dimer plus egressive neighbor
\begin{lemma}\libel{JQMRXTH}\libel{calc:pseudo-area3}
Let $(T_1,T_0)\in\PD$. Assume $T_0\Ra T_-$, with $T_-$ nonobtuse.
Then $\area\{T_0,T_1,T_-\} > 3\acrit + \epsM$.
\end{lemma}

\begin{proof} The proof uses computer interval arithmetic. See Section~\rif{sec:pseudo-area3}.
\end{proof}

\begin{corollary}\libel{lemma:egress'} 
Let $(T_1,T_0)\in\PD$.  Assume that $T_0\Ra T_-$, with $T_-$ nonobtuse.  Then
$\area(T_-) > \acrit+\epsM$.
\end{corollary}

\begin{proof}  
  By Lemma~\rif{calc:pseudo-area3} and the definition of pseudo-dimer,
\[
\area(T_-) = \area\{T_0,T_1,T_-\} - \area\{T_0,T_1\} 
> (3\acrit + \epsM) - 2\acrit = \acrit + \epsM.
\]
\end{proof}

\begin{corollary}\libel{lemma:m2-area}
  Suppose that $m_-(T_-)>0$.  Then $\area(T_-) > \acrit + \epsM$.
\end{corollary}

\begin{proof} 
  If $m_-(T_-)>0$, there exists $(T_1,T_0)\in\PD$ such that $T_0\Ra
  T_-$.  The result follows from the previous corollary.
\end{proof}

\begin{definition}[long isosceles]
  We say that a triangle is {\it long isosceles} if the two longest edges of
  the triangle have equal length.  We include equilateral triangles as
  a special case of long isosceles.
\end{definition}

The isosceles constraint occurs because of the way we have defined
certain structures (such as dimer pairs) in terms of a longest edge of a
triangle.  Isosceles triangles occur naturally as boundary cases in
such structures.

\begin{definition}[O2C]
  Let $T_0$ and $T_1$ be a pair of $P$-triangles that share an edge.
  We say that a triangle $T=T_0$ or $T=T_1$ 
  has outside double contact $(O2C)$ if the pentagon
  at the vertex of $T$ that is not shared with the other triangle in the
  pair has double contact.
\end{definition}

% COMPUTATION KUGAKIK | calc:dimer-isosc | No long-isosceles O2C peripheral in a dimer
\begin{lemma}\libel{KUGAKIK}\libel{calc:dimer-isosc}
  Let $(T_1,T_0)\in DP$.  Then $T_1$ is not both $O2C$ and long
  isosceles.
\end{lemma}

\begin{proof} The proof uses computer interval arithmetic. See Section~\rif{sec:dimer-isosc}.
\end{proof}

% COMPUTATION HUQEJAT | calc:pent3 | Three-triangle cluster bound
\begin{lemma}\libel{HUQEJAT}\libel{calc:pent3}
  Let $T_0,T_1^0,T_1^1$ be distinct
nonobtuse $P$-triangles. Assume that $T_1^i\Ra T_0$ and
$\area(T_1^i)\le\acrit$ for $i=0,1$. Then
\[
\area\{T_0,T_1^0,T_1^1\}>3\acrit+\epsM.
\]
\end{lemma}

\begin{proof} The proof uses computer interval arithmetic.  See Section~\rif{sec:pent3}.
\end{proof}

% COMPUTATION QPJDYDB | calc:pent4 | Four-triangle cluster bound
\begin{lemma}\libel{QPJDYDB}\libel{calc:pent4}
Let $T_0,T_1^0,T_1^1,T_1^2$ be distinct
nonobtuse $P$-triangles. Assume that $T_1^i\Ra T_0$ and
$\area(T_1^i)\le\acrit$ for $i=0,1,2$. Then
\[
\area\{T_0,T_1^0,T_1^1,T_1^2\}>4\acrit.
\]
\end{lemma}

\begin{proof} The proof uses computer interval arithmetic.  See Section~\rif{sec:pent4}.
\end{proof}

The final lemma in this section uses the notion of a \emph{large
angle}, which we now describe.  Let $A$
and $B$ be two pentagons in the plane.  Translate one of the pentagons
$B$ until the pentagons do not overlap and are situated so that some
vertex $\v_B$ of $B$ lies on the relative interior of an edge of $A$.
The four angles at $\v_B$ are $\pi$ (the straight angle), $3\pi/5$
(the interior angle), $\alpha\ge0$, and $\alpha'\ge0$, with
$\alpha+\alpha'=2\pi/5$.  The unordered pair $\alpha,\alpha'$ is
independent of the (constrained) choice of translation.  Clearly, the
same pair of angles is obtained if $A$ is translated rather than $B$.
In the limiting case when the pentagons are in slider contact,
$\{\alpha,\alpha'\}=\{0,2\pi/5\}$.

By ordering $(A,B)$, we obtain an order on
$(\alpha,\alpha')$ as follows.  After the translation described in the
previous paragraph, we may apply an orientation-preserving rigid
motion of the plane so that $A$ is centered at the origin and $B$ is
centered along the positive $x$-axis.  The cone over one angle $\alpha$
meets the upper half-plane in an unbounded region, and the cone over
the other angle $\alpha'$ meets the lower half-plane in an unbounded
region. In this way, we disambiguate $\alpha,\alpha'$.

\begin{definition}[large angle]\libel{def:large}  
  Let $T$ be a $P$-triangle.  Orient the edges of $T$ by the clockwise
  orientation of the triangle.  Let $(A,B)$ be a pair of pentagons in
  the $P$-triangle, ordered by the orientation on the edge $e$ of $T$
  between their centers.  Let $\alpha = \alpha(T,e)$, where $(\alpha,\alpha')$
  is the ordered pair of angles
  attached to the ordered pair $(A,B)$.  We say that the angle is
  {\it large} along $(T,e)$ if $\pi/5<\alpha(T,e) < 2\pi/5$.  (See
  Figure~\rif{fig:large}.)
\end{definition}

\tikzfig{large}{A Delaunay triangle $T$ 
with a large angle $\alpha$ along $(T,e)$.}
{
\begin{scope}[scale=0.8]
\threepent{0.00}{0.00}{50.58}{1.10}{1.49}{181.07}{1.73}{0.00}{217.07};
\draw (0,0) node[anchor=north] {$A$};
\draw (0,0) node[anchor=south] {$e\ $};
\draw (1.1,1.49) node[anchor=south] {$B$};
\draw[black] (50.58:1) -- ++ ($2*(50.58+72:1) - 2*(50.58:1)$);
\draw[black] (50.58:1) -- ++ ($2*(1.1,1.49)+2*(181.07:1) - 2*(50.58:1)$);
\draw (-0.8,1.4) node {$\alpha$};
\end{scope}
}

Let $T$ be a $P$-triangle and let $T'$ be the adjacent $P$-triangle
along edge $e$.  We have the invariant
$\alpha+\alpha'=\alpha(T,e)+\alpha(T',e)=2\pi/5$.  Hence if the angle
is large along $(T,e)$ then it is not large along $(T',e)$.

% COMPUTATION FHBGHHY | calc:large | Non-large-angle two-triangle estimate
\begin{lemma}\libel{FHBGHHY}\libel{calc:large}
  Let nonobtuse $P$-triangles $\{T_0,T_1\}$ be given (not necessarily a
  pseudo-dimer) such that $\area(T_1)\le \acrit$ and $T_1\Ra T_0$.
  Assume that there is a nonshared edge $e$ of $T_0$ of length greater
  than $1.8$ and such that the angle is not large along $(T_0,e)$.  Then
  $\area\{T_0,T_1\} > 2\acrit + \epsM$.
\end{lemma}
In particular, $\{T_0,T_1\}$ is not a pseudo-dimer.
\begin{proof} The proof uses computer interval arithmetic. See Section~\rif{sec:large}.
\end{proof}

\begin{corollary}\libel{lemma:large}
  If $(T_1,T_0)\in \PD$ and $e$ is the longest (that is, egressive)
  edge of $T_0$.  Then the angle is large along $(T_0,e)$.
\end{corollary}

\begin{proof} %
  The egressive edge has length greater than $1.8$ by Lemma~\rif{calc:pseudo2}.
  The edge is nonshared by Lemma~\rif{calc:pseudo1}.
  If the angle is not large, then we can apply Lemma~\rif{calc:large} to find that
  the area of the pseudo-dimer is greater than $2\acrit + \epsM$,
  which contradicts one of the defining properties of a pseudo-dimer.
\end{proof}

\section{Dimer Pairs}\libel{sec:dimer-pair}

The purpose of this section is to give a proof of the following
theorem.  This theorem is the principal optimization problem of this
article in the sense that all other optimizations deal with
configurations that are far from optimal.

\begin{theorem}\libel{thm:dimer-area}
  Let $(T_1,T_0)$ be a dimer pair.  (In particular, we assume that
  $T_1$ is subcritical, that $\area\{T_0,T_1\}\le 2\acrit$ and that
  $T_0$ and $T_1$ share a common longest edge.)  Then $(T_1,T_0)$ is
  the ice-ray dimer of area exactly $2\acrit$.
\end{theorem}

The proof will fill the entire section.  The strategy of the proof is
to give a sequence of area decreasing deformations to $(T_1,T_0)$,
until the ice-ray dimer is reached.

\noindent
{\bf Context and Notation.} We fix context and
notation that will be used throughout this section. Let
$(T_1,T_0)$ be a dimer pair.  The $P$-triangle $T_1$ has a pentagon
centered at each vertex.  We label the pentagons of $T_1$ as $A$, $B$,
$C$, with $A$ and $C$ shared with $T_0$.  We call $B$ the outer
pentagon of $T_1$.  Similarly, we label the pentagons of $T_0$ as $A$,
$C$, $D$, with outer pentagon $D$ of $T_0$.

\begin{lemma}  
  Let $(T_1,T_0)$ be a dimer pair, then every edge of $T_1$ and $T_0$
  has length less than $\kappa\sqrt8$.  In particular $T_1$ and $T_0$
  are both acute (and not merely nonobtuse).
\end{lemma}

\begin{proof}  
  The shared edge between $T_0$ and $T_1$ is the common longest edge
  of the two triangles.  It is enough to show that this edge has
  length less than $\kappa\sqrt8$.  This is an edge of a subcritical
  triangle $T_1$. The result follows from Lemma~\rif{lemma:right}.
\end{proof}

As long as a deformation of a dimer pair preserves the common-longest-edge
constraint and the total area bound $2\acrit$, it cannot reach a
first right triangle.  At such a first event both triangles are
nonobtuse, and one of them has area at most $\acrit$.  By
Lemma~\rif{lemma:right}, the common longest edge then has length
less than $\kappa\sqrt8$.  But it would also be the hypotenuse of
a right triangle with legs at least $2\kappa$, a contradiction.

The deformation of a general dimer pair to the ice-ray dimer takes
place in several stages.  We give a summary of the stages here, before
going into details.  Here is the proof sketch:

\begin{enumerate}
\item We deform the dimer pair so that each of $T_0$ and $T_1$ is
  $O2C$ or long isosceles.
\item We deform so that each of $T_0$ and $T_1$ is $O2C$.
\item We show that both triangles are triple contact.
\item Working with triple contact triangles, we compute that the
  condition $\area\{T_0,T_1\}\le 2\acrit$ implies that $(T_1,T_0)$
  lies in a small explicit neighborhood of the ice-ray dimer.
\item We construct a curve $\Gamma$ in the configuration space of
  geometric dimer configurations, with parameter $t$ such that $t=0$ defines the ice-ray
  dimer.
\item Working with triple contact triangles in a small explicit
  neighborhood of the ice-ray dimer, and for some small explicit
  constant $M$, each dimer pair can be connected by a path (in the
  ambient geometric dimer configuration space) to a point on the curve $\Gamma$ with
  parameter $|t|<M$.  A computation shows that the area of the dimer
  decreases along the path to $\Gamma$.  Thus, every area-minimizing
  dimer pair lies on the curve $\Gamma$.
\item The unique global minimum of the area function along $\Gamma$
  occurs at $t=0$; that is, the ice-ray dimer is the unique global
  minimizer along $\Gamma$ for $|t|<M$.
\end{enumerate}

\subsection{Reduction to $O2C$ or long isosceles (stage 1)}\libel{sub:long}

In this section, we show that each of $T=T_0$ and $T=T_1$ can be
deformed in an area decreasing way until $T$ is either $O2C$ (that is,
the outer pentagon has double contact) or long isosceles (that is, the
two longest edges of the triangle have the same length).

Fix $A$ and $C$ along the shared edge.  For $T=T_1$ the moving
outer pentagon is $B$; for $T=T_0$ it is $D$.  Start at a
configuration whose shared edge is uniquely longest, and follow an
area-nonincreasing motion only up to the first of two events: the
outer pentagon reaches double contact, or a nonshared edge becomes
as long as the shared edge.  The first event is $O2C$; the second
is the long-isosceles boundary considered in stage 2.  A configuration
already on that boundary is passed directly to stage 2.

This is easy to carry out. By Lemma~\rif{lemma:primary}, we may
choose a primary-contact configuration of no greater area, without
increasing the incident edges. We then move the outer
pentagon by translation (meaning no rotation) that preserves contact
with either $A$ or $C$ and that is non-increasing in triangle area
up to the first $O2C$ or longest-edge-tie event.  The motions are
applied separately to $T_1$ and $T_0$, with the same two shared
pentagons fixed.

% COMPUTATION GAIMXAD | sec:O2C | Long-isosceles terminal exclusions (stage 2)
\subsection{Reduction to $O2C$ (computer interval-arithmetic assisted) (stage 2)}\libel{sec:O2C}

In this subsection we show that each of $T=T_0$ and $T=T_1$ can be
deformed in an area decreasing way so that it is $O2C$.
In other words, we eliminate the long-isosceles possibility allowed in stage 1.

We begin with the case $T=T_0$.  If $T_0$ is a long isosceles
triangle, then $(T_1,T_0)$ is a boundary case that is covered by the
final sentence of Lemmas~\rif{calc:pseudo2} and ~\rif{calc:pseudo3}.
These lemmas give edge-length conditions that imply that $T_0$ is not
long isosceles.  For the rest of the proof, we disregard $T_0$ and
deal only with $T_1$.

We now consider the case $T=T_1$.  In view of the reductions of the
previous subsection, we may assume that $T_1$ has primary contact and
that $T_1$ is long isosceles.  We may further assume that translation
of the outer pentagon $B$ while maintaining contact (with $A$ or $C$)
in an area decreasing direction would violate the constraint that the
longest edge of $T_1$ is the shared edge.  (In other words, the
translation that decreases area would increase the edge length of the
long nonshared edge.)  We complete this stage of the proof by proving
non-subcriticality of $T_1$, which is contrary to the defining
conditions of a dimer pair.  The proof is computer interval-arithmetic
assisted.

In the deformations that follow, a stopping condition is that
primary contact is $O2C$.
We deform $T_1$ in an area decreasing
way into a configuration that can be easily computed.  Without loss of
generality, we assume that the outer pentagon $B$ is in contact with
pentagon $A$.  

We consider two cases, depending on whether the primary contact of $B$
has slider contact or midpointer contact with $A$.  We need two
non-anomaly lemmas, one for slider contact and one for midpointer
contact.  In both cases, the lemmas imply that the edge of contact
between $A$ and $B$ (whether slider or midpointer) is one of the long
edges of the isosceles triangle.

\begin{lemma}[slider-non-anomaly]  
  Let $T$ be a subcritical nonobtuse $P$-triangle with pentagons $A$,
  $B$, and $C$.  Assume that $B$ has slider contact with $A$.  Then
  the translation of $B$ along the slider contact in the direction to
  decrease the area also decreases the length $\dx{BC}$.
\end{lemma}

\begin{proof} 
  Assume to the contrary that the translation is increasing in
  $\dx{BC}$.  Choose coordinates so that the $x$-axis passes through
  $\c_A$ and $\c_C$, with the origin at $\c_A$, with $\c_C$ in the
  right half-plane, and with $\c_B$ in the positive half-plane
  (Figure~\rif{fig:slider}).  Our contrary assumption means the 
  line $\lambda$ through the edge of contact between $A$ and $B$ has
  positive slope.  Slider contact implies that the center $\c_B$ lies
  on the line $\lambda'$ parallel to $\lambda$ at distance $2\kappa$
  from the origin.  Every point on $\lambda'$ either has
  $y$-coordinate at least $2\kappa$ or negative $x$-coordinate.  If
  the $y$-coordinate is at least $2\kappa$ the triangle is not
  subcritical.  (The area is at least $2\kappa^2>\acrit$.) If the
  $x$-coordinate of $\c_B$ is negative, then the triangle $T$ is
  obtuse.
\end{proof}

\tikzfig{slider}{The center $\c_B$ of the nonobtuse triangle lands in
  the first quadrant and gives a triangle $(\c_A,\c_B,\c_C)$ of height
  and base both at least $2\kappa$.}
{
[scale=1.0]
\draw (0,0) -- (2,0);
\draw (0,0) -- (0,2);
\coordinate (V) at ($(0,0) + (-90+2*72+10:1)   $);
\draw (V) -- ++ (10:1.5) node[anchor=west] {$\lambda$};
\draw (V) -- ++ (180+10:1.5);
\coordinate (B) at ($(V) + (10:0.5) + (90+10- 2*72 - 180:1)$);
\draw (B) -- ++ (10:1.5) node[anchor=west] {$\lambda'$};
\draw (B) -- ++ (180+10:1.5);
\penp{B}{90+10};
\pen{0}{0}{-90+10};
\smalldot{B};
\smalldot{0,0};
\smalldot{2,0};
\draw (B) node [anchor=south west] {$\c_B$};
\draw (0,0) node [anchor=north] {$\c_A$};
\draw (2,0) node [anchor=north] {$\c_C$};
}

\begin{lemma}[midpointer-non-anomaly]  
  Let $T$ be a subcritical nonobtuse $P$-triangle with pentagons $A$,
  $B$, and $C$.  Assume that $B$ is not in contact with $C$.
  Assume that $A$ has midpointer contact with $B$, with
  $A$ pointing to $B$ at $\v_{AB}$.  Then there is a motion of $B$
  that is nonincreasing in the area and that preserves contact between $A$ and $B$,
  while also decreasing the distance $\dx{BC}$.
\end{lemma}

\begin{proof}
  Assume to the contrary that the translation is increasing in
  $\dx{BC}$.  Choose coordinates so that the $x$-axis passes through
  $\c_A$ and $\c_C$, with the origin at $\c_A$, with $\c_C$ in the
  right half-plane, and with $\c_B$ in the positive half-plane.
  Let $\c_B=(x_B,y_B)$ and $\v_{AB}=(x_{AB},y_{AB})$.
  The nonobtuse condition gives $x_B\ge0$.

  If the line $\lambda$ through the edge of $B$ containing $\v_{AB}$
  has nonpositive slope, then the translation of $B$ to the
  right along a line parallel to $\lambda$
  fulfills the conclusion of the lemma.  Thus, we may assume without
  loss of generality that the slope of $\lambda$ is positive.

  If $x_B\ge x_{AB}$, then the rotation of $B$ about the pointer $\v_{AB}$
  to the right fulfills the conclusion of the lemma.
  Thus, we may assume without loss of generality that $x_B<x_{AB}$.

  We claim that $y_B\ge 2\kappa$.  For the proof of this claim we may disregard $\c_C$.
  Rotate $A,B$ rigidly about $\c_A$ clockwise until the slope of $\lambda$ is $0$.
  This motion also rotates $\c_B$ clockwise and is decreasing in $y_B$.
  Thus, we may assume without loss of generality that $\lambda$ is horizontal.
  $\lambda$ is a separating line between pentagons $A$ and $B$.  The distance
  from $\c_B$ to $\lambda$ is $\kappa$. Since $\lambda$ is a support line of $A$,
  its distance from the $x$-axis is at least the inradius $\kappa$.  Thus, $y_B\ge 2\kappa$.

  As in the proof of the previous lemma, the lower bound on $y_B$ contradicts subcriticality.

\end{proof}

As a corollary of these two lemmas, if $d_{AB}$ is not a longest edge,
then there exists a deformation decreasing area and $d_{BC}$.  This
allows us to reduce the long isosceles triangle to a triangle such
that the long edges are $(\c_A,\c_B)$ and $(\c_A,\c_C)$.

Because we are now disregarding $T_0$, we may deform the
triangle $T_1$ by translating $C$, preserving the long isosceles constraint
and decreasing area, until $C$ has primary contact. Together with the contact between
$A$ and $B$, this gives double contact.

 We choose coordinates and
compute\calcref{iso\_2C,iso\_2C'} with interval
arithmetic to show that no $P$-triangle exists with the
following properties:
\begin{enumerate}
\item The triangle has double contact.
\item The pentagons $A$ and $B$ have contact,
  and the contact type between $A$ and $B$ is either slider contact or
  midpointer contact.
\item The triangle is long isosceles, with long edges $(\c_A,\c_B)$ and $(\c_A,\c_C)$.
\item The triangle is subcritical and acute.
\end{enumerate}
This
completes this reduction.

\subsection{Reduction to triple contact (stage 3)}\libel{sec:squeeze}

Before turning to stage 3, we describe a deformation that we
call \emph{squeezing}.  It involves two adjacent $P$-triangles
$T_1=(A,B,C)$ and $T_0=(A,D,C)$, as in the notation at the
start of this section.  The same construction applies to a single
triangle.  All pentagons move by translation, so their orientations
are fixed.  We use variants of this motion in
Sections~\rif{sec:pseudo-area}, \rif{sec:pseudo-area3}, and
\rif{sec:pent3}.

Let $T_0$ be a $P$-triangle with pentagons $A$, $D$, $C$ and $T_1$ a
$P$-triangle with pentagon $A,B,C$.  We assume that $\c_A$ and $\c_C$
lie on the $x$-axis with $\c_A$ to the left of $\c_C$, with $\c_B$ in
the upper half-plane, and $\c_D$ in the lower half-plane.  We assume
that the triangles $T_0$ and $T_1$ are initially acute.  The squeezing
motion starts with $A,C$ not in contact, and moves $\c_A$ and
$\c_C$ continuously towards each other.  A stopping condition is
contact between $A$ and $C$.

If $B$ is not in contact with another pentagon, the motion holds $B$ fixed.
If $B$ is in contact with $A$ or $C$ but not both, the motion
moves $\c_B$ along a horizontal line to maintain contact.
If $B$ is $O2C$, the motion translates $B$ to maintain its contacts
with $A$ and $C$. Its vertical direction depends on the contact branch.
The center $\c_D$ is moved in a corresponding way according to its contacts.

Choose a translating coordinate system in which
$\c_A=(0,0)$, $\c_C=(\ell,0)$, with $\ell>0$ decreasing,
and with $\c_B$ above the axis and $\c_D$ below it.  We
initially assume the two triangles are acute and $A,C$ are not
in contact.  On a branch maintaining two vertex--edge contacts of
$B$ with $A,C$, let $\mathbf{n}_{AB},\mathbf{n}_{CB}$ be normals to the two
contact edges.  These normals are constant on the branch.  Taking
$\ell(t)=\ell(0)-t$, the velocity $\u_B$ is determined by
\[
\mathbf{n}_{AB}\cdot \u_B=0,\qquad
\mathbf{n}_{CB}\cdot\bigl(\u_B+(1,0)\bigr)=0.
\]
The system has a unique solution when the two normals are linearly
independent.  If the two contact lines are parallel but not horizontal, we may
first slide $B$ along them, with $A,C$ fixed, in the direction to
decrease its height. If both contact lines are horizontal, we may
instead decrease $\ell$ while keeping the height of $B$ fixed. In
either case the area decreases without solving a full-rank system.
The full-rank linear system is just the derivative of the two affine contact
equations, so its integral maintains those contacts until their
vertex or edge labels change.  At a transition point from one branch to
another, the motion of $\c_B$ is continuous but not differentiable.
Transition between branches might be needed to obtain $A$--$C$ contact.
The same equations define $\u_D$
with the corresponding contact normals.

This completes our description of the squeezing motion, and we return
to the reduction to triple contact.  In this stage, we initially
assume that both triangles $T_0$ and $T_1$ are $O2C$.  We deform so
that both triangles are triple contact.  We briefly describe the
argument.  Our deformations will preserve the $O2C$ contacts.  By an
argument made in the second paragraph of Section~\rif{sec:O2C}, we may
assume without loss of generality that the triangle $T_0$ is not long
isosceles.  By Lemma~\rif{calc:dimer-isosc}, we have that $T_1$ is not
long isosceles.  Assuming that neither triangle is long isosceles, we
show that both triangles can be brought into triple contact.  Because
both triangles are already $O2C$, this amounts to decreasing the edge
between $\c_A$ and $\c_C$ until the pentagons $A$ and $C$ come into
contact.

% COMPUTATION LZMJIZU | lemma:squeeze-area | Squeezing area inequality (stage 3)
\begin{lemma}\libel{lemma:squeeze-area}
  Let $T$ be a nonobtuse $P$-triangle with pentagons $A$, $B$, and
  $C$, where $B$ has double contact.  Assume that the area of $T$ is  at most $\acrit+\epsN$.
  Then the normals $\mathbf n_{AB}$ and $\mathbf n_{CB}$ are not horizontal or vertical.
  Then the squeezing
  deformation decreases area.
\end{lemma}

\begin{proof}
  The proof is computer-assisted.
   Recall that $\dx{XY} = \dd{X}{Y}$, and let
  $\arc_X$ be the angle of the triangle at $\c_X$.  The area of $T$ is
   $\area(T) = d_{AC} d_{AB} \sin(\arc_A)/2$.

   Section~\rif{sec:two} gives a parametrization of two pentagons in contact.
   If $B$ is fixed, the contacts with $A$ (and $C$) can be rotated by multiples
   of $2\pi/5$ around $B$, while keeping the same parameters.
   The angle of $T$ at $\c_B$ lies in the interval
   \([\arc(2,2,2\kappa),\pi/2]\).  Since the width of this interval is less
   than $2\pi/5$, for each choice of contact parameters and triangle orientation, there is at most one choice
   of rotation parameter $2\pi m/5$ modulo $2\pi$ that gives an angle of $T$ at $\c_B$
   in the given interval.
   Orient the triangle $T$ with $\c_A$ at the origin, $\c_C=(\ell,0)$, with $\ell>0$ and
   $\c_B$ above the axis.
% COMPUTATION GJULAOH | lemma:squeeze-area | four contact-sign cases
   We check by a computer-assisted interval arithmetic calculation that under the conditions of the
   lemma ($O2C$, nonobtuseness, area at most $\acrit+\epsN$), the contact edge
   between $A$ and $B$ is not vertical and has negative slope.
   Moreover, $\c_B$ lies above the line through the contact edge and $\c_A$ below.
   By symmetry, the contact edge between $C$ and $B$ is not vertical and has a positive slope.
   The line through the contact edge separates $\c_C$ from $\c_B$, with $\c_B$ above the line.
   
  On any differentiable contact branch, put $h=d_{AB}\sin(\arc_A)$
  and $\ell=d_{AC}$.  The signed area derivative is
  $(\ell'h+\ell h')/2$.  Thus $\ell'<0$ and $h'\le0$
  give strict decrease without dividing by a slope.

  By the interval-calculation of contact edge slopes above, we may
  choose a horizontally translating
  frame in which $B$ moves directly upward by $\Delta y>0$,
  $A$ moves right by $\Delta x_A>0$, and $C$ moves left by
  $\Delta x_C>0$.  The transformed area is
  $(d_{AC}-\Delta x_A - \Delta x_C)
  (d_{AB}\sin(\arc_A)+\Delta y)/2$.
  Let $\sigma_A=dy/dx_A$ and $\sigma_C=dy/dx_C$ be the
  positive finite slopes on this branch.  Differentiating with
  respect to $y$ gives
\[
\frac{d\area(T)}{dy}
=\frac12\left[d_{AC}-d_{AB}\sin(\arc_A)
 \left(\frac1{\sigma_A}+\frac1{\sigma_C}\right)\right].
\]
  Thus squeezing decreases the area on this branch exactly when
\[
d_{AC} < \left(\frac{1}{\sigma_A}
+\frac{1}{\sigma_C}\right)d_{AB}\sin(\arc_A).
\]
The computer
interval-arithmetic calculation\calcref{calculation squeeze\_calc}
tests the disjunction
\[
\sigma_A<h/\ell,\qquad \sigma_C<h/\ell,\qquad\hbox{or}\qquad
\left(\frac1{\sigma_A}+\frac1{\sigma_C}\right)h>\ell.
\]
When both slopes are positive and finite, either of the first two
inequalities implies the third. Thus the tested disjunction gives
strict area decrease.

We defined $\sigma_A$ and $\sigma_C$ as derivatives, but in fact no
differentiation is required.  For example, consider $\sigma_A$.  The
pentagon $B$ has contact with $A$.  Thus, $B$ points into $A$ or $A$
points into $B$.  If $A$ points into $B$, then the point of contact
lies along an edge $e$ of $B$.  The squeeze transformation translates
$A$ and $B$ maintaining the contact.  Viewed from a coordinate system
that fixes $B$, the squeezing lemma translates $A$ parallel to the
line through $e$.  On the positive-slope branch, $\sigma_A$ is the absolute value of
the slope of the line through $e$.  Elementary coordinate calculations
described in the coordinate section of this article give an explicit
formula for the slope of this line.  There are two cases, depending on
which pentagon points to the other.  There are no difficulties in
carrying out the computer interval-arithmetic calculations with these explicit formulas.
At a vertex--vertex contact the formula must be applied to a
specified one-sided vertex--edge branch.
\end{proof}

\begin{lemma}\libel{lemma:dimer-acrit-epsN}
  Let $(T_1,T_0)$ be a dimer pair.
  Then $\area(T_0) < \acrit+\epsN$.
\end{lemma}

\begin{proof}
The triangle $T_1$ is subcritical, and
$\area\{T_0,T_1\}\le2\acrit$.  By Lemma~\rif{lemma:a0},
$\area(T_1)>\ao$.  Hence
\begin{equation}\libel{eq:dimer-acrit-epsN}
\area(T_0) = \area\{T_0,T_1\} - \area(T_1) 
< 2\acrit - \ao = \acrit+\epsN.
\end{equation}
\end{proof}

Thus the area hypothesis of Lemma~\rif{lemma:squeeze-area} holds
for both $T_0$ and $T_1$.  On an admissible squeezing branch,
individual area decrease preserves these bounds.  We stop at the
first contact change or longest-edge tie, using the exclusions of
stage 2 for the latter boundary.  Continuation to $A$--$C$
contact, with the $O2C$ contacts and nonobtuseness retained,
completes the reduction to $3C$ contact.

\subsection{Reduction to a small neighborhood of the ice-ray
% COMPUTATION DSKKWND | sec:nbd | Exclusion outside the ice-ray neighborhood (stage 4)
  dimer (stage 4)}\libel{sec:nbd}

From this stage forward, $T_0$ and $T_1$ are both triangles with
triple contact.  We show that the condition $\area\{T_0,T_1\}\le
2\acrit$ implies that $(T_1,T_0)$ lies in a small explicit
neighborhood of the ice-ray dimer.

The shared pentagons $A$ and $C$ are in contact.  By symmetry, we may
assume that $A$ points into $C$.  We have classified all $P$-triangles
with triple contact in Section~\rif{sec:contact}.  At this stage, we
rely heavily on this classification.  The triple contact type $\Delta$
has area at least $1.5 > \acrit+\epsN$ by Lemma~\rif{lemma:delta}.
This is too large an area to be part of a dimer pair.  There are eight
combinatorial ways that the pair $(A,C)$ with $A$ pointing to $C$ can
be extended to a triple contact triangle: a pinwheel, a pin-$T$
junction, an $LJ$-junction (3 ways), and a $TJ$-junction (3 ways).
There are three ways of extending the edge along $(A,C)$ to an
$LJ$-junction depending on which of the three triangle edges of the
$LJ$-junction is placed along $(A,C)$.  A similar remark applies to
$TJ$-junctions.  A pinwheel has cyclic symmetry, so it gives rise to a
single case.

We need an argument to show that only a single combinatorial type of
pin-$T$ triangle needs to be considered.  In this paragraph only, we
shift notation and refer to labels on pentagons in the pin-$T$
configuration shown in Figure~\rif{fig:cord-pint}.  We claim that if
the triangle has area less than $\acrit+\epsN$, then the edge
length $\dx{BC}$ in that figure is the unique longest edge of the
triangle.  This claim is established by computer interval-arithmetic calculation.  (See
Section~\rif{sec:one-pintx}.)  This means that the shared edge of the
pin-$T$ triangle is determined.

Because $T_0$ and $T_1$ both have triple contact, the dimer pair
$(T_1,T_0)$ lies in a four-dimensional configuration space.  The
strategy of the proof is to check by interval arithmetic that there does not
exist a dimer pair outside a small explicit neighborhood of the
ice-ray dimer.  In other words, the area constraints $\area(T_1)\le
\acrit$ and $\area\{T_0,T_1\}\le2\acrit$ are impossible to satisfy
when the longest edge on both triangles is the shared edge.  We run
over $64 = 8\times 8$ cases depending on the combinatorial types of
$T_0$ and $T_1$.  In each case, $T_1$ runs over a three-dimensional
configuration space.  Most of the $64$ cases do not contain the
ice-ray dimer.  In these cases, it is not necessary to specify a small
explicit neighborhood.  In the cases that do contain the ice-ray
dimer, the neighborhood is described in local coordinates.

The choice of coordinates depends on the case.  We refer to
Section~\rif{sec:coords} and the associated interval-arithmetic computer code
for more details about the coordinates and calculations.
For each of the contact types meeting the ice-ray configuration,
the shared-edge charts express these coordinates in terms of the
raw parameters.  The ice-ray dimer is the origin in these centered
coordinates.  Write
\[
U=\{(x_1,x_2,x_3,x_4): |x_i|<0.01\text{ for }1\le i\le4\}
\]
within these charts.  The exclusion outside the union of these
case-dependent neighborhoods is
what is needed to complete stage 4.

% COMPUTATION BKUSHRZ | sec:gamma | Convexity along Gamma (stage 7)
\subsection{Defining a curve (stage 5)}\libel{sec:gamma}

At this stage, we define a curve $\Gamma$ in the configuration space
of two adjacent triple-contact triangles sharing two pentagons.
The ice-ray dimer is defined by
parameter value $t=0$.  We represent the ice-ray dimer by a pair of
triple-contact triangles with shared pentagons $A$ and $C$, where $A$
points to $C$.  The curve $\Gamma$ is described by the translational shear motion that
slides along all four edges of contact, illustrated in Figures
\rif{fig:Gamma} and \rif{fig:Gamma'}.
The parameter $t$ is the signed distance between the pointer vertex
$A\to C$ and the midpoint of the receptor edge.
See Figure~\rif{fig:Gamma'}.

\tikzfig{Gamma'}{Three configurations along the curve $\Gamma$.}
{
\begin{scope}[scale=0.5,xshift=-6cm]
\threepentnoD{0.00}{0.00}{79.87}{0.93}{1.36}{259.87}{1.83}{0.00}{223.87};
\threepentnoD{0.00}{0.00}{-64.13}{1.03}{-1.50}{-244.13}
{1.83}{0.00}{-208.13};
\draw (0,0) node {$A$};
\draw (1.83,0) node {$C$};
\draw (0.5,-3.5) node {$t= 0.25$};
\smalldot{79.87-72:1};
\smalldot{$(1.83,0)+(223.87-3*72+180:0.809)$};
\end{scope}
\begin{scope}[scale=0.5]
\threepentnoD{0.00}{0.00}{72.00}{0.96}{1.43}{252.00}{1.81}{0.00}{216.00};
\threepentnoD{0.00}{0.00}{-72.00}{0.96}{-1.43}{-252.00}{1.81}{0.00}{-216.00};
\draw (0,0) node {$A$};
\draw (1.81,0) node {$C$};
\smalldot{1,0};
\draw (0.5,-3.5) node {$t=0$};
\end{scope}
\begin{scope}[scale=0.5,xshift=6cm]
\threepentnoD{0.00}{0.00}{-79.87}{0.93}{-1.36}{-259.87}{1.83}{0.00}{-223.87};
\threepentnoD{0.00}{0.00}{64.13}{1.03}{1.50}{244.13}{1.83}{0.00}{208.13};
\draw (0,0) node {$A$};
\draw (1.83,0) node {$C$};
\draw (0.5,-3.5) node {$t= -0.25$};
\smalldot{-79.87+72:1};
\smalldot{$(1.82,0)+(-223.87+3*72+180:0.809)$};
\end{scope}
}

In terms of the coordinates $(x_\alpha,\alpha)$ of
Section~\rif{sec:coords} for two pentagons in contact, the relative
position of $A$ and $C$ is described by the curve $\alpha=\pi/5$ and
$x_\alpha = \sigma+t$.  We call a $\Gamma$-configuration a
configuration on the curve $\Gamma$.  We define a $\Gamma$-triangle to
be a $P$-triangle that occurs as one of the two triangles in a
$\Gamma$-configuration.

% COMPUTATION LPEYXGH | sec:curve | Descent to the shear curve (stage 6)
\subsection{Reduction to the curve (stage 6)}\libel{sec:curve}

At this stage, we show that we can reduce to points on the curve in
the following sense.  Working with triple contact triangles in a small
explicit neighborhood $U$ of the ice-ray dimer ($|x_i|<M$, where
$M=0.01$, for $i=1,2,3,4$ in appropriate coordinates), each dimer pair
$(T_1,T_0)$ can be connected by a path (in the ambient configuration
space) to a $\Gamma$-dimer with parameter $|t|<0.01$.  The area
function is decreasing along this path.  We have two tasks.  First, we
construct the path $P$, and then we show that the area function
decreases along the path.  We will use $s$ as a local parameter on the
path ($s\mapsto P(s)$), which we need to keep separate from the
parameter $t$ for $\Gamma$.  We construct a path such that $s=0$
determines the initial dimer pair $(T_1,T_0)$ and such that the path
$P$ is defined for $s\in [0,s_0]$, for some $s_0\ge0$; the path is
constant if the initial point is already on $\Gamma$.

Section~\rif{sec:shared} accomplishes these tasks.  The path $P$ is
constructed and it terminates on $\Gamma$.  If the initial point of
the path lies in the neighborhood $U$, then the path stays in $U$
and terminates at a point on $\Gamma$ with parameter $|t|<M$.

Finally, we need to check that the area function is decreasing. This,
we prove by a computer interval-arithmetic calculation of the derivative of the area
function along the path at $s=0$.  For this, we use automatic
differentiation algorithms, as described in
Section~\rif{sec:auto-diff}.  This completes the reduction to points
on the curve $\Gamma(t)$.

\subsection{Global minimization along the curve (stage 7)}

The final stage of the proof of Theorem~\rif{thm:dimer-area} is the
minimization of the area of a $\Gamma$-dimer, for parameters
$|t|<0.01$.  We have now reduced to a straightforward minimization problem in a
single variable. 

A determinant is used to compute $\area\{T_0,T_1\}$ of a $\Gamma$-dimer.
Fix $A$ and keep the
orientations of all four pentagons fixed during the translational
shear. Let $\mathbf v_1$ be the unit vector from $\c_A$ towards
$\c_C(0)$, and let $\mathbf u_1$ be the perpendicular unit vector in
the direction in which $C$ moves as $t$ increases. Interchanging
$B$ and $D$ if necessary, assume that $B$ lies on the
positive-$\mathbf u_1$ side at $t=0$. We take determinants in the
ordered orthonormal frame $(\mathbf v_1,\mathbf u_1)$.

By the definition of $t$,
\[
\c_C(t)-\c_C(0)=t\mathbf u_1.
\]
The displacement of $B$ is parallel to its edge of contact with
$A$, while its displacement relative to $C$ is parallel to its
edge of contact with $C$. These displacements form a triangle
with angles $\pi/5,2\pi/5,2\pi/5$, in which the displacements of
$B$ and $C$ are the two equal sides. Thus $B$ slides through
signed distance $t$. The same argument applies to $D$. The
slider directions make angles $\pi/5$ with $\mathbf u_1$, on
opposite sides, giving
\[
\begin{aligned}
\c_B(t)-\c_B(0)&=t(-\sigma\mathbf v_1+\kappa\mathbf u_1),\\
\c_D(t)-\c_D(0)&=t(\sigma\mathbf v_1+\kappa\mathbf u_1).
\end{aligned}
\]

Let
\[
\mathbf u(t)=\c_C(t)-\c_A,\qquad
\mathbf v(t)=\c_B(t)-\c_D(t)
\]
be the diagonal vectors of the quadrilateral of centers. Then
\[
\mathbf u(t)=\mathbf u(0)+t\mathbf u_1,\qquad
\mathbf v(t)=\mathbf v(0)-2\sigma t\mathbf v_1.
\]
At the ice-ray configuration, $\mathbf u(0)$ is parallel to
$\mathbf v_1$ and $\mathbf v(0)$ is parallel to $\mathbf u_1$.
Consequently the mixed terms vanish, and
\[
\det(\mathbf u(t),\mathbf v(t))
=
\det(\mathbf u(0),\mathbf v(0))+2\sigma t^2.
\]

The two triangles lie on opposite sides of their shared edge,
so their total area is half the determinant. At
$t=0$, their total area is $2\acrit$. Hence
\[
\area\{T_0(t),T_1(t)\}
=
\frac12\det(\mathbf u(t),\mathbf v(t))
=
2\acrit+\sigma t^2.
\]
Since $\sigma>0$, the unique minimum along the portion of
$\Gamma$ under consideration occurs at $t=0$, the ice-ray
configuration.

An alternative derivation of area is given by Heron's formula. We give a brief sketch.
Write \(x=\bar x_\beta\). Along the given curve, the three squared side lengths of one triangle simplify to
\[
4\kappa^2+x^2,\qquad
4\kappa^2+\frac{x^2}{4\kappa^2},\qquad
(1+\kappa)^2+(\sigma-x)^2.
\]
Substitution into Heron's triangle-area formula,
followed by adding the area of the second triangle with parameter \(2\sigma-x\), gives exactly
\[
\area\{T_0,T_1\}
      =2a_{\rm crit}+\sigma(x-\sigma)^2
      =2a_{\rm crit}+\sigma t^2.
\]

This completes the proof of Theorem~\rif{thm:dimer-area}.

\section{Pseudo-Dimers}\libel{sec:pseudo}

We now determine the corrections on a pseudo-dimer $(T_1,T_0)$.
Lemma~\rif{lemma:pd-1} shows that the source $T_1$ has no area correction,
and Lemma~\rif{lemma:pd-b} shows that $b(T_0)>\acrit$.  These are
used in Section~\rif{sec:dimer-main} to exclude longer nonobtuse
attachment chains and in Section~\rif{sec:nonobtuse} to count the
corrections on small clusters.

For a nonobtuse triangle, the area is a monotonic function of its edge
lengths.  This makes it easy to give lower bounds on triangle areas.
Here are some simple area calculations that will be used.
\begin{align*}
\area(1.8,1.8,1.8) &> \acrit + 0.112\\
\area(1.8,1.8,1.72) &> \acrit + 0.069\\
\area(1.8,1.72,1.72) &> \acrit + 3\epsM\\
\area(1.8,1.8,2\kappa) &> \acrit + \epsM\\
\area(\kappa\sqrt8,2\kappa,1.72) &> \acrit +\epsN.
\end{align*}

\begin{lemma}\libel{lemma:pd-m2}  
If $(T_1,T_0)\in\PD$, then $m_-(T_0)=0$.
\end{lemma}

\begin{proof} 
  Assume for a contradiction that $m_-(T_0)>0$.  By the definition
  of $\M$, there exists $(S_1,S_0)\in\PD$ with
  $(S_0,T_0)\in\M$.  Let $e$ be the edge shared by $S_0$ and
  $T_0$.  It is the egressive edge of $S_0$, and has length greater
  than $1.8$ by Lemma~\rif{calc:pseudo2}.  The two non-egressive
  edges of $T_0$ have length less than $1.8$ by
  Lemma~\rif{calc:pseudo3}.  Thus, $e$ is also the egressive edge
  of $T_0$.  Corollary~\rif{lemma:large} makes both
  $\alpha(S_0,e)$ and $\alpha(T_0,e)$ large.  This contradicts
  $\alpha(S_0,e)+\alpha(T_0,e)=2\pi/5$.
\end{proof}

\begin{corollary}\libel{lemma:m1m2}  
  For every triangle $T$, either $m_+(T)=0$ or $m_-(T)=0$.
\end{corollary}

\begin{proof}  
  If $m_+(T)>0$, then by the definition of $\M$, there exists some $(T_1,T)\in\PD$.  
  Lemma~\ref{lemma:pd-m2} gives $m_-(T)=0$.
\end{proof}

\begin{lemma}[pseudo-dimer disjointness]  
  Assume $(T_1,T_0)\in \PD$ and $(T'_1,T'_0)\in\PD$ and
  $\{T_0,T_1\}\cap\{T'_0,T'_1\} \ne \emptyset$.  Then $T_0 = T'_0$.
\end{lemma}

\begin{proof}  
  If we dismiss the other three cases $T_1=T'_1$ and $T_1=T'_0$ and
  $T_0=T'_1$ of a nonempty intersection, then the conclusion $T_0 =
  T'_0$ will stand.

  Assume first that $T_1 = T'_1$.  Then $T_1\Ra T_0$ and $T_1\Ra
  T'_0$, which gives $T_0=T'_0$.

  Next assume that $T_1=T'_0$.  We have $T'_1\Ra T'_0 = T_1 \Ra T_0$.
  By Lemma~\rif{calc:pseudo1} and
  Lemma~\rif{calc:pseudo2}, this puts incompatible constraints on the
  length of the shared edge between $T_0$ and $T_1$.  It must have
  length less than $1.8$ and greater than $1.8$.

  The case $T_0=T'_1$ follows from the previous case by symmetry.
\end{proof}

\begin{lemma}[pseudo-dimer-obtuse]\libel{lemma:pd-obtuse}
  Assume $(T_1,T_0)\in\PD$.  Then each edge of $T_0$ and $T_1$ has
  length less than $\kappa\sqrt8$.  In particular, if $T$ is obtuse
  then $T\nRa T_0$ and $T \nRa T_1$.
\end{lemma}

\begin{proof} 
  The shared edge between $T_0$ and $T_1$ has length less than $1.8 <
  \kappa\sqrt8$ (Lemma~\rif{calc:pseudo1}).  The shared edge is the
  longest edge of $T_1$, so that each edge of $T_1$ has length less
  than $1.8$.  By Lemma~\ref{calc:pseudo3}, the only possibility is the egressive edge of $T_0$.
  But if the egressive edge of $T_0$ has length at least
  $\kappa\sqrt8$, then by Lemma~\ref{lemma:172},
\[
\area\{T_1,T_0\} > \ao + \area(\kappa\sqrt8,2\kappa,1.72) 
> \ao + (\acrit + \epsN) = 2\acrit.
\]
This contradicts the area condition in the definition of pseudo-dimer.
This completes the first claim of the lemma.

If $T$ is obtuse, then its longest edge has length at least
$\kappa\sqrt8$.  If $T\Ra T_i$, then the shared edge has length at
least $\kappa\sqrt8$, contrary to the first claim of the lemma.
\end{proof}

\begin{corollary}\libel{lemma:pd-n2}
If $(T_1,T_0)\in\PD$, then $n_-(T_0)=0$.
\end{corollary}

\begin{proof} 
  Suppose for a contradiction that $n_-(T_0)>0$.  Then $(T_+,T_0)\in
  \N$ for some $T_+$.  Because $T_0$ is nonobtuse,
  condition (4) in the definition of $\N_{\text{nonobtuse}}$
  gives an obtuse triangle $T$ such that $T\Ra T_0$.  This is contrary to the
  previous lemma.
\end{proof}

\begin{lemma}\libel{lemma:pd-1} 
  Let $(T_1,T_0)\in\PD$.  Then
  $n_+(T_1)=n_-(T_1)=m_+(T_1)=m_-(T_1)=0$.  In particular, $b(T_1) =
  \area(T_1)$.
\end{lemma}

\begin{proof}  
  We claim that $m_+(T_1)=m_-(T_1)=0$.  Otherwise by the definition of $\M$, we find that $T_1$ shares an
  egressive edge of length greater than $1.8$ with some pseudo-dimer. See Lemma~\ref{calc:pseudo2}.
  However, the longest edge of $T_1$ is its shared edge with $T_0$,
  which has length less than $1.8$ (by Lemma~\ref{calc:pseudo1}).  This gives the claim.

  Next we claim that $n_+(T_1)=0$.  Otherwise, $(T_1,T_0)\in \N$.
  Because $T_0$ is nonobtuse, the definition of $\N_{\text{nonobtuse}}(4)$
  implies that there exists an obtuse
  triangle $T$ such that $T\Ra T_0$.  This contradicts
  Lemma~\rif{lemma:pd-obtuse}.

  Finally, we claim that $n_-(T_1)=0$.  Otherwise, $(T_+,T_1)\in \N$ for
  some $T_+$.  Because $T_1$ is nonobtuse, by the definition of $\N_{\text{nonobtuse}}(4)$, this implies that there exists
  an obtuse triangle $T$ such that $T\Ra T_1$.  This contradicts
  Lemma~\rif{lemma:pd-obtuse}.
\end{proof}

\begin{lemma}\libel{lemma:pd-b}  
  Let $(T_1,T_0)\in\PD$.  Then $b(T_0) > \acrit$.  That is, $T_0$ is
  not $b$-subcritical.
\end{lemma}

\begin{proof}
  We recall from Corollary~\rif{lemma:pd-n2} that $n_-(T_0)=0$.  From
  Lemma~\rif{lemma:pd-m2} we have $m_-(T_0)=0$.  Thus, $b(T_0) =
  \area(T_0) + \epsN n_+(T_0) + \epsM m_+(T_0)$.  We have
  (by Lemma~\rif{calc:pseudo-area} and the definition of pseudo-dimer)
\[
\area(T_0) = \area\{T_0,T_1\} - \area(T_1) 
> (2\acrit-\epsM) -\acrit = \acrit - \epsM.
\]

We first treat the case that $n_+(T_0)>0$ or $m_+(T_0)>0$.  Then we
have $n_+(T_0)+m_+(T_0)\ge 1$.  Thus,
\[
b(T_0) \ge \area(T_0) + \epsM (n_+(T_0)+m_+(T_0)) > (\acrit - \epsM)
+ \epsM \ge \acrit.
\]
This completes this case.

For the remainder of the proof, we assume that $n_+(T_0)=m_+(T_0)=0$.  In
particular, we have $b(T_0) = \area(T_0)$.

Let $T_0\Ra T_-$.  We have $(T_0,T_-)\not\in\M$.  By the definition of $\M$ and Lemma \rif{calc:pseudo2}, 
the uniqueness property of $\M$ condition 4 specifically fails.  That is,
there exists $T'_1\ne T_1$ such that $(T'_1,T_0)\in\PD$.  The edges of
$T_0$ shared with $T_1$ and $T'_1$ have length at least $1.72$ and
less than $1.8$ (by Lemma \rif{calc:pseudo1} and Lemma \ref{lemma:172}).  The egressive edge of $T_0$ has length at least
$1.8$.  Thus, by monotonicity,
\[
b(T_0) = \area (T_0) \ge \area(1.72,1.72,1.8) > \acrit.
\]
This completes the proof.
\end{proof}

\section{Main inequality for dimers}\libel{sec:dimer-main}

Recall that Section~\rif{sec:dimer-pair} shows that there exists a unique
dimer pair (up to congruence): the ice-ray dimer.  In particular, if
$(T_1,T_0)$ is a dimer, then $\area(T_0)=\area(T_1)=\acrit$, and
$T_0\Ra T_1$, and $T_1\Ra T_0$, and $(T_0,T_1)$ is a dimer too.  Thus,
the relationship between $T_0$ and $T_1$ in a dimer pair is
symmetrical.  This section proves the following theorem.

\begin{theorem}\libel{thm:dimer} 
  Let $(T_1,T_0)$ be a dimer pair.  Then for $T\in\{T_0,T_1\}$, we have
  $m_+(T)=m_-(T) = n_+(T)=n_-(T) = 0$; and $b(T) = \area(T)$.
  Moreover, $\{T_0,T_1\}$ is a cluster, and the main inequality holds
  for $\{T_0,T_1\}$.
\end{theorem}

The proof will occupy the entire section.  Before treating dimers, we
treat the easy case of singletons.

\begin{lemma}\libel{lemma:singleton}  
  Assume that all the edges of a $P$-triangle $T$ have length less
  than $1.72$.  Then the cluster of $T$ is the singleton $\{T\}$ and
  $b(T) =\area(T)$.  Generally, if $\{T\}$ is a singleton
  cluster, then $b(T) > \acrit$ and the main inequality holds.
\end{lemma}

\begin{proof} 
  Assume that all the edges of $T$ have length less than $1.72$.  The
  constant $1.72$ is built into the definition of the constants
  $n_\pm,m_\pm$. This gives $n_+(T)=n_-(T) = m_+(T)=m_-(T)=0$.  Thus,
  $b(T)=\area(T)$.  We have $\area(T) > \acrit$ by
  Lemma~\rif{lemma:172}.  Because $T$ is not $b$-subcritical, there are
  no arrows $T\rab -$.

  We claim that there cannot exist an arrow $T_+\rab T$.  Otherwise,
  the longest edge of $T_+$ has length less than $1.72$ and again
  $T_+$ is not $b$-subcritical.  The claim follows and the cluster is
  a singleton.

  In general, if $\{T\}$ is any singleton cluster, there is no arrow
  $T\rab -$.  This implies that $T$ is not $b$-subcritical, so that
  $b(T) > \acrit$.  The result follows.
\end{proof}

\subsection{Interaction of dimers and pseudo-dimers}

Next, we show the disjointness of dimers from pseudo-dimers.

\begin{lemma}\libel{lemma:pd-nra-dimer}
  Let $(T_1,T_0)\in DP$ and let $(T'_1,T'_0)\in\PD$.  Then for
  $T\in\{T_0,T_1\}$, we have $T'_0\nRa T$.
\end{lemma}

\begin{proof} 
Assume that $T'_0\Ra T\in\{T_0,T_1\}$. The triangle $T$ is
nonobtuse, so Corollary~\rif{lemma:egress'} gives
$\area(T)>\acrit+\epsM$. On the other hand,
Theorem~\rif{thm:dimer-area} gives $\area(T)=\acrit$.
This is a contradiction.
\end{proof}

\begin{corollary}\libel{lemma:m2-dimer}  
  Let $(T_1,T_0)\in DP$. Then %
  $m_-(T_0)=m_-(T_1)=0$.
\end{corollary}

\begin{proof}  
  If $m_-(T_-)>0$, then by the definition of the set $\M$, there
  exists $(T_1',T_0')\in \PD$, such that $T_0'\Ra T_-$.  %
Lemma \rif{lemma:pd-nra-dimer} implies $T_-\not\in\{T_0,T_1\}$.
\end{proof}

\begin{lemma}\libel{lemma:m1-dimer}  
  Let $(T_1,T_0)\in DP$ and let $(T'_1,T'_0)\in\PD$.  Then
  $\{T_0,T_1\}\cap \{T'_0,T'_1\} = \emptyset$.  Moreover,
  $m_+(T_0) = m_+(T_1)=0$.
\end{lemma}

\begin{proof}
  The relationship between $T_0$ and $T_1$ in a dimer pair is
  symmetrical.  It is enough to show $T_1\not\in\{T_0',T_1'\}$.  We
  claim that $T_1\ne T'_0$.  Otherwise, $T'_0\Ra T_0$, which is
  contrary to Lemma~\rif{lemma:pd-nra-dimer}.

  We claim that $T'_1\ne T_1$.  Otherwise, if $T_1=T'_1$, then
  $T'_1=T_1\Ra T_0$ and $T'_1\Ra T'_0$, so that $T_0 = T'_0$, which
  gives $(T_1,T_0)=(T_1',T_0')\in DP\cap\PD=\emptyset$, which is impossible. This proves the disjointness result.

  If $m_+(T) >0$, then there exists $T_1''$ such that $(T_1'',T)\in
  \PD$.  This is impossible for $T\in\{T_0,T_1\}$ by the disjointness
  result.
\end{proof}

\subsection{Interaction with obtuse triangles}

\begin{lemma}\libel{lemma:t2b-nonobtuse}
  If $T_1$ is obtuse, $T_0$ is nonobtuse, and if $T_1 \Ra T_0$, then
  $T_0$ is not $b$-subcritical.
\end{lemma}

\begin{proof} 
  If $T_1$ is obtuse, then its longest edge, which is shared with
  $T_0$, has length at least $\kappa\sqrt8$.

  Assume for a contradiction that $T_0$ is nonobtuse and
  $b$-subcritical.  By Lemma~\rif{lemma:right}, $\area(T_0) > \acrit$.
  Thus, we must have $n_-(T_0)>0$ or $m_-(T_0)>0$.  We consider two
  cases, according to which of these inequalities occurs.

  Assume in the first case that $n_-(T_0)>0$; that is, $(T',T_0)\in
  \N$.  Recall that this implies $m_-(T_0)=0$ by
  Lemma~\rif{lemma:n2m2}.  We have $T_1\Ra T_0$ and $(T_1,T_0)\not\in
  \N$, because $T_1$ is obtuse.  It follows that $n_-(T_0)\le 2$.  In
  fact, $n_-(T_0)$ is equal to the number of nonobtuse $T$ with
  longest edge of length at least $1.72$ such that $T\Ra T_0$.  If
  $n_-(T_0)=1$, we have
  \[
b(T_0) \ge \area(T_0) - \epsN 
> \area(\kappa\sqrt8,1.72,2\kappa)  - \epsN > \acrit.
\]
If $n_-(T_0)=2$, we have
  \[
b(T_0) \ge \area(T_0) - 2\epsN 
> \area(\kappa\sqrt8,1.72,1.72)  - 2\epsN > \acrit.
\]
  This completes the first case.

  Finally, we consider the case that $m_-(T_0)>0$ (and $n_-(T_0)=0$).
  We have $m_-(T_0)\le 2$, because $T_1\Ra T_0$, and $T_1$ is obtuse,
  and cannot be part of a pseudo-dimer.  The following area estimate
  gives the result.
\[
  b(T_0) \ge \area(T_0) - 2\epsM\ 
  > \area(\kappa\sqrt8,1.8,2\kappa) -
  2\epsM > \acrit.
\]
The constant $1.8$ comes from the egressive edge of a pseudo-dimer in
the definition of $\M$ and $m_-$.  See Lemma~\rif{calc:pseudo2}.
\end{proof}

\begin{corollary}\libel{lemma:n2b} 
  If $n_-(T_-)>0$ with $T_-$ nonobtuse, then $b(T_-) > \acrit$.
\end{corollary}  

\begin{proof}  
  By the definition of $\N$ (nonobtuse target), There exists an obtuse
  triangle $T$ such that $T\Ra T_-$.  The result follows from the
  lemma.
\end{proof}

\begin{lemma}\libel{lemma:m2b}  
  Let $T$ be $b$-subcritical and nonobtuse.  Then $n_-(T)=0$.
  Moreover, assume that there exists $T'$ that is $b$-subcritical such
  that $T'\rab T$ or $T\rab T'$.  Then $\area(T)\le \acrit$.
\end{lemma}

\begin{proof}  
  By the contrapositive of the corollary, it follows that $n_-(T)=0$.

  Assume for a contradiction that $\area(T) > \acrit$.  We have
\begin{equation}\libel{eqn:2-ac}
\area(T) > \acrit \ge b(T) \ge \area(T) - \epsM m_-(T).
\end{equation}
Thus, $m_-(T) >0$.  By the definition of $\M$, there exists
$(T'_1,T'_0)\in\PD$ with $(T'_0,T)\in\M$ and $T'_0\Ra T$.  By
Corollary~\rif{lemma:egress'}, we have $\area(T) >
\acrit+\epsM$. Combined with Inequality ~\eqref{eqn:2-ac}, this gives
$m_-(T)\ge 2$.  Choose a second incoming pair, distinct from the
first.  Its source $T''_0\ne T'_0$ is the central triangle of a
pseudo-dimer $(T''_1,T''_0)\in\PD$, with $(T''_0,T)\in\M$ and
$T''_0\Ra T$.  The triangles $T'$, $T'_0$,
and $T''_0$ are distinct, because for example $b(T'_0) > \acrit \ge
b(T')$ (Lemma~\rif{lemma:pd-b}).  Because of the arrow $T'\rab T$ or $T\rab
T'$, the triangles $T$ and $T'$ belong to the same cluster.  The
shared edge between $T$ and $T'$ is the longest edge of the source
of this arrow.  If it had length less than $1.72$,
Lemma~\rif{lemma:singleton} would contradict the source being
$b$-subcritical.  It therefore has length at least $1.72$.  The
other two edges are greater than $1.8$ by
Lemma~\rif{calc:pseudo2}.  Then
\[
b(T) \ge \area(T) - m_-(T)\epsM > \area(1.8,1.8,1.72) - 3\epsM > \acrit.
\]
This  contradicts the assumption that $T$ is $b$-subcritical.
\end{proof}

\begin{lemma} \libel{lemma:rab-sequence} Assume that $T_1$ and $T_0$
  are both nonobtuse.  Then there does not exist a sequence $T_1\rab
  T_0 \rab T$, with $T_1\ne T$.
\end{lemma}

\begin{proof}  
  Assume for a contradiction that the sequence exists.  By the
  Lemma~\rif{lemma:m2b}, $n_-(T_0)=0$ and $\area(T_0)\le \acrit$.

  We claim that $m_+(T_0)=0$.  Otherwise, there exists
  $(T'_1,T_0)\in\PD$, and by Lemma~\rif{lemma:pd-b}, $T_0$ is not
  $b$-subcritical, contradicting the assumptions of the lemma.

  We claim that $m_-(T_0)=0$.  This follows by
  Lemma~\rif{lemma:m2-area} and the fact $\area(T_0)\le \acrit$.

We have that $n_+(T_0)=0$.  Otherwise, we reach the contradiction,
\[
\acrit \ge b(T_0) =\area(T_0) + \epsN n_+(T_0) > \ao + \epsN = \acrit.
\]
This shows that $\area(T_0) = b(T_0)$.

By Lemma~\rif{lemma:m2b}, we have $\area(T_1)\le \acrit$.
It follows that $(T_1,T_0)\in\PD$, and by Lemma~\rif{lemma:pd-b},
we reach a contradiction to the assumption that $T_0$ is $b$-subcritical. 
\end{proof}
  
We are finally in a position to prove Theorem~\rif{thm:dimer}.

\begin{proof}[Proof of dimer theorem~\rif{thm:dimer}.]
  Let $(T_1,T_0)$ be a dimer pair, and let $T\in \{T_0,T_1\}$.  We
  have proved that $m_+(T)=m_-(T)=0$ in Lemma~\rif{lemma:m1-dimer} and
  Corollary~\rif{lemma:m2-dimer}.

  We claim that $n_-(T_0)=n_-(T_1)=0$.  Otherwise, $T_1$ or $T_0$ has
  an edge of length at least $\kappa\sqrt8$, coming from an edge
  shared with an obtuse triangle, and the triangle $T_0$ or $T_1$ is
  not subcritical (Lemma~\rif{lemma:right}).

  We claim that $n_+(T_0)=n_+(T_1)=0$.  Every outgoing $\N$
  pair follows the unique attachment arrow.  In a dimer, $T_i$
  attaches to $T_{1-i}$.  Thus, after interchanging the indices,
  any nonzero outgoing term gives $(T_1,T_0)\in\N$, contrary
  to $n_-(T_0)=0$.

  This shows that $b(T_i)=\area(T_i)$, for $i=0,1$.  Because $T_0$ and
  $T_1$ are both subcritical, they are also $b$-subcritical, and fall
  into the same cluster.  This is the full cluster, for otherwise, we
  would have say $T\rab T_0\rab T_1$, which is impossible by
  Lemma~\rif{lemma:rab-sequence}, for $T$ nonobtuse.  And if $T$ is
  obtuse, this contradicts Lemma~\rif{lemma:t2b-nonobtuse}.
  This gives the proof.
\end{proof}

\subsection{Cluster structure}

We continue with our analysis of the clusters in a fixed saturated
packing of regular pentagons.  

\begin{lemma}\libel{lemma:notlong}  
  Let $T'$ and $T$ be $P$-triangles, such that $T'\Ra T$, where
  $T'$ is nonobtuse subcritical and $T$ is obtuse.  Then the edge
  of attachment is not the longest edge of $T$.
\end{lemma}

\begin{proof} We have seen that each edge of a nonobtuse subcritical
  triangle has length less than $\kappa\sqrt8$ (Lemma~\rif{lemma:right}) and that the longest
  edge of an obtuse Delaunay triangle has length at least
  $\kappa\sqrt{8}$.  These are incompatible conditions on a shared
  edge.
\end{proof}

\begin{lemma}\libel{lemma:3m2} 
If $m_-(T) = 3$, then $b(T) > \acrit + \epsN$.
\end{lemma}

\begin{proof} By Lemma~\rif{lemma:n2m2}, $n_-(T)=0$.  There is a
  longest (egressive) edge of a pseudo-dimer along each edge of $T$,
  each of length at least $1.8$ (Lemma~\rif{calc:pseudo2}).  This gives
\[
b(T) \ge \area(T) - 3\epsM 
\ge \area(1.8,1.8,1.8) - 3\epsM > \acrit + \epsN.
\]
\end{proof}

\begin{lemma}\libel{lemma:Nb}  
  If $(T_+,T_-)\in \N$, then $b(T_+)>\acrit$.
\end{lemma}

\begin{proof}  Assume for a contradiction that $b(T_+)\le\acrit$ and $(T_+,T_-)\in\N$. We have $n_+(T_+)\ge 1$, and
\[
\acrit \ge b(T_+) 
\ge \area(T_+) + \epsN (1 - n_-(T_+)) - \epsM m_-(T_+)
> \acrit - \epsN n_-(T_+) - \epsM m_-(T_+).
\]
So $n_-(T_+) > 0$ or $m_-(T_+)>0$. This gives two cases.  

Suppose that $n_-(T_+) >0$.  Recall that $T_+$ is nonobtuse.  Then
$(T,T_+)\in\N$ for some nonobtuse $T$ whose shared edge with $T_+$ has
length at least $1.72$.  Thus, by the definition of $\N$ (nonobtuse
target), there exists $T'$ obtuse such that $T'\Ra T_+$.  The shared
edge has length at least $\kappa\sqrt8$.  Also, $n_-(T_+)\le 2$
(because $(T',T_+)\not\in \N$).  By Lemma~\rif{lemma:n2m2}, we have
$m_-(T_+)=0$.  Then we reach a calculation
\[
b(T_+) \ge \area(T_+) + \epsN (1 - 2) 
\ge \area (2\kappa,1.72,\kappa\sqrt8) - \epsN > \acrit.
\]
This completes this case.

Finally, suppose that $m_-(T_+)>0$ and $n_-(T_+)=0$.  %
We have by Lemma~\rif{lemma:m2-area},
\[
b(T_+) \ge \area(T_+)+\epsN - \epsM m_-(T_+) 
> (\acrit + \epsM) + \epsN -3\epsM  \ge \acrit.
\]
\end{proof}

\begin{lemma}\libel{lemma:no-ao} 
  There is no arrow $T_1 \rab T_0$ with $T_1$ nonobtuse and $T_0$
  obtuse.
\end{lemma}

\begin{proof}  
  Assume for a contradiction that such a pair $(T_1,T_0)$ exists.  By
  Lemma~\rif{lemma:Nb}, $(T_1,T_0)\not\in\N$.  By the definition of
  $\N$ (obtuse target), the longest edge of $T_1$ has length less than
  $1.72$.  By Lemma~\rif{lemma:singleton}, $T_1$ forms a singleton
  cluster.  This contradicts $T_1\rab T_0$.
\end{proof}

\section{Obtuse Clusters}\libel{sec:obtuse}

In this section we prove the strict main inequality for clusters that
contain an obtuse triangle.  This will involve several cases, but in
every case the strict main inequality will be found to hold by a large
margin.  This allows us to use rather crude approximations of
area in this section.  In particular, we are able to disregard most
constraints on the shapes of Delaunay triangles imposed by the
pentagons.  Instead, we use generic features of the triangles such as
the fact that the circumradius is at most two and the edge lengths of
the triangle are at least $2\kappa$.

\begin{remark}\libel{rem:delaunay}
  Recall that the Delaunay property implies that two adjacent Delaunay
  triangles $T_1$ and $T_2$ have the property that $\alpha_1 +
  \alpha_2\le \pi$, where $\alpha_i$ is the angle of $T_i$ that is not
  at the shared edge of $T_1$ and $T_2$.  Indeed, when their union
  is a convex quadrilateral, the inscribed-angle criterion says that
  the opposite vertex lies on or outside the circumcircle exactly
  when $\alpha_1+\alpha_2\le\pi$.  If the union is nonconvex,
  its reflex angle forces the sum of these two opposite angles to be
  less than $\pi$.  The straight-angle limiting case gives the same
  non-strict inequality.  In particular, {\it two
    obtuse Delaunay triangles cannot be joined along an edge that is
    the longest on both triangles.}  The extreme case
  $\alpha_1+\alpha_2=\pi$ corresponds to the degenerate situation
  where $T_1$ and $T_2$ form a cocircular quadrilateral. When
  cocircular, either diagonal of the quadrilateral gives an acceptable
  Delaunay triangulation.
\end{remark}

We write $\areta(d_1,d_2,h)$ for the area of a triangle with two edges
$d_1,d_2$ and circumradius $h$.  In general, two noncongruent
triangles have data $d_1,d_2,h$ (Figure~\rif{fig:areta}).  We choose $\areta(d_1,d_2,h)$ to
give the area of that triangle such that its third edge $d_3$ is as
long as possible.

The following lemma shows that under quite general
conditions, we are justified in our decision to choose $d_3$ as long
as possible in the definition of the function $\areta(d_1,d_2,h)$.  It
is justified in the sense that the other choice does not usually give
a Delaunay triangle of a pentagon packing, according to the following
simple test.

\begin{lemma}\libel{lemma:areta}  
  Let $d_1$, $d_2$ and $\eta$ be positive real numbers.  Assume that
  $T$ and $T'$ are triangles with edge lengths $d_1,d_2,d_3$ and
  $d_1,d_2,d_3'$, and with the same circumradius $\eta$. Assume
  $2\kappa\le d_1\le d_2$.  Set $\theta =
  \arc(\eta,\eta,d_1)+\arc(\eta,\eta,2\kappa)$.  If $2\kappa \le d_3'
  < d_3$, then $\theta < \pi$ and $2\eta\sin(\theta/2) \le d_2$.
\end{lemma}

As a corollary, 
in contraposition, if $\theta\ge\pi$ or if $d_2 < 2\eta\sin(\theta/2)$,
then the triangle $T'\ne T$, with $d_3' < d_3$,
cannot satisfy the constraint $2\kappa\le d_3'$ of a  Delaunay triangle.

\begin{proof} 
  See Figure~\rif{fig:areta}.  Let $\p$, $\q$, and $\r$ (resp. $\p$,
  $\q$, and $\r'$) be the vertices of $T$ (resp. $T'$) on a common
  circle, with
\[
\norm{\p}{\q} =d_1, 
\quad\norm{\p}{\r}=\norm{\p}{\r'}=d_2, \quad\text{and } 
\norm{\q}{\r'} = d_3' \le d_3=\norm{\q}{\r}.
\]
The angle on the circumcircle from $\p$ to $\q$ is
$\arc(\eta,\eta,d_1)$, and $\theta$ is the angle on the circumcircle
from $\p$ to the first point $\s$ beyond $\q$ such that
$\norm{\s}{\q}=2\kappa$.  If $\theta\ge\pi$, a point $\r'\ne \r$
satisfying the constraints $2\kappa\le d_3'<d_3$ does not exist.  Assume $\theta < \pi$.  As
the figure indicates, $\norm{\p}{\r'}$ is minimized (as a function of
$d_2$) when $\r'=\s$, and $d_3' = 2\kappa$, the lower constraint.
Then $d_2 = \norm{\p}{\r'}\ge \norm{\p}{\s}= 2\eta\sin(\theta/2)$.
\end{proof}

\tikzfig{areta}{
There can be two positions $\r,\r'$ 
on the circumcircle for
the third vertex of the triangle.  
Here, $d_1 = \norm{\p}{\q}$.}{
[scale=1.3]
\draw (0,0) circle (1cm);
\draw (1,0) node[anchor=west] {$\p$} --  (70:1cm) node[anchor=south] {$\q$};
\draw (1,0) -- node[above=1pt] {$d_2$} (130:1cm) node[anchor=south] {$\r'$};
\draw (1,0) -- node[below=1pt] {$d_2$} (-130:1cm) node[anchor=north] {$\r$};
\draw (130:1cm) -- (-130:1cm);
\draw[thin,gray] (1,0) -- (-1,0);
\draw (100:1cm) node[above=1pt] {$\s$};
}

\begin{lemma}\libel{lemma:obtuse-area-monotone}
  Let $T$ be an obtuse triangle, with sides $a,b$ adjacent to its
  obtuse angle and circumradius $r$.  If $a\ge a_0>0$,
  $b\ge b_0>0$, and $r\le r_0$, then
  \[
  \area(T)\ge\areta(a_0,b_0,r_0).
  \]
\end{lemma}

\begin{proof}
  Increasing the circumradius keeping $a,b$ fixed increases the length of
  the edge opposite the obtuse angle.  This is decreasing in area.
  Decreasing $b$ while fixing the circumradius
  decreases the height of the triangle measured from the base of length $a$,
  and thus decreases area.
  Decreasing $a$ while fixing the circumradius
  decreases the height of the triangle measured from the base of length $b$,
  and thus decreases area.
  These operations preserve obtuseness.
\end{proof}

\begin{lemma}\libel{lemma:t2b}
  If $T_1$ is obtuse, and $T_1 \Ra T_0$, then $T_0$ is not
  $b$-subcritical.
\end{lemma}

\begin{proof}  
  If $T_0$ is nonobtuse, then this is Lemma~\rif{lemma:t2b-nonobtuse}.

  Assume that $T_0$ is obtuse.  By basic properties of Delaunay
  triangles (Remark~\rif{rem:delaunay}), obtuse Delaunay triangles never join
  along an edge that is the longest on both triangles.  Thus, $T_1$
  attaches to $T_0$ along an edge adjacent to the obtuse angle of
  $T_0$.  The shared edge has length at least $\kappa\sqrt8$.  
  Recall that the circumradius of a Delaunay triangle is at most $2$,
  because of the saturation assumption.  The sides adjacent to the
  obtuse angle of $T_0$ are at least $2\kappa$ and
  $\kappa\sqrt8$.  Lemma~\rif{lemma:obtuse-area-monotone} gives
  $\area(T_0)\ge\areta(2\kappa,\kappa\sqrt8,2)$.
  Since $T_1$ is obtuse, $(T_1,T_0)\notin\N$.  Its shared edge
  occupies one of the three edges of $T_0$, so $n_-(T_0)\le2$.  Then a numerical calculation (using %
  Lemma~\rif{lemma:obtuse-b}) gives
\[
b(T_0) \ge \area(T_0) - \epsN n_-(T_0) \ge 
\areta(2\kappa,\kappa\sqrt{8},2) - 2\epsN > \acrit.
\] 
Here the obtuse-angle hypothesis identifies the branch of $\areta$
directly, as in Lemma~\rif{lemma:obtuse-area-monotone}.
\end{proof}

For an obtuse triangle we use Lemma~\rif{lemma:obtuse-area-monotone}
to justify the branch and monotonicity of $\areta$.  When obtuseness
is not prescribed, the test following Lemma~\rif{lemma:areta}
distinguishes the branches.

\begin{lemma}\libel{lemma:sequence}  There does not exist a three
  term sequence $- \rab -\rab -$ where the three triangles are
  distinct.
\end{lemma}

\begin{proof}  
  Assume for a contradiction, that such a sequence exists.  By
  Lemma~\rif{lemma:t2b}, there does not exist a sequence $- \rab -
  \rab -$, where the first triangle is obtuse.  Thus, we may assume
  that the first triangle is nonobtuse.  By Lemma~\rif{lemma:no-ao},
  there does not exist $T_1\rab T_0$, where $T_1$ is nonobtuse and
  $T_0$ is obtuse.  Thus, we may assume that every triangle in the
  sequence is nonobtuse.  This is impossible by
  Lemma~\rif{lemma:rab-sequence}.
\end{proof}

If $\C$ is a cluster that is not a singleton, then there is
some arrow $T_1\rab T_0$.  We have the following structure
theorem for clusters.

\begin{theorem}\libel{lemma:c-weak}
  Let $\C$ be a cluster, and let $T_1\rab T_0$ be an
  arrow between triangles in $\C$.
  Then 
  \begin{equation}\libel{eqn:c-weak}
  \C = \{T_0\}\cup \{T \mid T\rab T_0\}.
  \end{equation}
\end{theorem}

\begin{proof}  
  We use Lemma~\rif{lemma:sequence}.  Assume that $T\rab T_0$.  There
  is no arrow $T'\rab T$, with $T'\ne T_0$, because that would also
  produce a sequence $T' \rab T \rab T_0$ of three distinct triangles.
  There is a unique arrow out of $T$.  Thus, we have accounted for all
  of the arrows in and out of $T$.

  There is no arrow $T_0\rab T''$, with $T''\ne T_1$, because that
  would produce a sequence $T_1\rab T_0\rab T''$ of  three 
  distinct triangles (Lemma~\rif{lemma:sequence}).  We have accounted for all the arrows in and
  out of $T_0$.  Thus, the full cluster has been identified.
\end{proof}

\begin{corollary}\libel{lemma:card4}
  Every cluster is finite of cardinality at most $4$.
\end{corollary}

\begin{proof} At most three triangles attach to $T_0$.
\end{proof}

The following theorem is the main result of this
section.

\begin{theorem}\libel{thm:obtuse}  
  Let $\C$ be any cluster that contains an obtuse triangle.  Then 
   the strict main inequality (\rif{eqn:strict-main})
  holds for $\C$.
\end{theorem}

We prepare for the proof of the theorem with some lemmas.

\begin{lemma} \libel{lemma:obtuse-source} 
  Let $T_1\rab T_0$ be an
  arrow between two triangles in a cluster that contains an obtuse
  triangle.  Then $T_1$ is obtuse.  Moreover, for every obtuse
  triangle $T$ in the cluster, $b(T) = \area(T) -\epsN n_-(T)$.
\end{lemma}

\begin{proof} 
  Assume for a contradiction that $T_1$ is nonobtuse.  By
  Lemma~\rif{lemma:no-ao} applied to the arrow $T_1\rab T_0$, the
  triangle $T_0$ is nonobtuse.  By assumption and the structure
  theorem for clusters (Theorem~\rif{lemma:c-weak}), there exists $T'\rab T_0$, where $T'$ is
  obtuse.  The singleton lemma (Lemma~\rif{lemma:singleton}) implies
  that the longest edge of $T_1$, which is shared with $T_0$ has
  length at least $1.72$.  By the definition of $\N$ (nonobtuse
  target) and Lemma~\rif{lemma:obtuse-b}, we have $(T_1,T_0)\in\N$.  By Lemma~\rif{lemma:Nb}, we have
  $b(T_1)>\acrit$, and $T_1$ is not $b$-subcritical.  Thus, we obtain
  a contradiction to $T_1\rab T_0$.

 Thus, $T_1$ is obtuse.
  Moreover, if $T$ is obtuse, then $b(T) = \area(T) -\epsN n_-(T)$ by
  Lemma~\rif{lemma:obtuse-b}.
\end{proof}

Let $\C$ be a cluster.  By {\it external edge} of the cluster, we mean an
edge of a triangle in the cluster that is not shared with another
triangle in the cluster.  Let $\tn(\C)$ be the total number of
external edges of length at least $1.72$ of the cluster $\C$.  We have
\[
\tn(\C) =\sum_{T\in \C} \tn(T),
\]
where $\tn(T)$ is the number of edges of $T$ of length at least $1.72$
that are external edges of its cluster.
Define $\tb(T) := \area(T)-\epsN\tn(T)$.  We use $\tb(T)$ to give an
easily computed lower bound given in the following lemma.

\begin{lemma} \libel{lemma:tb}
  Let $\C$ be a cluster of cardinality at least $2$ 
  that contains an obtuse triangle.
Then 
\[
\sum_{T\in\C} b(T)\ge \sum_{T\in\C} \tb(T).
\]
\end{lemma}

\begin{proof}
  Let $\N_{ext}\subseteq \N$ be the subset consisting of pairs
  $(T_+,T_-)$ such that at least one of $T_+$ and $T_-$ is not in
  $\C$.  Define $\M_{ext}\subseteq \M$ similarly.  If $(T_+,T_-)\in\N\setminus\N_{ext}$,
  where $\{T_+,T_-\}$ meets $\C$, then $T_+,T_-\in\C$ and the pair $(T_+,T_-)$
  contributes $+\epsN$ to the value of $b(T_+)$ and $-\epsN$ to the
  value of $b(T_-)$.  These contributions cancel.  Similar comments
  apply to $(T_+,T_-)\in\M\setminus \M_{ext}$.  Thus,
\[
\sum_{T\in \C} b(T) = \sum_{T\in \C} b_{ext}(T),
\]
where $b_{ext}(T)$ is defined as $b(T)$, but using $\N_{ext}$ and
$\M_{ext}$ instead of $\N$ and $\M$. It is enough to show that
\[
b_{ext}(T) \ge \tb(T),
\]
for all $T\in\C$.

Let $T\in\C$.  To prove the inequality for $T$, it is enough to show
that $\tn(T)\ge n_{ext,-}(T)+m_{ext,-}(T)$.  In fact, this inequality gives
\begin{align*}
b_{ext}(T) &\ge \area(T) - \epsN n_{ext,-}(T) - \epsM m_{ext,-}(T) \\
&\ge \area(T) - \epsN (n_{ext,-}(T)+m_{ext,-}(T)) \\
&\ge \area(T) - \epsN \tn(T)\\
&= \tb(T).
\end{align*}
The sum $n_{ext,-}(T)+m_{ext,-}(T)$ counts pairs
$(S,T)\in\N_{ext}\sqcup\M_{ext}$ with $S\notin\C$.
Each determines a shared external edge of length at least $1.72$.
There is only one neighboring triangle across a given edge, and
$\N\cap\M=\emptyset$.  Thus, distinct counted pairs determine
distinct external edges.  This gives the required injection into
the edges counted by $\tn(T)$ and completes the proof.
\end{proof}

\begin{lemma}\libel{lemma:nonobtuse-radius-area}
  Let $T$ be a nonobtuse triangle whose sides are at least $a>0$.
  If its circumradius $r$ satisfies $r\ge r_0\ge a/\sqrt2$, then
  \[
  \area(T)\ge\frac {a}2\sqrt{4r_0^2-a^2}.
  \]
\end{lemma}

\begin{proof} Begin by fixing the circumradius.  Fixing any one
  edge length and varying the other two edge lengths, for fixed circumradius, the maximum
  area is obtained when the two edge lengths are equal, and the
  minimum area is obtained when one of the two edge lengths is as
  small as possible.  Thus, for a nonobtuse triangle of fixed circumradius,
  the minimum is at a boundary: either a right triangle with hypotenuse $2r$
  and legs $a,\sqrt{(2r)^2-a^2}$, or a triangle with two edges $a$.
  The assumption $a\le r\sqrt2$ rules out the latter possibility among
  acute triangles.  Thus, the right-triangle extremal gives
     \[
   \area(T)\ge\frac {a}2\sqrt{4r^2-a^2}
   \ge\frac {a}2\sqrt{4r_0^2-a^2}.
   \]
\end{proof}

\begin{proof}[proof of Theorem~\rif{thm:obtuse}.]
  Let $\C$ be any cluster that contains an obtuse triangle.
  The theorem asserts that the strict main inequality holds for $\C$.
  The proof involves several relatively simple cases.  

  We recall that each Delaunay triangle has edge lengths at
  least $2\kappa$ and circumradius at most $2$.

  If the cluster is a singleton $\{T\}$, where $T$ is obtuse, then the
  singleton lemma (Lemma~\rif{lemma:singleton}) gives the result.  We
  now assume that $\C$ is not a singleton.  By
  Theorem~\rif{lemma:c-weak}, the cluster $\C$ has the form of
  Equation~\rif{eqn:c-weak} for some triangle $T_0$.  Each $T$ such
  that $T\rab T_0$ is obtuse by Lemma~\rif{lemma:obtuse-source}.

  We break the proof into six cases depending on whether $T_0$ is
  nonobtuse, and depending on $\card(\C)\in \{2,3,4\}$.  In each case
  we prove inequality
\begin{equation}\libel{eqn:n}
 \area(\C) > \acrit \card(\C) + \epsN \tn(\C).
\end{equation}
By Lemma~\rif{lemma:tb},    this implies that
\[
\sum_{T\in\C} b(T) \ge \sum_{T\in\C} \tb(T) 
\ge \area(\C) -  \epsN\tn(\C) 
> \acrit \card(\C),
\]
which is the strong main inequality.

{\bf Case 1. The triangle $T_0$ is a nonobtuse triangle, $T_1$ is obtuse, and
  $\C=\{T_0,T_1\}$.}  The triangle $T_0$ has a vertex $\v$ that is
not shared with $T_1$.  By the Delaunay property, $\v$ lies outside
the circumcircle of $T_1$.  The triangles $T_0$ and $T_1$ form a
quadrilateral $Q$ whose diagonal is the shared edge of $T_0$ and
$T_1$.  We deform the quadrilateral $Q$ to decrease its area while
maintaining the following constraints:
\begin{enumerate}
\item The vertex $\v$ lies on or outside the circumcircle of
  $T_1$. The circumradius of $T_1$ is at most $2$.
\item The edge length of the $i$th edge of $Q$ is at least
  $d_i\in\{2\kappa,1.72\}$,
where  the number of $d_i$ that equal $1.72$ is $\tn(\C)$; and
\item $T_1$ is not acute and its largest angle is opposite
  the shared edge with $T_0$.
\item $T_0$ is not obtuse.
\end{enumerate}
We drop all other constraints as we deform. (In particular, we do not
enforce the nonoverlapping of pentagons in the $P$-triangles.)  We
continue to deform $Q$ until one of the following stopping conditions hold:
\begin{enumerate}
\item $Q$ is cocircular; or
\item For all $i=1,2,3,4$, the $i$th edge of $Q$ has reached its lower
  bound $d_i$.
\end{enumerate}
We may assume that $\area(Q)\le2\acrit+\epsN\tn(\C)$. First hold $T_1$ and the shared edge fixed and decrease the nonshared sides of $T_0$ until they reach their bounds or $Q$ becomes cocircular. If $T_0$ becomes right at a shared vertex, then
\[
\area(Q)\ge\kappa^2\sqrt8+\areta(2\kappa,2\kappa,2)>2\acrit+4\epsN,
\]
a contradiction. A right angle at the nonshared vertex is either preceded by cocircularity or is itself the cocircular endpoint. Thus no other nonobtuse boundary stops this step.

Next hold $T_0$ and the shared edge fixed and decrease the nonshared sides of $T_1$. Its area decreases and its circumradius increases, until a side bound, cocircularity, or radius $2$ is reached. In this step the sides of $T_0$ other than the shared side are already at most $1.72$. Nonobtuseness gives a shared length at most $1.72\sqrt2$. If $\eta(T_1)=2$, its two acute base angles are each at least $\arcsin(\kappa/2)$ and have sum at most $\pi/2$. The sine rule then gives the shared length at least
\[
4\sin(2\arcsin(\kappa/2))=4\kappa\sqrt{1-\kappa^2/4}>1.72\sqrt2,
\]
which is impossible. Continuing the contraction of whichever nonshared side has not yet attained its bound gives the two stated terminal alternatives.

In the first subcase (cocircularity), we drop the final constraints
(acute/obtuse) and continue area decreasing deformations for $Q$ under
the constraint of a fixed circumcircle.  We note that the area of a
cocircular quadrilateral $Q$ depends only on the lengths of the edges
and not on their cyclic order on $Q$.  We may thus rearrange the edge
order as we deform.  For a given circumcircle, the area is minimized
when three of the edges attain their lower bound $d_i$.  By suitable
reordering of the edges, we may assume that $Q$ is an isosceles
trapezoid and that the fourth (free) edge is parallel to its opposite
edge on $Q$.  Write its equal legs as $a,a$,
its fixed base as $b$, and its free base as $b+2t$.
Its area and squared
circumradius are
\[
\area(Q_t)=(b+t)\sqrt{a^2-t^2},\qquad
\eta(Q_t)^2=\frac{a^2(a^2+b^2+2bt)}{4(a^2-t^2)}.
\]
The latter is increasing in $t$, and
\[
\area(Q_t)''=-\frac{ba^2+3a^2t-2t^3}{(a^2-t^2)^{3/2}}<0.
\]
Thus the minimum along this trapezoid family occurs at an endpoint:
the free side reaches its lower bound, or $\eta(t)=2$.  When
$\eta(Q)=2$, we relax the edge lengths constraints further to allow
three edges to have length $d=2\kappa$.
When the three fixed sides all have length $d=2\kappa$ and
$\eta(Q)=2$, the area is
\[
d^3(1-d^2/16)^{3/2}>3.24>2\acrit+4\epsN.
\]
The case in which every side reaches its lower bound is included
in the calculations below, including the special case when $\tn=0$.

With the circumcircle fixed, vary two arcs while keeping
the other arcs fixed. Increasing the difference between the two
varying arcs decreases their contribution to the area until a side
bound is reached. Repeating this operation gives three active side
bounds. We may then reorder the sides, since only the side bounds
and the circumradius constraint are being retained.

After reordering, let $d_4$ be the prescribed lower bound for the free side. With
the notation already used above, its permitted range is
\[
\frac{d_4-b}{2}\le t<a.
\]
In particular $t\ge(2\kappa-1.72)/2>-0.051$; it need not be
nonnegative. The displayed formulas for area and circumradius are
unchanged. The numerator $ba^2+3a^2t-2t^3$ is positive on this
whole range. For $t<0$ use $b+3t>0$; for $t\ge0$ use $t^2<a^2$.
Thus area is strictly concave. The squared circumradius is strictly
increasing, since its derivative has the sign of
$(b+t)(a^2+bt)$, which is positive on the same range.
The minimum therefore occurs when the free side reaches $d_4$ or
when the circumradius reaches $2$.

At the first extremal all four sides equal their given bounds.
Brahmagupta's formula for the area of a cyclic quadrilateral applies
without any acute/obtuse or diagonal constraint, so the cyclic column
of the table below gives the required bound, including $\tn=0$. We do
not use its right-triangle column in this subcase. At the second
extremal we relax all four side bounds to $2\kappa$; the same
fixed-circle argument makes three of them equal to $2\kappa$, and the
displayed radius-$2$ area bound applies.

In the second subcase (every edge attains its minimal length $d_i$),
the four edge lengths are fixed.  We drop the constraint that $T_0$ is
not obtuse. Put $z=x^2$, where $x$ is the diagonal length.  A triangle with
fixed adjacent sides $a,b$ has area
\[
\frac14\sqrt{4a^2b^2-(a^2+b^2-z)^2}.
\]
This is concave in $z$ on its domain.  The area of $Q$, the sum of two
such functions, is therefore concave in the \emph{squared}
diagonal length.  On a feasible interval the minimum is at an
endpoint.  The proposed endpoint reduction gives the diagonal as
small as possible (that is, $T_1$ is a right triangle -- when this
satisfies the constraint that $\v$ is outside the circumcircle of
$T_1$) or as large as possible (that is, $Q$ is cocircular).  The
cocircular case has already been considered.
Here are explicit bounds at the proposed endpoints.  Set $d_0=2\kappa$
and $d_1=1.72$, and let $k=\tn$.  The cyclic area is given by
Brahmagupta's formula with $k$ sides $d_1$ and $4-k$ sides $d_0$.
For a right $T_1$ with legs $a,b$ and other two sides $c,d$,
the total area is
\[
\frac{ab}{2}+\frac14
\sqrt{4c^2d^2-(c^2+d^2-a^2-b^2)^2}.
\]
We use only the orders satisfying $a^2+b^2\le c^2+d^2$, as
required by the Delaunay condition at the right endpoint.  Substituting
$a,b,c,d\in\{d_0,d_1\}$ gives the following strict lower bounds on
area and strict upper bounds on the required right side:
\[
\begin{array}{c|cc|c}
k&\text{cyclic}&T_1\text{ right}&2\acrit+k\epsN\\ \hline
0&2.6180&2.6180&2.5808\\
1&2.6991&2.6979&2.6341\\
2&2.7830&2.7783&2.6875\\
3&2.8694&2.8682&2.7408\\
4&2.9583&2.9583&2.7942
\end{array}
\]
Thus each listed endpoint satisfies (\rif{eqn:n}).

{\bf Case 2. The triangle $T_0$ is a nonobtuse triangle, and
  $\C=\{T_0,T_1,T_1'\}$.}  
The long edges of the obtuse triangles
$T_1$ and $T_1'$ have length at least $\kappa\sqrt{8}$.

We consider a subcase where $\eta(T_1)\le1.7$ and
$\eta(T_1')\le1.7$.  Let $k_0\in\{0,1\}$ count the long external
edges of $T_0$, and let $k_T\in\{0,1,2\}$ count those of each
$T\in\{T_1,T_1'\}$.  Here long means length at least $1.72$.
Set $d_0=1.72$ if $k_0=1$, and $d_0=2\kappa$ otherwise.
For a peripheral triangle, set $d_T=1.72$ if $k_T\ge1$, and
$d_T=2\kappa$ otherwise.  Monotonicity on nonobtuse triangles gives
\[
\area(T_0) \ge \area(d_0,\kappa\sqrt{8},\kappa\sqrt{8}) >
   \begin{cases}1.73,&\text{if } k_0=0\\
     1.73+\epsN, & \text{if } k_0=1 \end{cases}.
\]
For $T\in\{T_1,T_1'\}$,
Lemma~\rif{lemma:obtuse-area-monotone} gives
\begin{equation}\libel{eqn:173}
\area(T) \ge \areta(2\kappa,d_T,1.7) >
   \begin{cases}1.08,&\text{if } k_T=0\\
     1.08+2\epsN, & \text{if } k_T\ge1 \end{cases}.
\end{equation}
If $k_T=1$, only one short side is known to be at least $1.72$;
the other is only at least $2\kappa$.  This explains the asymmetric
arguments of $\areta$.  If $k_T=2$, the same weaker estimate still
pays for both long external edges.  Thus,
$\area(T_0)>1.73+\epsN k_0$ and
$\area(T)>1.08+\epsN k_T$.  Since
$\tn(\C)=k_0+k_{T_1}+k_{T_1'}$, these bounds give
inequality (\rif{eqn:n}):
\[
\area(T_0) +\area\{T_1,T_1'\}  > 1.73 + 2(1.08) + \epsN \tn > 
3 \acrit +  \epsN \tn.
\]

By symmetry, we may now assume that $\eta(T_1) \ge 1.7$.  
The areas of $T_1$ and $T_1'$ are at least
\begin{equation}\libel{eqn:968}
\areta(2\kappa,2\kappa,2) > 0.968.
\end{equation}
Let $e$ be the shared edge of length $x$, and let $\theta_i$ be the angle of
$T_i$ opposite $e$.  We have $\theta_0\le\pi/2<\theta_1$ and
$\theta_0+\theta_1\le\pi$.  Thus,
$\sin\theta_0\le\sin\theta_1$, and the sine rule gives
\[
\eta(T_0)=\frac{x}{2\sin\theta_0}
\ge\frac{x}{2\sin\theta_1}=\eta(T_1)\ge1.7.
\]
Apply Lemma~\rif{lemma:nonobtuse-radius-area} with $a=2\kappa$
and $r_0=1.7$.  Its hypothesis $r_0\ge a/\sqrt2$ holds, so
\[
\area(T_0)\ge\kappa\sqrt{4(1.7)^2-4\kappa^2}>2.41.
\]
There are five external edges, and we have $\tn(\C)\le 5$.
This completes this case:
\[
\area(T_1) + \area(T_1') + \area(T_0) 
>
2(0.968) + 2.41 > 3\acrit + 5 \epsN \ge 3\acrit + \epsN \tn.
\] 

{\bf Case 3. The triangle $T_0$ is a nonobtuse triangle, and
  $\C=\{T_0,T_1,T_1',T_1''\}$.}

This case is almost identical to case 2.  We use the same bounds
(Equations (\rif{eqn:173}) and (\rif{eqn:968})) on $\area(T)$ as
before, for $T = T_1, T_1', T_1''$.  We can improve the bound on the
area of $T_0$:
\[
\area(T_0) 
\ge \area(\kappa\sqrt{8},\kappa\sqrt{8},\kappa\sqrt{8}) > 2.2668.
\]
Moreover, in the subcase where $\eta(T_0)\ge1.7$, retain the
nonobtuseness of $T_0$ and put $r=\kappa\sqrt8$.
Since $1.7\ge r/\sqrt2$, Lemma~\rif{lemma:nonobtuse-radius-area}
gives
\[
\area(T_0)\ge\frac r2\sqrt{4(1.7)^2-r^2}>2.6.
\]

In this case, $\tn\le 6$.  Proceeding as before, we get
\[
\area\{T_0,T_1,T_1',T_1''\} > 
\begin{cases}
2.2668 + 3 (1.08) \\
2.6 + 3(0.968)
\end{cases}
> 4\acrit +  \epsN \tn.
\] 

This completes the proof for cases involving a nonobtuse triangle
$T_0$.  In the remaining cases, we assume that $T_0$ is obtuse.  
In the remaining cases, every triangle in $\C$ is obtuse.

{\bf Case 4. The triangle $T_0$ is an obtuse triangle, and
  $\C=\{T_0,T_1\}$.}

In this case, $\tn\le 4$.  It will not be necessary to create subcases
according to whether short edges are at least $2\kappa$ or $1.72$.  We
relax the lower bound on the short edges to
$2\kappa$. %
Drop the Delaunay triangle constraint between triangles $T_0$ and $T_1$,
but keep the constraint that the nonshared vertex of $T_1$ must lie on
or outside the circumcircle of $T_0$. 

We minimize area by flattening $T_0$ by stretching its long edge until
$\eta(T_0)=2$.  We further decrease area, keeping the circumradius
fixed, by contracting the shorter edge not shared with $T_1$, until
that edge has length $2\kappa$.

Now impose a stopping condition on deformation, when $T_1$ and $T_0$ become
cocircular; in other words, when $\eta(T_1)=2$.
Next continue to minimize area by contracting an edge of $T_1$, keeping
the shared edge fixed, and
keeping its circumradius fixed, until an edge has length $2\kappa$.
Then, allowing the circumradius of $T_1$ to increase, we continue until
both shorter edges have length $2\kappa$ or until the circumradius
reaches $2$.

First assume that both shorter edges of $T_1$ have length $2\kappa$.
We have reduced to a one-parameter family of quadrilaterals.  We can
choose the parameter to be the length $x$ of the diagonal, the common
edge of $T_1$ and $T_0$.  The parameter $x$ ranges between
$\kappa\sqrt{8}$ and $x_{\max}\approx 2.9594$, determined by the
condition $\eta(2\kappa,2\kappa,x_{\max}) = 2$.  We check numerically
that
\[
\area(T_1) + \area(T_0) \ge \area(2\kappa,2\kappa,x) +
\areta(2\kappa,x,2) > 2\acrit + 4 \epsN \ge 2\acrit + \epsN \tn.
\] 

Next, assume the circumradius of $\eta(T_1)$ reaches $2$, then we
have a cocircular quadrilateral that can be treated as in Case 1.  In
particular, the minimizing cocircular quadrilateral has three edges of
length $2\kappa$ and circumradius $2$, which has area
\[
\area\{T_0,T_1\}
>  2\acrit+4\epsN \ge 2\acrit+\epsN\tn.
\]

{\bf Cases 5 and 6. The triangle $T_0$ is an obtuse triangle, and
  $\card(\C)\ge3$} %

By Remark~\rif{rem:delaunay}, there is no arrow $T \rab T_0$ in $\C$
such that the shared edge is the long edge of $T_0$.  In particular,
there cannot exist (Case 6) with $\card(\C)=4$ with %
every triangle obtuse.  Thus, Case 5 ($\card(\C)=3$) is the last case to be
considered.

We have $\tn\le 5$.
The area of $T_0$ is at least
$\areta(\kappa\sqrt{8},\kappa\sqrt{8},2) > 2.45$.  
Using our earlier estimates (\rif{eqn:968})
for $\area(T)$, for $T=T_1,T_1'$, we
have
\[
\area\{T_1,T_1',T_0\} > 2 (0.968) + 2.45 > 3\acrit + \epsN\tn.
\] 
(As mentioned earlier, each use of the function $\areta$ is justified
by a calculation based on Lemma~\rif{lemma:areta}.)
This completes the proof of the theorem.
\end{proof}

\section{Nonobtuse Clusters}\libel{sec:nonobtuse}

In this section we prove the main inequality for clusters in which
every triangle is nonobtuse.  By Corollary~\rif{lemma:n2b} and
Lemma~\rif{lemma:Nb}, if $T$ is $b$-subcritical and nonobtuse, then
$n_+(T)=n_-(T)=0$.

\begin{lemma}\libel{lemma:m2} 
  Let $T\rab T_0$. Assume that $T$ and $T_0$ are nonobtuse.  Then
  $m_-(T)=0$.
\end{lemma}

\begin{proof}  
  Assume for a contradiction that $m_-(T)>0$.  We have just observed
  that $n_+(T)=n_-(T)=0$.  By Lemma~\rif{lemma:m1m2}, we have
  $m_+(T)=0$.  By Lemma~\rif{lemma:3m2}, we have $m_-(T)\le 2$.

  We have $m_-(T)=2$.  Otherwise, if $m_-(T)=1$, we have a
  contradiction (by Corollary~\rif{lemma:m2-area}):
\[
\acrit \ge b(T) \ge \area(T) - \epsM 
> (\acrit + \epsM) - \epsM = \acrit.
\]

Because $m_-(T)=2$, there exist two pseudo-dimers $(T'_1,T'_0)$ and
$(T''_1,T''_0)$ such that $T'_0\Ra T$ and $T''_0\Ra T$.  The shared
edges $e'$ and $e''$ have length at least $1.8$.  Moreover, the angles
are large along $(T_0',e')$ and $(T_0'',e'')$, but not large along
$(T,e')$ and $(T,e'')$.  (See Definition~\rif{def:large} and
Lemma~\rif{lemma:large}.)  Lemma~\rif{lemma:2small} (below) and the
estimate
\[
\acrit \ge b(T) \ge \area(T) - 2\epsM 
> (\acrit + 2\epsM) - 2\epsM = \acrit
\]
complete the proof.
\end{proof}

\begin{lemma} \libel{lemma:2small}
  Let $T$ be a nonobtuse $P$-triangle.  Suppose that two of its edges
  $e'$ and $e''$ have length at least $1.8$ and that the angles are
  not large along $(T,e')$ and $(T,e'')$.  Then $\area(T) > \acrit +
  2\epsM$.
\end{lemma}

\begin{proof}  
  If the third edge has length at least $1.63$, then the result easily
  follows:
\[
\area(T) \ge \area(1.8,1.8,1.63) > \acrit + 2\epsM.
\]
We remark that $1.63$ is close to the minimum edge length
$2\kappa\approx 1.618$.  This leaves hardly any flexibility in the
relative position of the two pentagons along this edge.

We may assume without loss of generality that the third edge has
length in the range $[2\kappa,1.63]$.  Let the pentagons at the
vertices of $T$ be $A$, $B$, and $C$, with $d_{AB}\le 1.63$.  For the
moment, we disregard the pentagon $C$ and parallel translate $B$,
decreasing $d_{AB}$ until $A$ and $B$ come into contact.  We assume
without loss of generality that $B$ points into $A$ at $\v_B$.  Draw
the configuration as in Figure~\rif{fig:not-large} with a vertical
receptor edge $e$ on $A$.  There are two cases, depending on whether
$\v_B$ lies above or below the midpoint of the edge $e$.  (The pointer
cannot be at the midpoint of $e$ by Lemma~\rif{lemma:mid-172}, and
$1.72 > 1.63 \ge d_{AB}$.)  Let $\gamma\in[0,2\pi/5)$ be the angle
formed by edges of $A$ and $B$ at the pointer $\v_B$ as in
Figure~\rif{fig:not-large}.

We have the constraint
\[
\cos(\pi/5 - \gamma) =\sin(\gamma+3\pi/10) \le 1.63 - \kappa.
\]
This constraint expresses the fact the distance from $\c_B$ to the
edge $e$ of pentagon $A$ can be at most $1.63-\kappa$.  The constraint
implies that
\[
\gamma \ge \pi/5 +
\arccos (1.63-\kappa) > 2\pi/5 - 0.021\quad \text{ or }\quad
\gamma\le \pi/5 - \arccos(1.63-\kappa) < 0.021,
\]
according to whether $\v_B$ is below or above the midpoint of $e$.

Now we return to the original $P$-triangle with pentagons $A$, $B$,
$C$ in their original position.  Because $\gamma$ was obtained after a
parallel translation of $B$, it equals the incidence angle of lines
through edges of the original $A$ and $B$.  Changing the choice of
edges on $A$ and $B$, we find an incidence angle  $\gamma' \in
[4\pi/5-0.021,4\pi/5+0.021]$.  See Figure~\rif{fig:not-large}.

We form a triangle with angles $\alpha$, $\beta$, and $\gamma'$ by
extending edges of $A$, $B$, and $C$.  The assumption that angles are
not large along $(T,e')$ and $(T,e'')$ gives $\alpha\ge \pi/5$ and
$\beta\ge \pi/5$.  The angle sum of the triangle gives a contradiction
\[
\pi=\alpha+\beta+\gamma' 
\ge  \pi/5 + \pi/5 + (4\pi/5 - 0.021) > \pi.
\]
\end{proof}

\tikzfig{not-large}{The angle between the nearly vertical edges of
  $A$ and $B$ is $\gamma$.  On the left, $\gamma= 0.02$ and $
  2\kappa<d_{AB}\approx 1.6296 < 1.63$.  There is almost no play in
  the configuration.  The  figure on the right does not satisfy the
  constraints of the proof of Lemma \rif{lemma:2small}, but
  illustrates the notation.}
{
\begin{scope}[scale=1]
\pen{0}{0}{35.87};
\pen{1.623}{0}{-34.72+180};
\draw (0,0) node {$A$};
\draw (1.623,0) node {$B$};
\draw (0.8,1) node {$\gamma'$};
\draw ($(1.623,0)+(-34.72+180:1)$) node[anchor=north west] {\ $\v_B$};
\smalldot{$(1.623,0)+(-34.72+180:1)$};
\end{scope}
\begin{scope}[scale=1,xshift=5cm]
\pen{0}{0}{34.86};
\pen{1.756}{0}{-17.68+180};
\draw (0,0) node {$A$};
\draw (1.756,0) node {$B$};
\draw ($(1.756,0) + (-17.68+180:1)$) --  ++ 
($1.5*(-17.68+180+72:1) - 1.5*(-17.68+180:1)$) node[anchor=east] {$\gamma$};
\draw ($(1.756,0) + (-17.68+180:1)$) --  ++ 
($1.5*(34.86-72:1) - 1.5*(34.86:1)$) ;
\coordinate (A) at (34.86:1);
\coordinate (Aa) at (34.86+72:1);
\draw ($1.21*(A) - 0.21*(Aa)$) node[anchor=south] {$\gamma'$} 
-- ($-0.4*(A) + 1.4*(Aa)$);
\draw ($(1.756,0)+(-17.68+180:1)$) node[anchor=north west] {\ $\v_B$};
\smalldot{$(1.756,0)+(-17.68+180:1)$};
\end{scope}
}

Next, we turn our attention to the target $T_0$ of an arrow $T\rab T_0$,
where both $T$ and $T_0$ are nonobtuse.  By
Lemma~\rif{lemma:singleton}, the shared edge has length at least
$1.72$.

\begin{lemma}\libel{lemma:n2-0} 
  Let $T_0$ be a triangle in a cluster $\C$ containing only nonobtuse
  triangles.  Then $n_-(T_0)=0$.
\end{lemma}

\begin{proof} 
  If $n_-(T_0)>0$, then by Corollary~\rif{cor:obtuse-b-arrow} there
  exists an obtuse triangle $T'$ such that $T'\rab T_0$ and the
  cluster $\C$ contains an obtuse triangle.
\end{proof}

\begin{lemma}\libel{lemma:T1-all} 
  Let $T\rab T_0$. Assume that $T$ and $T_0$ belong to a cluster $\C$
 containing only  nonobtuse triangles.  Then
  $m_+(T)=0$.  Moreover, $n_+(T)=n_-(T)=m_+(T)=m_-(T)=0$ and
  $b(T)=\area(T)\le \acrit$.
\end{lemma}

\begin{proof}
  If $m_+(T)>0$, then there exists a pseudo-dimer $(T'_1,T'_0)$ such
  that $T=T'_0$.  Then Lemma~\rif{lemma:pd-b} implies the
  contradiction that there is no arrow $T\rab T_0$.

  The final statement is a summary of the preceding series of lemmas.
   If $n_+(T)>0$, then $(T,T_0)\in\N$ and $n_-(T_0)>0$, which is contrary
  to Lemma~\rif{lemma:n2-0}.    The equalities $n_-(T)=m_-(T)=0$ are
  Lemmas~\rif{lemma:n2b} and~\rif{lemma:m2}.
\end{proof}

\begin{lemma}\libel{lemma:4}
  Let $\C$ be a cluster consisting of nonobtuse triangles.  Assume
  that the cardinality of $\C$ is four.  Then the strict main
  inequality holds for $\C$.
\end{lemma}

\begin{proof}  
  Let $\C = \{T_0\}\cup \{T_1^i \mid i=1,2,3\}$.  We have $b(T_1^i) =
  \area(T_1^i)\le \acrit$ by Lemma~\rif{lemma:T1-all}.  By
  Lemma~\rif{lemma:n2-0}, we have
  $n_-(T_0)=0$.

  We claim $m_-(T_0)=0$.  The three distinct peripheral triangles
  occupy all three edges of $T_0$.  An incoming $\M$ pair
  $(S,T_0)$ must use one of these edges, so $S=T_1^i$ for some $i$.
  By the definition of $\M$, this $T_1^i$ is the central triangle
  of a pseudo-dimer.  Lemma~\rif{lemma:pd-b} gives
  $b(T_1^i)>\acrit$, contrary to $T_1^i\rab T_0$.

  We claim $n_+(T_0)=0$.  Otherwise, if $(T_0,T_1^i)\in\N$, then we
  get $n_-(T_1^i)>0$, which is contrary to Lemma~\rif{lemma:T1-all}.

  Hence all the negative coefficients $n_-,m_-$ are zero on the cluster:
  $b(T_1^i)=\area(T_1^i)$ and $b(T_0) \ge \area(T_0)$.  The result now
  follows from Lemma~\rif{calc:pent4}.
\end{proof}

\begin{lemma}\libel{lemma:3}
  Let $\C$ be a cluster consisting of nonobtuse triangles.  Assume
  that the cardinality of $\C$ is three.  Then the strict main
  inequality holds for $\C$.
\end{lemma}

\begin{proof}  
  Let $\C = \{T_0\}\cup \{T_1^i \mid i=1,2\}$.  We have $b(T_1^i) =
  \area(T_1^i)\le \acrit$ by Lemma~\rif{lemma:T1-all}.  By
  Lemma~\rif{lemma:n2-0}, we have
  $n_-(T_0)=0$.

  We claim that $m_-(T_0)\le1$.  Two edges of $T_0$ are occupied
  by the two peripherals.  If $m_-(T_0)\ge2$, the incoming pairs
  use two distinct edges, so at least one source is a peripheral
  $T_1^i$.  By the definition of $\M$ it is the central triangle
  of a pseudo-dimer, and Lemma~\rif{lemma:pd-b} gives
  $b(T_1^i)>\acrit$, contrary to $T_1^i\rab T_0$.

This gives
\[
b(T_0) \ge \area(T_0) - \epsM
\]
By Lemma~\rif{calc:pent3}, we have
\[
\sum_{T\in\C} b(T) \ge (\area(T_0) - \epsM) +  \area(\C\setminus \{T_0\}) 
= -\epsM+ \area(\C) > 3\acrit.
\]
This is the strict main inequality for $\C$.
\end{proof}

\begin{lemma}\libel{lemma:2} 
  Let $\C$ be a cluster consisting of nonobtuse triangles.  Assume
  that the cardinality of $\C$ is two.  If the cluster is not a dimer
  pair, then the strict main inequality holds for $\C$.
\end{lemma}

\begin{proof}  
  Let $\C = \{T_1,T_0\}$, with $T_1\rab T_0$. By
  Lemma~\rif{lemma:T1-all}, $b(T_1)=\area(T_1)\le \acrit$.

  We assume that $\C$ is not a dimer pair $(T_1,T_0)$.

  We consider the case of a pseudo-dimer.  If $(T_1,T_0)\in\PD$, then
  $m_-(T_0)=n_-(T_0)=0$ (by Lemmas~\rif{lemma:pd-m2}
  and~\rif{lemma:pd-n2}).  Thus, by Lemma~\rif{calc:pseudo-area} and
  Lemma~\rif{lemma:T1-all},
\[
b(T_1)+b(T_0) 
\ge \area\{T_1,T_0\} + \epsM (n_+(T_0) + m_+(T_0)) 
> 2\acrit - \epsM  + \epsM (n_+(T_0)+m_+(T_0)).
\]
The main inequality follows if we show that $n_+(T_0)>0$ or
$m_+(T_0)>0$.  Assume for a contradiction that $n_+(T_0)=m_+(T_0)=0$.
Pick $T$ such that $T_0\Ra T$.  The condition $n_+(T_0)=0$ implies
that $(T_0,T)\not\in\N$.  The longest edge of $T_0$ is at least $1.72$
by Lemma~\rif{calc:pseudo2}.  According to the definition of $\M$, we
have $m_+(T_0)>0$ unless uniqueness fails: there exists $T'_1\ne T_1$
such that $(T'_1,T_0)\in\PD$.  Lemma~\rif{lemma:pd-1} gives
$b(T'_1)=\area(T'_1)\le\acrit$.  Since $T'_1\Ra T_0$, it follows
that $T'_1\rab T_0$.  Thus, $T'_1$ is a third triangle of $\C$,
a contradiction.  This completes the case of a pseudo-dimer.

By Lemma~\rif{lemma:n2-0}, we have $n_-(T_0)=0$.  Thus,
by Lemma~\rif{lemma:T1-all},
\begin{equation}\libel{eqn:2}
b(T_0) \ge \area(T_0)  - \epsM m_-(T_0)
\quad b(T_1) = \area(T_1) \le \acrit.
\end{equation}

We may assume that $\area\{T_0,T_1\}> 2\acrit$.  Otherwise,
$(T_1,T_0)$ is a dimer pair or a pseudo-dimer, and these cases have
already been handled.  We assume for a contradiction that the
strict main inequality is false:
\[
2\acrit \ge b(T_1) + b(T_0).
\]

Combined with the Inequalities~\rif{eqn:2},
this gives
\begin{align*}
2\acrit &\ge b(T_1) + b(T_0) \\
&\ge \area(T_1) + \area(T_0)  - \epsM m_-(T_0)\\
&> 2\acrit  - \epsM m_-(T_0).
\end{align*}
This implies that  $m_-(T_0) > 0$.

We consider the case $m_-(T_0)= 3$.  By Lemma~\rif{lemma:3m2}, we have
$b(T_1) + b(T_0) > \ao + (\acrit + \epsN) = 2\acrit$, which completes
this case.

For any incoming $\M$ pair $(S,T_0)$, the source $S$ is the
central triangle of a pseudo-dimer.  Lemma~\rif{lemma:pd-b} gives
$b(S)>\acrit$, whereas $b(T_1)\le\acrit$.  Hence $S\ne T_1$,
and the edge shared with $S$ is not the edge shared with $T_1$.
Its length is greater than $1.8$ by Lemma~\rif{calc:pseudo2}.

We consider the case $m_-(T_0)=2$.  In this case, there are two
pseudo-dimers that share a longest edge with $T_0$.  These edges have
length at least $1.8$. The third edge is shared with $T_1$ and has
length at least $1.72$.  Then
\[
b(T_1)+b(T_0) 
\ge \ao + (\area (1.8,1.8,1.72) - 2\epsM) > 2\acrit.
\]

Finally, we consider the case $m_-(T_0)=1$.  There exists a
pseudo-dimer whose egressive edge $e$ is shared with $T_0$.
As just shown, $e$ is nonshared with $T_1$ and has length greater
than $1.8$.  The angle is large on the pseudo-dimer side by
Corollary~\rif{lemma:large}, so the complementary-angle invariant
makes it not large along $(T_0,e)$.  Together with
$T_1\Ra T_0$ and $\area(T_1)\le\acrit$, these are precisely the
hypotheses of Lemma~\rif{calc:large}.  That lemma
implies
\[
b(T_1) + b(T_0) \ge \area\{T_1,T_0\} - \epsM 
> (2\acrit  +\epsM) - \epsM = 2\acrit.
\]
This completes the proof.
\end{proof}

We are ready to give a proof of the pentagonal ice-ray conjecture. We
repeat the statement of the theorem (Theorem~\rif{thm:main}) from the
introduction of the article.

\begin{theorem}   
  No packing of congruent regular pentagons in the Euclidean plane has
  density greater than that of the pentagonal ice-ray.  The pentagonal
  ice-ray is the unique periodic packing of congruent regular
  pentagons that attains optimal density.
\end{theorem}

We combine the proof with a proof of Theorem~\rif{thm2:main}.

\begin{proof}
  By the main inequality in Lemma~\rif{lemma:main} applied to
  $a=\acrit$, together with Corollary~\rif{cor:equal} it is enough to
  give a proof of Theorem~\rif{thm2:main}.  Specifically, we show that
  every cluster in every saturated packing is finite of cardinality at
  most $4$.  This is Corollary~\rif{lemma:card4}.  If $\C$ is a dimer
  pair, then $\C$ is the ice-ray dimer in the pentagonal ice-ray and
  the (weak) main inequality holds for $\C$, with equality exactly for
  the ice-ray dimer.  This is Theorem~\rif{thm:dimer} and
  Theorem~\rif{thm:dimer-area}.  If $\C$ is not a dimer pair, then we
  show that $\C$ satisfies the strict main inequality.  If $\C$ has an
  obtuse triangle, then this is found in Theorem~\rif{thm:obtuse}.  If
  every triangle in $\C$ is nonobtuse, and if $\C$ has cardinality
  $4$, $3$, $2$, or $1$, then the result is found in
  Lemmas~\rif{lemma:4}, \rif{lemma:3}, \rif{lemma:2},
  \rif{lemma:singleton}.

  This completes the proof of the main theorem.
\end{proof}

\appendix
\section{Appendix on Explicit Coordinates}\libel{sec:coords}
\libel{sec:appendix}

\subsection{Two pentagons in contact}\libel{sec:two}

Let $A$ and $B$ be pentagons in contact, with $B$ the pointer at
vertex $\v_B$ to the receptor pentagon $A$.  Label vertices
$(\u_A,\w_A,\u_B,\v_B,\w_B)$ of $A$ and $B$ as in
Figure~\rif{fig:theta}.  Let $x=x_\alpha=\norm{\v_B}{\w_A}$ and $\alpha
= \angle(\v_B,\u_B,\u_A)$.  We have $0\le x_\alpha\le 2\sigma$ and
$0\le \alpha\le 2\pi/5$.

\tikzfig{theta}{coordinates for a pair of pentagons in contact}{
[scale=1.2]
\pen{0}{0}{0};
\pen{1.616}{0.598}{168.54};
\draw[blue] (0,0) node[anchor=south,black] {$B$}
 -- (1.616,0.598) node[anchor=south,black] {$A$};
\draw (0,0) node[anchor=north,black] {$\c_B$};
\smalldot{0,0};
\draw (1.616,0.598) node[anchor=north,black] {$\c_A$};
\smalldot{1.616,0.598};
\draw (72:1) node[anchor=south] {$\w_B$};
\draw (0:1) node[anchor=west] {$\v_B$};
\draw (-72:1) node[anchor=north] {$\u_B$};
\smalldot{72:1};
\smalldot{0:1};
\smalldot{-72:1};
\smalldot{0.636,0.797};
\draw (0.636,0.797) node[anchor=west] {$\w_A$}
node[anchor=north west] {$x$};
\smalldot{1.124,-0.2727};
\draw (1.124,-0.2727) node[anchor=north west] {$\u_A$}
node[anchor=north east] {$\alpha$};
}

Let $\bl = \bl(x_\alpha,\alpha) = \norm{\c_A}{\c_B}$, viewed as a
function of $x_\alpha$ and $\alpha$.  We omit the explicit formula for
$\bl$, but it is obtained by simple trigonometry.  Under the symmetry
$\u_A\leftrightarrow \w_A$, $\u_B\leftrightarrow\w_B$, we have the
transformation
\begin{equation}\libel{eqn:sym}
x\leftrightarrow 2\sigma-x,
\quad \alpha\leftrightarrow 2\pi/5-\alpha
\end{equation}
and
\[
\bl(x,\alpha) = \bl(2\sigma-x,2\pi/5 - \alpha).
\]
We use $(x_\alpha,\alpha)$ as the standard coordinates on the
configuration space of two pentagons in contact.  These coordinates
are particularly convenient, because they represent the configuration
space as a rectangle $[0,2\sigma]\times[0,2\pi/5]$.

\subsection{Angles of two pentagons in contact}\libel{sec:xalpha}

Let $A$ and $B$ be two pentagons in contact, given in coordinates by
$(x_\alpha,\alpha)$ as in the previous subsection.  Referring to
Figure~\rif{fig:theta}, $\v_B$ is the pointer vertex of $B$ into $A$,
and we have an angle $\theta' = \angle(\c_B,\c_A,\v_B)$ that specifies
the location of $\v_B$ relative to the segment $(\c_B,\c_A)$.  The
oriented angle $\angle(\c_B,\c_A,\v)\in[0,2\pi/5]$ to another vertex
$\v=\u_B,\w_B,\ldots$ of $B$ is $\theta' + 2\pi k/5$ for some integer
$k$.  We thus consider $\theta'$ as an angle defined modulo
$2\pi\ring{Z}/5$.  By subtracting a multiple of $2\pi/5$, we choose
the angle to lie in the range $[-\pi/5,\pi/5]$.  With these
conventions, in the figure, $\theta'$ is positive.

Referring to Figure~\rif{fig:theta}, $\u_A$ is a vertex of $A$, and we
have an angle $\theta = \angle(\c_A,\c_B,\u_A)$ that gives the
location of $\u_A$ relative to the segment $(\c_A,\c_B)$.  The angle
$\angle(\c_A,\c_B,\u)$ to another vertex of $A$ is $\theta + 2\pi k/5$
for some integer $k$.  Adjusting by a multiple of $2\pi/5$, we choose
$\theta$ to lie in the range $[0,2\pi/5]$.

We stress our convention that $\theta'$ refers to the angle on the
pointer pentagon and that $\theta$ refers to the angle on the receptor
pentagon.  It can be easily checked that 
\begin{equation}
\theta+\theta' \equiv \alpha \mod (2\pi\ring{Z}/5).
\end{equation}
(We remark in passing that we can define angles $\theta$, $\theta'$,
and $\alpha$ even when $A$ and $B$ are not in contact in such a way
that this relation still holds.)  We may consider $\theta$ and
$\theta'$ functions of the standard variables $(x_\alpha,\alpha)$.  With
our conventions $\theta'\in [-\pi/5,\pi/5]$ and $\theta\in [0,2\pi/5]$
on the range of these angles, $\theta'$ and $\theta$ are both
determined as continuous functions of $(x_\alpha,\alpha)$.  To obtain
continuity, it is necessary for $\theta$ to take both values $0$ and
$2\pi/5$, even though these values are equal modulo $2\pi\ring{Z}/5$.

In general, when two pentagons $A$ and $B$ come into contact,
sometimes $A$ is the pointer and sometimes $B$ is the pointer.  We
describe an extended coordinate system $(x_\alpha,\alpha)$, with a
domain $\alpha\in[0,4\pi/5]$ and $x_\alpha\in[0,2\sigma]$ of twice the
size that unifies both pointer directions.  See
Figure~\rif{fig:extended}.  When $\alpha\le2\pi/5$, the coordinates are
precisely as before.  Note that configurations with $\alpha=2\pi/5$
are ambiguous, with $B$ pointing into $A$ with coordinates
$(x_\alpha,\alpha)$, or with $A$ pointing into $B$ with coordinates
$(x'_\alpha,\alpha')=(2\sigma - x_\alpha,0)$.  When $\alpha > 2\pi/5$, we
let ($x_\alpha,\alpha)$ represent the configuration with $A$ pointing
into $B$ and coordinates $(x'_\alpha,\alpha') =
(2\sigma-x_\alpha,\alpha-2\pi/5)$.  The configuration of two pentagons
in contact depends continuously on the coordinates $(x_\alpha,\alpha)$.
The dependence is analytic except along $\alpha=2\pi/5$.

Using the continuous dependence of the configuration on the
coordinates, we may uniquely extend the functions $\theta$ and
$\theta'$ to continuous functions on the extended domain.  The
functions still represent the inclination angle of a vertex of $A$
(resp. $B$) with respect to the edge $(\c_A,\c_B)$.  However, when
$\alpha\ge 2\pi/5$, the range of $\theta'$ is $[0,2\pi/5]$ and the
range of $\theta$ is $[\pi/5,3\pi/5]$.  Also, $\theta$ becomes
the coordinate related to the pointer vertex (of $A$ into $B$).
Here is an explicit formula
for the extension on the domain $\alpha\ge 2\pi/5$:
\begin{align*}
\dx{AB}(x_\alpha,\alpha) &=\dx{AB}(2\sigma-x_\alpha,\alpha-2\pi/5),\\
\theta'(x_\alpha,\alpha) 
&= \theta(2\sigma-x_\alpha,\alpha-2\pi/5),\\
\theta(x_\alpha,\alpha) 
&= 2\pi/5+ \theta'(2\sigma-x_\alpha,\alpha-2\pi/5).
\end{align*}
By unifying both directions of pointing, extended coordinates lead to
a significant reduction in the number of cases to be considered.  In
fact, a single calculation can involve multi-triangle configurations
with several ($k$) edges in contact, and without extended coordinates
this leads to $2^k$ times the number of cases.

\tikzfig{extended}{Extended coordinates give a continuous transition
  from pairs $(A,B)$ with pointer $B\to A$ to pointer $A\to B$.  The
  functions $\theta,\theta'$ are continuous in $\alpha,x_\alpha$.  The
  function $\theta$ is the angle $\angle(\c_A,\c_B,\v)$, and $\theta'
  = \angle(\c_B,\c_A,\w)$. The vertices $\v$ and $\w$ transition into
  and away from the pointer vertex.  In these figures, $\alpha$ varies
  and $x_\alpha=0.8=\norm{\v}{\w}$ is fixed.}
{
\begin{scope}[scale=1]
\pen{0}{0}{41.16};
\pen{1.745}{0}{180-16.51};
\draw[blue] (0,0) node[black,anchor=east] {$A$} 
-- (1.745,0) node[black,anchor=west] {$B$};
\draw (0.8,-2) node {$\alpha = 2\pi/5 - 0.25$};
\smalldot{41.16-72:1};
\draw (41.16-72:1) node[anchor=east] {$\v~$~};
\smalldot{$(1.745,0)+(180-16.51:1)$};
\draw ($(1.745,0)+(180-16.51:1)$)  node[anchor=west] {$\w$};
\end{scope}
\begin{scope}[scale=1,xshift=4cm]
\pen{0}{0}{49.067};
\pen{1.66}{0}{180-22.93};
\draw (0.8,-2) node {$\alpha = 2\pi/5$};
\draw[blue] (0,0) node[black,anchor=east] {$A$} 
-- (1.66,0) node[black,anchor=west] {$B$};
\smalldot{49.067-72:1};
\draw (49.067-72:1) node[anchor= east] {$\v$~};
\smalldot{$(1.66,0)+(180-22.93:1)$};
\draw ($(1.66,0)+(180-22.93:1)$)  node[anchor=west] {$~\w$};
\end{scope}
\begin{scope}[scale=1,xshift=8cm]
\pen{0}{0}{55.489};
\pen{1.745}{0}{180-30.83};
\draw (0.8,-2) node {$\alpha = 2\pi/5+0.25$};
\draw[blue] (0,0) node[black,anchor=east] {$A$} 
-- (1.745,0) node[black,anchor=west] {$B$};
\smalldot{55.489-72:1};
\draw (55.489-72:1) node[anchor=east] {$\v$};
\smalldot{$(1.745,0)+(180-30.83:1)$};
\draw ($(1.745,0)+(180-30.83:1)$)  
node[anchor=west] {$~\w$};
\end{scope}
}

\subsection{Two pentagons in contact, alternative
  coordinates}\libel{sec:banana}

Inversely, $d_{AB} = \norm{\c_A}{\c_B}$, $\theta$, and $\theta'$
determine $(x_\alpha,\alpha)$.  In fact, any two of $d_{AB}$,
$\theta$, $\theta'$ determine $(x_\alpha,\alpha)$ up to finite
ambiguity.  In general, there can be two configurations of pentagons
for a given $\dx{AB}$ and $\theta'$. See Figure~\rif{fig:banana2}.
They can be distinguished by a boolean variable giving the sign of $h
:= x_\alpha-\sigma$.

\tikzfig{banana2}{With given fixed center $\c_A$ and fixed pointer
  pentagon $B$, there might be two pentagons $A$ in contact with
  $B$. They are distinguished by a boolean variable indicating whether
  the pointer $B\to A$ lies above or below the midpoint of the edge on
  $A$.}
{
\begin{scope}[scale=1,xshift=5cm]
\pen{0}{0}{29.68};
\pentpink{0}{0}{-15.887}{1}{black!50};
\pen{1.82}{0}{180-5.73};
\draw (0,0) node[black] {$A$} ;
\draw (1.82,0) node[black] {$B$};
\end{scope}
}

The interval-arithmetic computer code implements a function that generates a configuration
of two pentagons $(A,B)$ in contact as a function of $\dx{AB}$ and
$\theta'$, for contact type $B\to A$ and $h\ge0$.  (This function does
not use extended coordinates.)  Our proof of the pentagonal ice-ray
conjecture depends on having fast, numerically-stable algorithms for
computing configurations in terms of these variables on intervals.  It
is a matter of simple trigonometry to express $\theta$ in terms of
$\dx{AB}$ and $\theta'$.  However, it is somewhat more work to give
good interval bounds on $\theta$ as a function of intervals $\dx{AB}$
and $\theta'$.  The implementation of this function is based on
detailed information about the image of $(d_{AB},\theta')$ as
functions on the set of pairs $(A,B)$ with $B\to A$ and $h\ge 0$.  The
image is a convex region in the plane with a piecewise analytic
boundary.

\subsection{Pentagon existence test}\libel{sec:pet}

The values $(\dx{AB},\theta,\theta')$ can be defined for a pair of
nonoverlapping pentagons, even when $A$ and $B$ do not touch.  The
interval-arithmetic computer code implements a test to determine whether a given triple
$(d,\theta,\theta')$ is equal to a triple
$(\dx{AB},\theta,\theta')$ associated with some pair $A,B$ of
pentagons in contact.  More generally, when $(d,\theta,\theta')$
are interval-valued variables, the computer code implements a test to
determine if there exist pentagons $A$ and $B$ that do not overlap and
whose values lie in the given intervals.  The interested reader can
consult the computer code for the details of the test.

\subsection{Coordinates for zero, single, and double contact
  triangles}

The configuration space of $P$-triangles in which no pair of pentagons
is in contact is six dimensional.  Let $A$, $B$, $C$ be the pentagons
centered at the vertices $\c_A$, $\c_B$, $\c_C$ of the triangle.  For
example, we can use the six coordinates
$\dx{AB},\dx{BC},\dx{AC},\theta_{ABC},\theta_{BCA},\theta_{CAB}$.  The
three edge lengths $\dx{XY}$ of the triangle determine the triangle up
to congruence, and the three angles $\theta_{XYZ}$ determine the
inclination of each pentagon $X$ relative to the edge $(\c_X,\c_Y)$ of
the triangle.  See Figure~\rif{fig:p-triangle}.  Other quantities can
be easily computed from these six coordinates.  For example, to
compute $\theta_{ACB}$, the angle of pentagon $A$ relative to the edge
$(\c_A,\c_C)$, we use the relation
\begin{equation}
\arc_A + \theta_{ABC} + \theta_{ACB} \cong 0\mod 2\pi/5.
\end{equation}
where $\arc_A$ is the angle of the triangle at vertex $\c_A$.  
Similar equations hold for the angles at $\c_B$ and $\c_C$.

\tikzfig{p-triangle}{Angle conventions. 
  The edge lengths of the Delaunay triangle are
  $d_{AC}$, $d_{AB}$, and $d_{BC}$.  The angles at $\c_A$ are
  $\theta_{ABC}$, $\arc_A$, and $\theta_{ACB}$.  The angles at $\c_B$
  are $\theta_{BAC}$, $\arc_B$, and $\theta_{BCA}$. The angles at
  $\c_C$ are $\theta_{CBA}$, $\arc_C$, and $\theta_{CAB}$.}
{
\begin{scope}[scale=1.9,xshift=5cm]
\def\x{1.2}
\threepent
{0.00}{0.00}{96.88}
{\x*0.95}{\x*1.52}{252.34}
{\x*1.68}{0.00}{210.61};
\draw(0,0) node[anchor=east] {$A$};
\draw(\x*0.95,\x*1.52) node[anchor=south] {$B$};
\draw(\x*1.68,0) node[anchor=west] {$C$};
\smalldot{0,0};
\smalldot{\x*0.95,\x*1.52};
\smalldot{\x*1.68,0};
\draw(0,0) -- ++ (96.88:1);
\draw(0,0) -- ++ (96.88- 2*72:1);
\draw(\x*0.95,\x*1.52) -- ++ (252.34+72:1);
\draw(\x*0.95,\x*1.52) -- ++ (252.34-72:1);
\draw(\x*1.68,0) -- ++ (210.61 - 2*72:1);
\draw(\x*1.68,0) -- ++ (210.61:1);
\draw(25:0.65) node {$\arc_A$};
\draw(75:0.7) node {$\theta_{ABC}$};
\draw(-20:0.6) node {$\theta_{ACB}$};
\draw($(\x*0.95,\x*1.52)+(-45:0.7)$) node {$\theta_{BCA}$};
\draw($(\x*0.95,\x*1.52)+(-95:0.65)$) node {$\arc_B$};
\draw($(\x*0.95,\x*1.52)+(210:0.55)$) node {$\theta_{BAC}$};
\draw($(\x*1.68,0)+(90:0.65)$) node {$\theta_{CBA}$};
\draw($(\x*1.68,0)+(140:0.60)$) node {$\arc_C$};
\draw($(\x*1.68,0)+(195:0.6)$) node {$\theta_{CAB}$};
\end{scope}
}

Each $P$-triangle determines a six-tuple $(\dx{AB},\dx{BC},\ldots)$.
We can algorithmically test whether a six-tuple of real numbers (or of
intervals) comes from $P$-triangles by using the triangle inequality
and the test of Section~\rif{sec:pet}.

The configuration space of $P$-triangles in which a single pair
$(A,B)$ of pentagons comes into contact is five-dimensional.  Let $A$,
$B$, $C$ be the pentagons centered at the vertices $\c_A$, $\c_B$,
$\c_C$ of the triangle.  Assume that $A$ is in contact with $C$, with
pointer $A\to C$.  For example, we can use the five coordinates
$\dx{AB},\dx{BC},\dx{AC},\theta_{ACB},\theta_{BAC}$.  The three edge
lengths $\dx{XY}$ of the triangle determine the triangle up to
congruence, and the two angles $\theta_{XYZ}$ determine the
inclination of each pentagon relative to the triangle.  It is clear
that the $\theta_{ACB}$ fixes the inclination of $A$ and that
$\theta_{BAC}$ fixes the inclination of $B$.  The variable
$\theta_{CAB}$ is computed from $\dx{AC}$ and $\theta_{ACB}$ by the
procedure in Section~\rif{sec:banana}.  Other quantities can be easily
computed from these quatitites. Again, we can algorithmically
test whether a given five-tuple $(\dx{AB},\dx{BC},\ldots)$ 
comes from a $P$-triangle.

The configuration space of $P$-triangles in which two separate pairs
$(A,B)$ and $(B,C)$ of pentagons comes into contact is
four-dimensional.  Let $A$, $B$, $C$ be the pentagons centered at the
vertices $\c_A$, $\c_B$, $\c_C$ of the triangle.  We can use the
extended coordinates $(x_\alpha,\alpha)$ of Section~\rif{sec:xalpha}
to give the relative position of $A$ and $B$ and similar coordinates
$(x_\gamma,\gamma)$ to give the relative position of $B$ and $C$.
These four coordinates uniquely determine the $P$-triangle up to
finite ambiguity.  The ambiguity is resolved by specifying which edge
of $B$ is in contact with $A$ and which edge is in contact with $C$.
These four coordinates determine other quantities such as lengths
$\dx{AB}$ and $\dx{BC}$ and angles $\theta_{BAC}$, $\theta_{ABC}$,
$\theta_{BCA}$, $\theta_{CBA}$.  Once again, we can algorithmically
test whether a given four-tuple $(x_\alpha,\alpha,x_\gamma,\gamma)$
comes from a $P$-triangle.

\subsection{Triple contact}

If we have coordinates on a $3C$-triangle that determine the variables
$(x_\alpha,\alpha)$ for each of the pairs $\{A,B\}$, $\{B,C\}$, and
$\{A,C\}$ of pentagons, then we may use the function $\bl$ of
Section~\rif{sec:two} to calculate the edge lengths and area of the
$3C$-triangle.

The configuration space of $3C$-triangles is three dimensional,
obtained by imposing three contact constraints between pairs of
pentagons on the six-dimensional configuration space of all
$P$-triangles.

% COMPUTATION SMPXAOD | sec:shared | Shared-edge chart restrictions
\subsection{$3C$ triangles with a shared edge}\libel{sec:shared}

Some calculations deal with a dimer pair of $3C$ triangles sharing a
triangle edge and two pentagons $\bar{A}$ and $\bar{C}$, say with
$\bar{A}$ pointing into $\bar{C}$.  (We place bar accents on symbols
in this subsection for compatibility with the sections that follow.)
Let $\bar{B}$ and $\bar{D}$ be the two outer pentagons of the dimer
pair.  The configuration space of dimers in which both triangles have
triple contact is four dimensional.  In this situation, it is best to
develop coordinate systems that make efficient use of the shared
information.  Associated with the pentagons $(\bar{A},\bar{C})$ in
contact are two variables $(\bar x_\beta,\bar\beta)$ (that we rename
from $(x_\alpha,\alpha)$ in Section \rif{sec:xalpha}).  It is generally
advantageous to make $(\bar x_\beta,\bar\beta)$ two of the four
coordinates on the dimer pair.  We supplement this with one further
angle $\bar\alpha$ to determine the position of $\bar B$ and a further
angle to determine $\bar D$.

We focus on the $P$-triangle $(\bar A,\bar B,\bar C)$, with
coordinates $(\bar\alpha,\bar\beta,\bar x_\beta)$.  Similar
considerations apply to the other $P$-triangle $(\bar A,\bar D,\bar
C)$.  We assume that $\bar A$ points into $\bar C$.

For each triple contact, the following sections give further details
about its coordinate systems.  There will be a {\it shared edge}
coordinate system for each contact type and each pair of pentagons
$(\bar A,\bar C)$ such that $\bar A$ points into $\bar C$.

We call a $\Gamma$-triangle to be a $P$-triangle that appears as one
of the two triangles (with given shared edge) in the curve $\Gamma$ of
Section~\rif{sec:gamma}.  We set up coordinates in a uniform manner so
that regardless of the contact type of the $3C$ triangle, the curve
$\Gamma$ is parameterized by variable $t$ and is given as
\[
t = \bar x_\beta  -\sigma, \quad
\bar\alpha = 0,\quad
\bar\beta = \pi/5.
\]
(This is the curve restricted to the $P$-triangle $(\bar A,\bar B,\bar
C)$, with a similar description on the other triangle $(\bar A,\bar
D,\bar C)$ in the dimer.)

When the configuration space of a given $3C$-type contains
$\Gamma$-triangles and the ice-ray triangle, we also describe a path
$P$ with parameter $s$ from an arbitrary triangle in that
configuration space to a $\Gamma$-triangle.  Again, we set up
coordinates uniformly so that the formulas are independent of the
contact type of the $P$-triangle.  Coordinates will be defined in such
a way that for all points of the domain, the relation $\bar\alpha\ge
0$ holds.  The path $P$ from an arbitrary point in the domain
$(\bar\alpha_0,\bar\beta_0,\bar x_{\beta,0})$ to the curve $\Gamma$ is
defined by functions $s\mapsto (\bar\alpha(s),\bar\beta(s),\bar
x_\beta(s))$ (with parameter $s\ge 0$), where
\begin{align*}
\bar\alpha(s)&=\max\{\bar\alpha_0-s,0\},\\
\bar\beta(s)&=\frac\pi5+\op{sgn}(\bar\beta_0-\pi/5)
                  \max\{|\bar\beta_0-\pi/5|-s,0\},\\ %
\bar x_\beta(s)&=\bar x_{\beta,0}.
\end{align*}
We take $\op{sgn}(0)=0$.  Thus, a coordinate stops when it reaches
its target, while the other coordinates continue.  In particular,
$\bar x_\beta$ remains constant.  The path is continuous and
piecewise linear, and terminates when $\bar\alpha=0$ and
$\bar\beta=\pi/5$.  These
termination conditions are the defining conditions of the image of
$\Gamma$.  Thus, the path leads to $\Gamma$ in all cases.

If the initial configuration satisfies
\[
|\bar\alpha_0|\le M,\quad 
|\bar\beta_0-\pi/5|\le M,\quad |\bar x_{\beta,0}-\sigma|\le M,
\]
then these inequalities for $(\bar\alpha(s),\bar\beta(s),\bar
x_\beta(s))$ hold along the path $P$.

We write $\bar\alpha(s,\bar\alpha_0)$, $\bar\beta(s,\bar\beta_0)$, and
$\bar x_\beta(s,\bar x_{\beta,0})$ to show the dependence of the path
$P$ on the initial point.  The dimer pair has a second triangle $(\bar
A,\bar D,\bar C)$ whose initial configuration has coordinates
$(\bar\alpha_1,\bar\beta_1,\bar x_{\beta,1})$, where by Equation
(\rif{eqn:sym}) we have
\[
\bar\beta_0+\bar\beta_1 = 2\pi/5,\quad 
\bar x_{\beta,0}+\bar x_{\beta,1} = 2\sigma.
\]
The equations for the path on the two triangles then satisfy
\[
\bar\beta(s,\bar\beta_0)+\bar\beta(s,\bar\beta_1)=2\pi/5,\quad
\bar x_{\beta}(s,\bar x_{\beta,0}) + \bar x_{\beta}(s,\bar x_{\beta,1}) = 2\sigma.
\]
Indeed, $|\pi/5-\bar\beta_0|=|\pi/5-\bar\beta_1|$, and
$\bar\beta(s,\bar\beta_0)+\bar\beta(s,\bar\beta_1)$ is constant
in $s\ge0$, %
and both coordinates stop at the same time.
The shared
length coordinates remain constant.  By the
two-pentagon contact coordinates, these two identities describe the
same shared pair of pentagons from either triangle.  Thus, the paths
are coherent along $(\bar A,\bar C)$ and determine a path of the full
dimer $(\bar A,\bar B,\bar C,\bar D)$ that terminates at the curve
$\Gamma$.

\subsection{Triple contact at a $\Delta$-junction}

We describe a coordinate system on $3C$-triangles of $\Delta$-junction
type.  As indicated in Figure~\rif{fig:cord-delta}, we use coordinates
$(x_\alpha,\alpha,\beta)$, where $x_\alpha$ is a length and $\alpha$
and $\beta$ are each angles between lines through edges of pentagons
in contact.  We assume that $B$ points into $A$ and into $C$ and that
$A$ points into $C$. The length $x_\alpha$ is the (small) distance
between the nearly coincident vertices of pentagons $B$ and $C$.  The
coordinates satisfy the conditions $0\le\beta\le\alpha\le\pi/5$,
$\alpha+\beta\le \pi/5$, and $x_\alpha\in[0, 2\sigma -
\sigma/\kappa]$.  Starting from these coordinates, we define $\gamma$
by $\alpha+\beta+\gamma=\pi/5$, and angles $\alpha',\beta',\gamma'$ of
the inner triangle $\Delta$ by
\begin{align}\libel{eqn:abc}
\alpha+\alpha' &= 2\pi/5,\\
\beta+\beta' &= 2\pi/5,\nonumber\\
\gamma+\gamma' &= 2\pi/5.\nonumber
\end{align}
The edges $y_\alpha$, $y_\beta = 2\sigma$, and $y_\gamma$ of the
triangle $\Delta$ opposite the angles $\alpha'$, $\beta'$, $\gamma'$,
respectively are easily computed by the law of sines.  Define
$x_\beta$ by $x_\alpha+y_\gamma+x_\beta=2\sigma$, and $x_\gamma$ by
$y_\alpha+x_\gamma=2\sigma$.  The value $x_\beta$ is the distance
between the nearly coincident vertices of pentagons $A$ and $C$, and
$x_\gamma$ is the distance between the nearly coincident vertices of
pentagons $A$ and $B$.  The edges of the $3C$ Delaunay triangle have
lengths
\[
\bl(x_\alpha,\alpha'),\quad \bl(x_\beta,\beta'),\quad \bl(x_\gamma,\gamma').
\]
By Lemma~\rif{lemma:delta}, triangles of type $\Delta$ have area too
large to be relevant for calculations of dimers.  There is no need to
describe the shared coordinates $(\bar \alpha,\bar\beta,\bar x_\beta)$
in this case.

\tikzfig{cord-delta}{Coordinates for $\Delta$-types}
{
[scale=1.0]
\threepentnoD
{0.00}{0.00}{-5.16}%
{0.99}{1.70}{235.38}%
{1.98}{0.00}{183.43}; %
\draw (0,0) node {$C$};
\draw (0.99,1.70) node {$B$};
\draw (1.98,0.0) node {$A$};
\draw (1.0,0.6) node {$\Delta$};
\draw (-5.16:1) -- ++ (126 - 5.16 :2.5) node[anchor=south] {$\alpha$};
\draw (-5.16:1) -- ++ (- 180 + 126 - 5.16 :1.5) node[anchor=west] {$\beta$};
\draw ++ (1.98,0) ++ (183.43 - 72:1) 
-- ++ (54 + 3.43:1) node[anchor=south] {$\gamma$};
\draw (72-5.16:1) -- ++ (126 + 90 - 5.16:0.5);
\draw (0.99,1.70) ++ (235.38:1) node[anchor = south west] {$x_\alpha$} 
-- ++ (126+90 - 5.16:0.5);
}

\subsection{Triple contact at a pinwheel type}\libel{sec:pinwheel}

We describe a coordinate system on $3C$-triangles of pinwheel type.
We assume that $C$ points into $B$, that $B$ points into $A$, and that
$A$ points into $C$.  As indicated in Figure~\rif{fig:cord-pinwheel},
we use coordinates $(\alpha,\beta,x_\gamma)$.  The angles $\alpha$ and
$\beta$ are angles between pentagon edges on touching pentagons.  The
value $x_\gamma$ is the distance between the pointer vertex of
pentagon $A$ and the pointer vertex of pentagon $B$.  The coordinates
satisfy constraints: $0\le\alpha$, $0\le\beta$, $\alpha+\beta\le
\pi/5$, and $0\le x_\gamma\le 2\sigma$.  Define $\gamma$ by
$\alpha+\beta+\gamma=\pi/5$.  The angles $\alpha'$, $\beta'$, and
$\gamma'$ of the inner background triangle $P$ of the pinwheel are
given by Equation~\rif{eqn:abc}.  The edge lengths $x_\alpha$,
$x_\beta$, $x_\gamma$ of the inner triangle $P$ are easily computed
from $(\alpha,\beta,x_\gamma)$ by the law of sines.  The edges of the
$3C$-triangle have lengths
\[
\bl(x_\alpha,\alpha),\quad 
\bl(x_\beta,\beta),\quad \bl(x_\gamma,\gamma).
\]

\tikzfig{cord-pinwheel}{Coordinates for pinwheel type}{
\begin{scope}[xshift=4cm,scale=1.2]
\threepentnoD{0.00}{0.00}{46.69}%
{0.82}{1.53}{218.09}%
{1.73}{0.00}{163.28};%
\draw(0,0) node {$A$};
\draw(0.82,1.53) node {$C$};
\draw(1.73,0) node {$B$};
\draw(1.8,1.2) node {$\alpha$};
\draw(-0.2,1.0) node {$\beta$};
\draw(0.85,-0.7) node {$\gamma$};
\draw(0.9,0.5) node {$P$};
\draw(0.5,0.5) node {$x_\gamma$};
\smalldot{0.772,0.288}; %
\smalldot{0.686,0.728}; %
\end{scope}
}

\subsubsection{Pinwheel type with a shared edge}

Now we specialize this discussion of Section \rif{sec:shared} to
pinwheels.  Pinwheels have a rotational symmetry, so we may assume
without loss of generality that the shared edge is $A\to C$ (meaning,
$A$ and $C$ are the shared pentagons and $A$ points to $C$).  This
choice of shared edge gives $(\bar A,\bar C) = (A,C)$.  The shared
variables are $(\bar x_\beta,\bar\beta)=(x_\beta,\beta)$.  The
nonshared variable is $\bar\alpha=\alpha$.

\subsection{Triple contact at a $LJ$-junction type}

We describe a coordinate system on $3C$-triangles of $LJ$-junction
type.  As indicated in Figure~\rif{fig:cord-L}, we use coordinates
$(\alpha,\beta,x_\alpha)$.  The angles $\alpha$ and $\beta$ are each
formed by edges of two pentagons in contact.  Let $x_\alpha$ be the
distance between the pointer vertex of $C$ to $A$ and the pointer
vertex of $B$ to $C$.  The coordinates satisfy relations:
$\alpha,\beta\in [0,2\pi/5]$, $\pi/5\le\alpha+\beta\le 3\pi/5$, and
$0\le x_\alpha\le 2\sigma$.  Define $\gamma$ by
$\alpha+\beta+\gamma=3\pi/5$.  The three acute angles $\alpha'$,
$\beta'$, and $\gamma'$ of the inner $L$-shaped quadrilateral are
given by Equation~\rif{eqn:abc}.  The edge lengths of the $L$-shaped
quadrilateral are easily computed by triangulating the quadrilateral
into two triangles and applying the law of sines.  (Triangulate $L$ by
extending the line through the edge of $A$ containing the pointer
vertex of $C$ into $A$.)  This gives $x_\beta$, the distance between
the pointer vertex of $C$ to $A$ and the inner vertex of $A$.  This
gives $x_\gamma$, the distance between the pointer vertex of $B$ and
the inner vertex of pentagon $A$.  As before, the edges of the
$3C$-triangle have lengths
\[
\bl(x_\alpha,\alpha),\quad \bl(x_\beta,\beta),\quad \bl(x_\gamma,\gamma).
\]

\tikzfig{cord-L}{Coordinates for $LJ$-junction type}{
\begin{scope}[xshift=8cm,scale=1.2]
\threepentnoD{0.00}{0.00}{80.86}
{0.97}{1.58}{232.21}
{1.82}{0.00}{223.24};
\draw (0,0) node {$C$};
\draw (0.97,1.58) node {$B$};
\draw (1.82,0) node {$A$};
\draw (2.0,1.2) node {$\gamma$};
\draw (-0.1,1.2) node {$\alpha$};
\draw (0.85,-0.8) node {$\beta$};
\draw (0.85,0.75) node {$L$};
\draw (0.55,0.4) node {$x_\alpha$};
\smalldot{$(0,0)+(80.86-72:1)$};
\smalldot{$(0.97,1.58)+(232.24:1)$};
\end{scope}
}

\subsubsection{$LJ_1$-junction type}

Now we specialize to $LJ$-junctions with a shared edge $C\to A$.
Then $(\bar A,\bar B,\bar C)=(C,B,A)$.

The shared variables are $(\bar \beta,\bar x_\beta)=(\beta,x_\beta)$.
The nonshared variable is $\bar\alpha=\gamma' = 2\pi/5-\gamma\ge0$.
These coordinates are numerically stable.  We compute other angles and
edges by triangulating the $L$ by extending the receptor edge of $A$
and using the law of sines.

\subsubsection{$LJ_2$-junction type}

We specialize to $LJ$-junctions with a shared edge $B\to C$.
In this case, $(\bar A,\bar B,\bar C)=(B,A,C)$.

The shared variables are $(\bar x_\beta,\bar\beta)=(x_\alpha,\alpha)$.
The nonshared variable is $\bar\alpha=\beta$.  These coordinates are
numerically stable.  They are exactly the standard variables given
above for a general $LJ$-junction.

\subsubsection{$LJ_3$-junction type}\libel{sec:one-ljx}

We specialize to $LJ$-junctions with a shared edge $B\to A$.
In this case, $(\bar A,\bar B,\bar C)=(B,C,A)$.

The shared variables are $(\bar x_\beta,\bar\beta)=(x_\gamma,\gamma)$.
The nonshared variable is $\beta$.  If $\beta > 0.9$ a
calculation shows\calcref{one\_ljx} that the triangle area is greater
than $\acrit+\epsN$. Thus, in treating triangles of smaller area than this,
we may assume
$\beta\le0.9$.  
Under the additional assumption $\beta\le0.9$, these coordinates
are numerically stable.  We
compute other lengths and angles by triangulating the $L$-region by
extending the receptor edge of $A$.

There does not exist a dimer such that one of its triangles
is a $LJ_3$ configuration with shared edge $B\to A$.

\subsection{Triple contact at a $TJ$-junction type}

We describe a coordinate system on $3C$-triangles of $TJ$-junction
type.  As indicated in Figure~\rif{fig:cord-T}, we use coordinates
$(\alpha,\beta,x_\gamma)$.  The angles $\alpha$ and $\beta$ are each
formed by edges of two pentagons in contact.  The length $x_\gamma$ is
the distance between the pointer vertex of $B$ into $A$ and the inner
vertex of $A$.  The coordinates satisfy:
$\alpha,\beta\in[\pi/5,2\pi/5]$, $3\pi/5\le \alpha+\beta\le 4\pi/5$,
$0\le x_\gamma\le 2\sigma$.  Define $\gamma$ by
$\alpha+\beta+\gamma=\pi$.  Three of the angles $\alpha'$, $\beta'$,
and $\gamma'$ of the inner irregular $TJ$-shaped pentagon $P$ are
given by Equation~\rif{eqn:abc}.  The edge lengths of the $TJ$-shaped
pentagon are easily computed by triangulating $P$ into three triangles
and applying the law of sines.  (Triangulate by extending the edge of
$P$ shared with $A$ that ends at the pointer vertex of $A$ into $C$
and by extending the edge of $P$ shared with $C$ that contains pointer
vertex of $B$ into $C$.)  This gives $x_\alpha$, the distance between
the pointer vertex of $B$ to $C$ and the inner vertex of $C$.  This
gives $x_\beta$, the distance between the pointer vertex of $A$ to $C$
and the inner vertex of $C$.  As before, the edges of the
$3C$-triangle have lengths
\[
\bl(x_\alpha,\alpha),\quad \bl(x_\beta,\beta),\quad \bl(x_\gamma,\gamma).
\]

\tikzfig{cord-T}{Coordinates for $TJ$-junction types}{
\begin{scope}[xshift=12cm,scale=1.2]
\threepentnoD{0.00}{0.00}{114.48} %
{0.90}{1.59}{237.18} %
{1.66}{0.00}{219.24}; %
\draw (0,0) node {$A$};
\draw (0.9,1.59) node {$B$};
\draw (1.66,0) node {$C$};
\draw (0.9,0.9) node {$P$};
\draw (0,1.1) node {$\gamma$};
\draw (0.75,-0.9) node {$\beta$};
\draw (1.9,1.2) node {$\alpha$};
\draw (0.45,0.5) node {$x_\gamma$};
\smalldot{0.737,0.675};
\smalldot{0.358,0.75};
\end{scope}
}

\subsubsection{$TJ_1$-junction type}\libel{sec:one-tjx}

We specialize to $TJ$-junctions with a shared edge $A\to C$.  In this
case, $(\bar A,\bar B,\bar C)=(A,B,C)$.

The shared variables are $(\bar x_\beta,\bar\beta)=(x_\beta,\beta)$.
The nonshared variable is $\bar\alpha=\alpha$.  If $\beta < 1.0$ a
calculation %
shows\calcref{one\_tjx} that the triangle area is greater than $\acrit+\epsN$.
Thus, for a triangle of area at most $\acrit+\epsN$, we may assume $\beta\ge1.0$.
Under this additional
assumption, these coordinates are numerically stable.
We compute other lengths and angles by triangulating the $TJ$-region
by extending the receptor edge of $A$ and extending the receptor edge
of $C$.

There does not exist a dimer such that one of its triangles
is a $TJ_1$ configuration with shared edge $A\to C$.

\subsubsection{$TJ_2$-junction type}

We specialize to $TJ$-junctions with a shared edge $B\to A$. In this
case, $(\bar A,\bar B,\bar C)=(B,C,A)$.

The shared variables are $(\bar x_\beta,\bar\beta)=(x_\gamma,\gamma)$.
The nonshared variable is $\bar\alpha =\beta$.  From
$\alpha+\beta+\gamma=\pi$, we obtain $(\alpha,\beta,x_\gamma)$, which
are the standard coordinates described above for $TJ$-junction
triangles.  These coordinates are numerically stable.

There does not exist a dimer such that one of its triangles
is a $TJ_2$ configuration with shared edge $B\to A$.

\subsubsection{$TJ_3$-junction type}

We specialize to $TJ$-junctions with a shared edge $B\to C$.
In this case, $(\bar A,\bar B,\bar C)=(B,A,C)$.

The shared variables are $(\bar x_\beta,\bar\beta)=(x_\alpha,\alpha)$.
The nonshared variable is $\bar\alpha=\beta'=2\pi/5-\beta$.  These
coordinates are numerically stable.  We compute other lengths and
angles by triangulating the $TJ$-region by extending the receptor edge
of $A$ (of $B\to A$) and extending the receptor edge of $C$ (of $A\to
C$).

\subsection{Triple contact at a pin-$T$ junction type}\libel{sec:pint}
We describe a coordinate system on $3C$-triangles of pin-$T$ junction
type.  As indicated in Figure~\rif{fig:cord-pint}, we use coordinates
$\alpha$, $\beta$, and $x_\alpha$.  The angles $\alpha$ and $\beta$
are each formed by edges of two pentagons in contact.  The length
$x_\alpha$ is the distance between the nearly coincident vertices of
$B$ and $C$.  The coordinates satisfy $\pi/5 \le \alpha \le 2\pi/5$,
$\pi/5\le \beta \le 2\pi/5$, and $3\pi/5 \le \alpha+\beta$.
Lemma~\rif{lemma:0605} shows that $0\le x_\alpha\le 0.0605$.  Define
$\gamma$ by $\alpha+\beta+\gamma = \pi$.  Three of the angles
$\alpha'$, $\beta'$, and $\gamma'$ of the inner irregular $T$-shaped
pentagon are given by Equation~\rif{eqn:abc}.  The edge lengths of the
$T$-shaped pentagon $P$ are easily computed by triangulating $P$ into
three triangles and applying the law of sines.  (Triangulate by
extending the two edges of $P$ that meet at the pointer vertex of $C$
into $A$.)  This gives $x_\beta$, the distance between the pointer
vertex of $C$ to $A$ and the inner vertex of $A$.  This gives
$x_\gamma$, the distance between the pointer vertex of $A$ to $B$ and
the pointer vertex of $B$ into $C$.  As before, the edges of the
$3C$-triangle have lengths
\[
\bl(x_\alpha,\alpha),\quad \bl(x_\beta,\beta),\quad \bl(x_\gamma,\gamma).
\]

\tikzfig{cord-pint}{Coordinates for pin-$T$ junction types. Although
  it is difficult to tell from the figure, $A$ points into $B$, $B$
  into $C$, and $C$ into $A$.  The parameter $\beta'\approx 0$
  measures the incidence angle between the nearly horizontal edges of
  $A$ and $C$.  The region bounded by the three pentagons is a
  $T$-shaped pentagon, with stem between the nearly parallel edges of
  $A$ and $C$ and two arms along $B$.  The arm between $B$ and $C$ is
  almost imperceptible.  The distance between the neighboring vertices
  of $B$ and $C$ is $x_\alpha$.  The figure on the right distorts the
  pentagons to make the incidence relations more apparent.}
{
\begin{scope}[xshift=4.5cm,yshift=1.5cm]
\threepentnoD{0.00}{0.00}{90} %
{1.40}{-1.377}{0} %
{-0.35}{-1.628}{-18.89}; %
\draw (0,0) node {$C$};
\draw (1.4,-1.37) node {$B$};
\draw (-0.35,-1.62) node {$A$};
\draw (0.36,-2.4) node {$\gamma$};
\draw (-1.1,-0.5) node {$\beta$};
\draw (1.0,-0.3) node {$\alpha$};
\draw (0.4,-0.49) node {$x_\alpha$};
\draw (1.4 - 0.809,-1.377 - 0.5878) -- +(0,-0.75);
\end{scope}
\begin{scope}[xshift=9.5cm,yshift=1cm]
\draw[red] (0.2,-1.2) -- ++ (90:1.4) -- ++ (10:1); %
\draw[red] (0.2,-1.2) -- ++ (-10:1);
\draw[red] (-1,0) -- ++(1,0) -- ++ (45:1); %
\draw[red] (-1,0) -- ++(100:1);
\draw[red] (-1,0) -- + (-10:0.9) -- (0.2,-0.8) -- ++ (-120:1); %
\draw[red] (-1,0) -- ++(180-10:0.3) -- ++(-110:1);
\draw (0.9,-0.4) node {$B$};
\draw (-0.5,0.7) node {$C$};
\draw (-0.7,-0.7) node {$A$};
\end{scope}
}

\begin{lemma}\libel{lemma:0605}
  Let $T$ be a $3C$ triangle of type pin-$T$.  The coordinates
  $\alpha$, $\beta$, and $x_\alpha$ satisfy the relation
\[
x_\alpha \sin(2\pi/5) \le 2\sigma (\sin(\alpha+\pi/5) - \sin(\beta+\pi/5)).
\]
In particular, $x_\alpha \le 0.0605$.
\end{lemma}

\begin{proof} Let $\v_{AB}$ be the pointer vertex of $A$ to $B$, and
  let $\v_{BC}$ be the pointer vertex of $B$ to $C$.  Let $\v$ and
  $\v_{BC}$ be the endpoints of the edge of $B$ containing $\v_{AB}$.
  We represent $T$ as in Figure~\rif{fig:cord-pint}, with the lower edge of
  $C$ along the $x$-axis.  Because $\v_{AB}$ lies on the segment between
  $\v$ and $\v_{BC}$, the $y$-coordinate $y(\v_{AB})$ of $\v_{AB}$ is
  nonpositive and lies between the $y$-coordinates $y(\v)$ and
  $y(\v_{BC})$.  We have
\begin{align*}
y(\v) &= x_\alpha \sin(2\pi/5) - 2\sigma\sin(\alpha+\pi/5)\\
y(\v_{AB}) & = -x_\beta \sin(\beta') - 2\sigma\sin(\beta+\pi/5).
\end{align*}
Using $x_\beta \sin(\beta')\ge 0$ and $y(\v) \le y(\v_{AB})$, we
obtain the claimed inequality.

Recall that $\pi/5\le \beta \le 2\pi/5$.  In particular, we have
$\sin(\alpha+\pi/5) \le 1$ and $\sin(\beta+\pi/5)\ge \sin(2\pi/5)$.
This gives
\[
x_\alpha \le 2\sigma(1/\sin(2\pi/5) - 1) < 0.0605.
\]
\end{proof}

\subsubsection{Pin-$T$-junction type}\libel{sec:one-pintx}

A computer interval-arithmetic calculation
shows\calcref{one\_pintx} that the longest edge in a pin-$T$ junction is always $B\to C$.  We specialize to
pin-$T$ junctions with a shared edge $B\to C$.  We use the standard
coordinates, with $(\bar x_\beta,\bar\beta)=(x_\alpha,\alpha)$ shared.
An ice-ray triangle does not have type pin-$T$.

This completes our discussion of coordinates used for computations.

\section{Appendix on Computer Calculations}\libel{appendix:cc}

The code for the computer interval-arithmetic assisted proofs is written in Objective
Caml.  There are about five thousand lines of code, available for
download from github \cite{Git}.  The 2016 computer calculations for this
article took about 60 hours in total to run on an Intel quad 2.6 GHz
processor with 3.7GB memory, running the Ubuntu operating system. The
hashtables occupied between 1 and 2 GB of memory.

Mathematica was also used for secondary tasks such as computing numerical
constants and generating figures for this article.

\subsection{Interval arithmetic}

The computations use interval arithmetic with outward bounds on
floating-point operations. The accompanying code release uses the
CRlibm-backed OCaml interval implementation, descended from the
package of Alliot and Gotteland \cite{All}, with a small documented
compatibility adapter. Its source, dependencies, runtime settings,
tests, and execution records are identified in the accompanying
repository. Intervals are stored by their lower and upper endpoints.
A domain or numerical-instability failure is not a successful
exclusion; it must be handled by justified subdivision or reported
as an incomplete calculation.

\subsection{Automatic differentiation}\libel{sec:auto-diff}

Recall that there are several ways to compute derivatives by computer,
such as numerical approximation by a difference quotient
$(f(y)-f(x))/(y-x)$, symbolic differentiation (as in computer algebra
systems), and automatic differentiation.  In this project, we
differentiate functions of a single variable using automatic
differentiation.  The value of a function and its derivative are
represented as a pair $(f,f')$ of intervals (the $1$-jet of the
function at an interval-valued point $x_0$), where $f$ is an interval
bound on the function at $x_0$, and $f'$ is an interval bound on the
derivative of the function at $x_0$.  More complex expressions can be
built from simpler expressions by extending arithmetic operations
$(+)$, $(-)$, $(*)$, $(/)$ to $1$-jets.  For example,
\begin{align*}
(f,f') + (g,g') &= (f+g,f'+g'),\\
(f,f') * (g,g') &= (f g,f g' + f' g),\\
(f,f') / (g,g') &= (f / g, (f' g - f g')/g^2).
\end{align*}
where the component-wise arithmetic operations on the right are
computed by interval arithmetic.  Automatic differentiation extends
standard functions $F$ to functions $F^D$ on $1$-jets.
For example,
\[
\op{sqrt}^D(f,f') = (\sqrt{f}, f'/(2\sqrt{f})),\quad 
\sin^D (f,f') = (\sin(f), \cos(f) f').
\]

Sections~\rif{sec:dimer-pair} and Section~\rif{sec:coords} describe
the proof of the local minimality of the ice-ray dimer.  We review
that argument here with an emphasis on automatic differentiation.
Automatic differentiation allows us to show that the ice-ray dimer is
the unique minimizer of area in an explicit neighborhood of the
ice-ray dimer.  In Section~\rif{sec:dimer-pair}, we give a curve
$\Gamma$ (in the configuration space of dimers) with parameter
$t\in\ring{R}$ that passes through the ice-ray dimer point at $t=0$.

For any point $\mathbf{x}_0$ in an explicit neighborhood of the ice-ray dimer
in the dimer configuration space, that section describes a path from
$\mathbf{x}_0$ to a point on the curve $\Gamma$.  Using automatic
differentiation algorithms, we show by computer that area decreases as
we move along $P$ from $\mathbf{x}_0$ towards $\Gamma$.  Thus, the area
minimizer, lies on $\Gamma$ for some parameter $|t|\le M$.

By symmetry in the underlying geometry, it is clear that the area
function has derivative zero along $\Gamma$ at $t=0$. By taking a
second derivative with automatic differentiation, we find that the
second derivative of the area function along this curve is positive
when $|t|\le M$.  Thus, the ice-ray dimer is the unique area minimizer
on this curve within this explicit neighborhood.

We remark that many of the functions that are used in the proof of the
pentagonal ice-ray conjecture are not differentiable.  Our use of
automatic differentiation is restricted to a small neighborhood of the
ice-ray dimer, where all the relevant functions are analytic.

\subsection{Meet-in-the-middle}\libel{sec:mitm}

A common algorithmic technique for reducing the time complexity of an
algorithm through greater space complexity is called {\it
  meet-in-the-middle} (MITM).  This is very closely related to the
{\it linear assembly algorithms} used in the solution to the Kepler
problem~\cite{hales2003some}. (Both methods break the problem into
subproblems of smaller complexity that are later recombined.  MITM
stores the data for recombination in a hashtable.  Linear assembly
encodes the data for recombination as linear programs.  We did not
try to solve the pentagonal ice-ray conjecture with linear programming
techniques. Such an approach might also work.)

Some introductory examples of meet-in-the-middle
algorithms can be found at the blog post \cite{mitm}.
A simple example from there is to find if there are four numbers in a
given finite set $S$ of integers that sum to zero, where repetitions
of integers are allowed.  If we calculate all possible sums of four
integers, testing if each is zero, then there are $n^4$ sums, where
$n$ is the cardinality of $S$.  The MITM solution to the problem
computes and stores (in a hashset) all sums $a+b$ of unordered pairs
of elements from $S$.  We then search the hashset for a collision,
meaning a sum $a+b$ that is the negative of another sum $c+d$: $a+b =
-(c+d)$.  Any such collision gives $a+b+c+d=0$.  The MITM solution
involves the computation of $n^2$ sums $a+b$, rather than $n^4$, for a
substantial reduction in complexity.  MITM techniques have numerous
applications to cryptography, and it is there that we first
encountered the technique.  See for example,
\cite{dinur2014dissection}
which applies MITM to a general class of dissection problems,
including the Rubik's cube.

We obtain computational bounds on the area of clusters of Delaunay
triangles (or more accurately, $P$-triangles).  Each of our clusters
will be assumed to consist of one triangle (called the {\it central
triangle}), flanked by one, two, or three additional triangles along
its edges.  We call the flanking triangles {\it peripheral triangles}.
The aim of each computation is to give a lower bound on the sum of the
areas of the triangles, subject to a collection of constraints.  Two
types of constraints are allowed: (1) constraints that can be
expressed in terms of a single triangle, and (2) assembly constraints.
An assembly constraint states that the central $P$-triangle fits
together with a flanking triangle.  In more detail, the central
triangle $T_0$ shares an edge and two pentagons $A$, $B$ with a
flanking triangle $T_1$.  Associated with $T_0$ are parameters
$d^0_{AB}$, $\theta^0_{ABC}$, $\theta^0_{BAC}$ giving the edge length
of the common edge with $T_1$, and the inclination angles of the
pentagons $A$ and $B$ with respect to that common edge.  Similarly,
associated with $T_1$ are parameters $d^1_{AB}$, $\theta^1_{ABC}$,
$\theta^1_{BAC}$.  The assembly constraint along the common edge is
\begin{equation}\libel{eqn:assembly}
d^0_{AB} = d^1_{AB},\quad 
\theta^0_{ABC} = -\theta^1_{ABC},\quad \theta^0_{BAC} = -\theta^1_{BAC}.
\end{equation}
The negative sign comes from the opposite orientations of the common
edge with respect to $T_0$ and $T_1$.

A cluster consisting of one central $P$-triangle and $k$ flanking
triangles is a point in a configuration space of dimension $6 + 3k$.
If we cover the configuration space with cubes of edge-length
$\epsilon$, then there are order $(1/\epsilon)^{6+3k}$ cubes.  This is
generally beyond our computational reach when $k>0$.

We can use MITM techniques to reduce to order $(1/\epsilon)^6$ cubes,
and this puts all our computations (barely) within the reach of a
laptop computer in 2016.  Specifically, we fix an edge size $\epsilon$ and
cover the configuration space of peripheral triangles by cubes of size
$\epsilon$, calculating area, edge lengths, inclination angles, and
other relevant quantities (using interval arithmetic) over each cube.

The idea of MITM is to place the peripheral triangle data into a hash
table, keyed by the variables $d^1_{AB}$, $\theta^1_{ABC}$,
$\theta^1_{BAC}$ that are shared with the central triangle through
Equation~\rif{eqn:assembly}.  We view Equation~\rif{eqn:assembly} as
the analogue of the collision condition $(a+b) = -(c+d)$ in the simple
example of MITM given above.  Of course, the variables $d^1_{AB}$
etc. are represented as intervals with floating point endpoints and
cannot be used directly as keys to a hashtable.  Instead we discretize
the keys in such a way that a collision of real numbers implies a
collision of the keys.

Once the hash is created, we divide the configuration space of central
triangles into cubes and compute the relevant quantities (edge
lengths, triangle area, and inclination angles) over each cube using
interval arithmetic.  The areas of the peripheral triangles are
recovered from the hash table and combined with the area of the
central triangle to get a lower bound on the sum of the areas of the
triangles in the cluster.

The entire process is iterated for smaller and smaller $\epsilon$
until the desired bound on the areas of the triangles in the cluster
is obtained.  Each time $\epsilon$ is made smaller, only the
peripheral cubes that were involved in a key collision with a central
cube are carried into the next iteration for subdivision into smaller
cubes.  Only the cubes with suitably small triangle area bounds are
carried into the next iteration.  In practice, to achieve our bounds
in Section~\rif{sec:calc}, the smallest $\epsilon$ that was required
was approximately $0.00024$.

\subsection{Preparation of the inequalities}\libel{sec:calc-prep}

Section~\rif{sec:calc} gives a sequence of inequalities that have been
proved by MITM algorithms.  We organize these calculations to permit a
more or less uniform proof of all of them.  The triangle $T_0$ is  the
central triangle.  The other triangles $T_1,\ldots$ in the cluster are the
peripheral triangles.

We prove each inequality by contradiction. Specifically, we negate the
conclusion and add it to the set of assumptions.  Then we show in each
case that the domain defined by the set of assumptions is empty. In
the computer code, we implement {\it out-of-domain} functions that
return true when the interval input lies entirely outside the given
domain.

We deform the cluster of triangles to make computations easier.  Note
that in every case except for the nonobtuse triangle $T_-$ in
Lemma~\rif{calc:pseudo-area3}, the peripheral triangles are all
subcritical.  We prove the lemmas of Section~\rif{sec:calc} in the
sequential order given in that section.  In particular, we may
assume the previous lemmas to simplify what is to be proved in those
that follow.

\begin{lemma}\libel{lemma:peripheral-o2c}
  We can assume without loss of generality that the nonobtuse triangle $T_-$ in
  Lemma~\rif{calc:pseudo-area3} is $O2C$.
\end{lemma}

\begin{proof} 
  Assume for a contradiction, that we cannot deform into a $O2C$.
  Let
  the three pentagons of $T_-$ be $A,B,C$, where $A,C$ are shared with
  the central triangle $T_0$.
  The edge shared between $T_-$ and $T_0$ has length less than $\kappa\sqrt8$.
  This implies that the angle of $T_-$ opposite the shared edge is acute.
  By Lemmas~\rif{lemma:primary}
  and~\rif{lemma:primary2}, we can assume that $B$ has primary
  contact.  By translating $B$, we can continue to deform $T_-$,
  decreasing its area, until it is a right triangle (because $O2C$ is
  assumed not to occur).  By Lemma~\rif{calc:pseudo2}, the shared edge
  with $T_0$ has length at least $1.8$, so a right angle gives the
  proof of Lemma~\rif{calc:pseudo-area3} in this case:
\[
\area\{T_0,T_1,T_-\} > 2\ao + \area(T_-) 
\ge 2\ao + 1.8\kappa > 3\acrit + \epsM.
\]
\end{proof}

We note that an earlier lemma (Lemma~\rif{lemma:2C}) shows that
subcritical peripheral triangles can be assumed to be $O2C$. Thus, we
will assume in the computations that all peripheral triangles are
$O2C$.

In the rest of this section, we describe how each calculation has been
prepared, to reduce the dimension of the configuration space, prior to
computation.

\subsubsection{Calculation~of Lemma~\rif{calc:pseudo1}} \libel{sec:pseudo1}

We repeat the statement of Lemma~\rif{calc:pseudo1}.  Let
$(T_1,T_0)\in \PD$.  Then the edge shared between $T_0$ and $T_1$ has
length less than $1.8$.  %

Negating the conclusion, we may assume that the shared edge between
triangles has length at least $1.8$.  Let $A,B,C$ be the pentagons of
$T_0$, with $A,C$ shared with $T_1$.  We drop the condition
$T_0\nRa T_1$: the selected longest edge of $T_0$ is no longer
required to be nonshared with $T_1$.  We retain only the weaker
consequence $\max(\dx{AB},\dx{BC})\ge1.8$, which follows from
the original pseudo-dimer condition and the assumed shared-edge
length at least $1.8$.  We may assume that $B$ has primary contact. By symmetry, we may assume contact between $A$ and $B$.  By Lemma~\rif{lemma:primary}, as
long as $B$ is not $O2C$, we may continue to move $B$ (decreasing area
as always) until $\max(\dx{AB},\dx{BC})=1.8$.  The calculation reduces
to three subcases:

\begin{enumerate}
\item $T_0$ has $O2C$ contact.
\item $B$ has midpointer contact along $AB$, with $\dx{AB}=1.8$.  (We
  eliminate the case $\dx{BC}=1.8$ because this would give using
  Lemma~\rif{lemma:mid-172},
\[
\area\{T_0,T_1\} > \ao + \area(T_0) \ge \ao + \area(1.8,1.8,1.72) > 2\acrit.
\]
This is contrary to the pseudo-dimer area bound.)
\item $B$ has slider contact with $A$ and $\max(\dx{AB},\dx{BC})=1.8$.
\end{enumerate}

\subsubsection{Calculation~of Lemma~\rif{calc:pseudo2}} \libel{sec:pseudo2}

We repeat the statement of Lemma~\rif{calc:pseudo2}.
Let $(T_1,T_0)\in \PD$.  The longest edge of $T_0$
has length greater than $1.8$. %

Negating the conclusion, we may assume that all edges of $T_0$ have
length at most $1.8$.  As in the previous calculation, we let $A,C$ be
the shared pentagons, and let $B$ be the third pentagon of $T_0$.
By the arguments of Section~\rif{sec:O2C}, we may assume that $T_1$
is $O2C$ and that no edge of $T_1$ is longer than its shared edge.

Now fix $A,C$ and all of $T_1$.
Minimize \(\area(T_0)\) over the position and
orientation of \(B\), subject to the local nonoverlap conditions and
\begin{equation}\label{eqn:pseudo2-constraint}
d_{BA}\le1.8,\qquad d_{BC}\le1.8,\qquad \max(d_{BA},d_{BC})\ge d_{AC}.
\end{equation}
Because of the upper bound $1.8 < \kappa\sqrt8$ on edges, the minimizing
$T_0$ is acute.
If $B$ has no primary contact, rotate it slightly about its center
to break its possible single contact. Hold fixed whichever of
$d_{BA},d_{BC}$ is larger, and move $\c_B$ a small distance along the
corresponding circle towards the shared line. This decreases the
area and the other incident distance, while preserving the upper
bounds and the lower bound on their maximum. The upper side bound
$1.8<\kappa\sqrt8$ excludes a right-angle obstruction. This
contradicts minimality, so $B$ has primary contact.  Relabeling $A$ and $C$ if necessary, we may assume
without loss of generality that $B$ and $A$ are in contact.
If the contact is not $O2C$, then a minimum
is attained when some length constraint on $d_{BA}$ or $d_{BC}$
is extremal. (Otherwise, slide $B$ along the contact or rotate $B$ around the
pointer vertex in a direction to further decrease area. Or in the degenerate
case when the slide is horizontal, slide $B$ keeping area fixed until a
length is extremal.)
Therefore the minimization leads to the alternatives:
\begin{enumerate}
\item  $B$ has $O2C$ contact, or
\item $B$ has slider or midpoint contact with $A$, and
  equality holds for some inequality in (\rif{eqn:pseudo2-constraint}).
\end{enumerate}

\subsubsection{Calculation~of Lemma~\rif{calc:pseudo-area}} \libel{sec:pseudo-area}

We repeat the statement of Lemma~\rif{calc:pseudo-area}.
Let $(T_1,T_0)\in\PD$.  Then $\area\{T_0,T_1\} > 2\acrit - \epsM$.

Negating the conclusion, we may assume that
$\area\{T_0,T_1\}\le2\acrit-\epsM$.  Let $A,B,C$ be the pentagons of $T_0$, with $A,C$
shared with $T_1$.  We deform $B$, decreasing area, until primary contact.  If it is not
$O2C$, then we continue to translate $B$.  We never encounter the long
isosceles constraint while translating $B$ because if the longest edge
and shared length have equal lengths, we have $\area\{T_0,T_1\} >
2\acrit$ by Lemmas~\rif{calc:pseudo1} and~\rif{calc:pseudo2}.  Thus,
we always reduce to $O2C$ on $T_0$.  We may continue with a squeeze
transformation (Section~\rif{sec:squeeze}) until $T_1$ is long
isosceles or $3C$.  We consider two long isosceles subcases, depending
on which of the two nonshared edges of $T_1$ has the same length as the
shared edge.

The squeeze transformation may result in $T_0$ and $T_1$ becoming
triple contact.  The calculations for dimer pairs in triple contact
were carried out with assumptions that were sufficiently general to
include this case of pseudo-dimer triple contact.

\subsubsection{Calculation~of Lemma~\rif{calc:pseudo-area3}}\libel{sec:pseudo-area3}

We repeat the statement of Lemma~\rif{calc:pseudo-area3}.  Let
$(T_1,T_0)\in\PD$. Assume $T_0\Ra T_-$ with $T_-$ nonobtuse.  Then $\area\{T_0,T_1,T_-\} >
3\acrit + \epsM$. %

By Lemma~\rif{lemma:2C} for $T_1$ and
Lemma~\rif{lemma:peripheral-o2c} for $T_-$, we may assume that both
peripheral triangles are $O2C$.  Negating the conclusion, assume
$\area\{T_0,T_1,T_-\}\le3\acrit+\epsM$.

Subtracting the pair-area bound of
Lemma~\rif{calc:pseudo-area} gives
\begin{align*}
\area(T_-) &= \area\{T_0,T_1,T_-\} - \area\{T_0,T_1\} \\
&< (3\acrit+\epsM) - (2\acrit - \epsM) = \acrit + 2\epsM < \acrit+\epsN.
\end{align*}

We let $A,B,C$ be the pentagons in the central triangle $T_0$, where
$A$ is shared with $T_0,T_1,T_-$ and $B$ is shared between $T_0$ and
$T_-$.  We make a deformation which moves $B$ by translation along the
fixed line through $(\c_A,\c_B)$ and keeps all other pentagons of
$T_1$ and $T_0$ fixed.  A stopping condition occurs when $B$ comes
into contact with $A$ or $C$.  From the perspective of the triangle
$T_-$, it undergoes a squeezing transformation, as described in
Section~\rif{sec:squeeze}.  The pentagon at the nonshared vertex of
$T_-$ is translated to maintain $02C$ contact.  By
Lemma~\rif{lemma:squeeze-area}, squeezing $T_-$ decreases its area.
Furthermore, the translation of $B$ toward $\c_A$ decreases the area
of $T_0$.  The area of $T_1$ remains fixed.  Thus, the deformation
decreases $\area\{T_0,T_1,T_-\}$.

Renaming pentagons of the central triangle, we assume that $\bar A$
and $\bar C$ are in contact, with $\bar A$ pointing to $\bar C$.  We
consider six cases: each permutation on three letters determines a
choice for the edge of $T_0$ shared with $T_1$, and a choice of the
edge of $T_0$ shared with $T_-$.

\subsubsection{Calculation~of Lemma~\rif{calc:dimer-isosc}}\libel{sec:dimer-isosc}

We repeat the statement of Lemma~\rif{calc:dimer-isosc}.  Let
$(T_1,T_0)\in DP$.  Then $T_1$ is not both $O2C$ and long isosceles.

Negating the conclusion, we may assume that $T_1$ is $O2C$ and long
isosceles and that $T_0$ is $O2C$.  We perform the calculation without
further preparation.

\subsubsection{Calculation~of Lemma~\rif{calc:large}}\libel{sec:large} %

We repeat the statement of Lemma~\rif{calc:large}.  Let $\{T_0,T_1\}$
be given nonobtuse $P$-triangles (not necessarily a pseudo-dimer) such that
$\area(T_1)\le \acrit$ and $T_1\Ra T_0$.  Assume that there is a
nonshared edge $e$ of $T_0$ of length greater than $1.8$ and that the
angle is not large along $(T_0,e)$.  Then $\area\{T_0,T_1\} > 2\acrit
+ \epsM$. %

Negating the conclusion, we may assume that $\area\{T_0,T_1\}\le
2\acrit + \epsM$.  We assume that $T_0$ has a nonshared edge of length
at least $1.8$ and that the angle is not large along that edge.
Let $A,B,C$ be the pentagons of $T_0$, with $B$ the outer pentagon,
and such that $e$ is the edge $(\c_B,\c_C)$.
By Lemma~\rif{lemma:172}, we may assume that the length of the
shared edge $(\c_A,\c_C)$ has length at least $1.72$.
By earlier reductions (Section~\rif{sec:O2C}), we may assume
that $T_1$ is $O2C$.  The arguments there were made for dimer pairs, but the argument
is sufficiently general to apply to $T_1$.

With the orientations fixed, choose a minimum-area admissible
position of $B$, retaining $d_{BC}\ge1.8$, nonoverlap,
nonobtuseness, and the original area upper bound. These closed
constraints give a compact set. The translations considered below
are not required to decrease edge lengths. We will show that a
minimum is jammed, meaning that at least two of the following
unilateral constraints are met.
\begin{enumerate}
\item The edge $e$ has length $1.8$.
\item Pentagons $B$ and $A$ have contact.
\item Pentagons $B$ and $C$ have contact.
\end{enumerate}
 
A right-angle boundary is impossible. The shared edge has length
greater than $1.72$ and less than $2.1<\kappa\sqrt8$. It cannot be
the hypotenuse of a right triangle of pentagon centers, and is
therefore a leg. Thus, if $T_0$ were right, we would have
\[
\area\{T_1,T_0\} > \ao + \frac12(2\kappa)(1.72) > 2 \acrit + \epsM.
\]
If $B$ satisfies at most one unilateral constraint, it can be translated
horizontally (with horizontal axis $(\c_A,\c_C)$) in the direction that breaks a
possible constraint,
then the translation can be continued vertically to further decrease area until
at least one unilateral constraint is again met.
A horizontal supporting contact would give height at least
$2\kappa$ and the already excluded lower bound $\ao+1.72\kappa$
for the pair area. At any other position satisfying at most one
of the remaining unilateral constraints, the local translation
just described decreases area. Hence a minimum satisfies at least
two of them and is one of the stated terminal configurations.

\subsubsection{Calculation~of Lemma~\rif{calc:pent3}}\libel{sec:pent3} %

We repeat the statement of Lemma~\rif{calc:pent3}. Let
$T_0,T_1^0,T_1^1$ be distinct nonobtuse $P$-triangles such that
$T_1^i\Ra T_0$ and $\area(T_1^i)\le\acrit$ for $i=0,1$. Then $\area\{T_0,T_1^0,T_1^1\} > 3\acrit +
\epsM$. %

We assume that $\area\{T_0,T_1^0,T_1^1\}\le 3\acrit + \epsM$.
By earlier reductions (Section~\rif{sec:O2C}), we may assume
that $T_1^0,T_1^1$ are $O2C$.  The arguments there were made for dimer pairs, but the argument
is sufficiently general to apply to $T_1^i$.
The reduction may be assumed to preserve the shared edge between $T_1^i$ and $T_0$
as a longest edge of $T_1^i$, for $i=0,1$.

Let $A_1^i$, $B_1^i$, $C$ be the pentagons of $T_1^i$,
and let $A_1^0,A_1^1,C$ be the pentagons of $T_0$.

Considering $i=0,1$ in turn, we deform the cluster by fixing all pentagons
except $A_1^i$ and $B_1^i$, and moving them by continuous translation.
The motion moves $\c_{A_1^i}$ linearly directly toward $\c_C$ and $\c_{B_1^i}$
moves according to the squeeze deformation, Lemma~\rif{lemma:squeeze-area},
in an area decreasing way.  This linear motion also decreases the area
of $T_0$ and fixes the area of the remaining triangle.
The motion stops as soon as at least one of the
following two stopping conditions is met:
\begin{enumerate}
\item $A_1^i$ comes in contact with another pentagon of $T_0$.
\item The triangle $T_1^i$ becomes long isosceles.
\end{enumerate}
Applying the motion to $i=0,1$, at least one of the conditions holds for each $i$.

Since $T_1^i$ are subcritical, by Lemma~\rif{lemma:right}, the edges
$(\c_{A_1^i},\c_C)$ have length initially less than $\kappa\sqrt8$ and
are not increased by the deformations.  It follows that the angles of $T_0$
at $\c_{A_1^i}$ cannot become right.  Further, the angle of $T_0$ at $\c_C$
remains constant.  Hence $T_0$ remains a nonobtuse triangle throughout
the motions.

\subsubsection{Calculation~of Lemma~\rif{calc:pent4}} \libel{sec:pent4}

We repeat the statement of Lemma~\rif{calc:pent4}. Let
$T_0,T_1^0,T_1^1,T_1^2$ be distinct nonobtuse $P$-triangles such that
$T_1^i\Ra T_0$ and $\area(T_1^i)\le\acrit$ for $i=0,1,2$. Then $\area\{T_0,T_1^0,T_1^1,T_1^2\} >
4\acrit$. %

Each $T_1^i$ is $O2C$.  We carry out the calculation as a single case,
using MITM as usual as described above.  There is an $S_3$-symmetry %
that we can exploit to reduce the number of cases. 


      
%\input{version1}
      
      \bibliography{pentagon} 
%\bibliographystyle{plain}
      
      \bibliographystyle{plainnat}
%\bibliography{announce_refs}



%% shell:>makeindex index/Index
%% shell:>makeindex index/Notation
%% edit index/Notation.ind to put an \indexspace before Greek entries.

%\printindex{index/Index}{General index}
%\printindex{index/Notation}{Notation index}

%\input{frontmatter}

\end{document}
